\SGAchapter{VI\textsubscript{B}}{群概形的一般性质}{Généralités sur les schémas en groupes}
\SourcePage{335}
\begin{center}J.-E. Bertin\end{center}
本讲并不对应研讨班中的任何一次口头报告，目的是汇集一些关于\Term{group-scheme}{群概形}的常用技术结果。\footnote{本讲自油印本以来作了相当大的修订，尤其第 10、11 节已全部重写。\textbf{编者注（1）：}第 5 节也已重写，参见编者注（34）和（46）。}

\SGAsection{1}{域上局部有限型群的态射}
\textbf{1.1。}设 $A$ 是一个 Artin \Term{local-ring}{局部环}，$G,H$ 是两个 $A$-群，\footnote{\textbf{编者注（2）：}回顾第 VI$_A$ 讲开头所采用的约定：$A$-群指 $A$-\Term{group-scheme}{群概形}（这样的\Term{scheme}{概形}是\Term{separated}{分离}的，参见第 VI$_A$ 讲 0.3）。}$u:G\to H$ 是一个 $A$-\Term{group-morphism}{群态射}。于是 $u$ 诱导\Term{group-morphism}{群态射} $u(A):G(A)\to H(A)$。\footnote{\textbf{编者注（3）：}将“\Term{homomorphism}{同态}”改成了“\Term{morphism}{态射}”。}由于 $H(A)$ 通过\Term{right-translation}{右平移}\Term{action}{作用}于 $H$，$u(A)$ 通过\Term{restriction}{限制}给出 $G(A)$ 在 $H$ 上的\Term{action}{作用}。这个\Term{action}{作用}与\Term{morphism}{态射} $u$ 以及 $G(A)$ 通过\Term{right-translation}{右平移}在 $G$ 上的\Term{action}{作用}相容。$G(A)$ 在 $G$ 的\Term{strictly-rational-point}{严格有理点}（定义见第 VI$_A$ 讲 0.4）上\Term{transitive-action}{传递地作用}，所以这些点“对于 $u$ 都有相同的行为”；由此产生下列性质。
\begin{SGAresult}{命题}{1.2}
设 $u:G\to H$ 是\Term{locally-finite-type}{局部有限型} $A$-群之间的\Term{quasi-compact}{拟紧}\Term{group-morphism}{群态射}。则\Term{set}{集合} $u(G)$ 在 $H$ 中是闭的，并且
\[
\dim G=\dim u(G)+\dim\operatorname{Ker}u.
\]
\end{SGAresult}
\footnote{\textbf{编者注（4）：}这个陈述也见第 VI$_A$ 讲 2.5.2。}
由于 $u$ 与 $G,H$ 的\Term{inversion-morphism}{取逆态射}交换，像 $u(G)$ 在 $H$ 的\Term{inversion-morphism}{取逆态射}下\Term{invariant}{不变}；因而它在 $H$ 中的\Term{closure}{闭包} $\overline{u(G)}$ 也\Term{invariant}{不变}。另一方面，设 $L$ 是 $H\times_A H$ 中两个\Term{projection}{投影}都属于 $u(G)$ 的点的\Term{set}{集合}。显然 $L$ 是\Term{morphism}{态射}
\SourcePage{336}
$u\times_Au:G\times_AG\to H\times_AH$ 的像；因而 $H$ 的\Term{multiplication-morphism}{乘法态射}将 $L$ 映入 $u(G)$，即 $u(G)\cdot u(G)=u(G)$。下面的\Term{lemma}{引理} 1.2.1 则说明，$L$ 在 $H\times_AH$ 中的\Term{closure}{闭包} $\overline L$，正是两个\Term{projection}{投影}都属于 $\overline{u(G)}$ 的点所组成的\Term{set}{集合}。因此
\[
\overline{u(G)}\cdot\overline{u(G)}=\overline{u(G)}.
\]
所以，以 $\overline{u(G)}$ 为底\Term{space}{空间}的 $H$ 的\Term{reduced}{约化}\Term{subscheme}{子概形}，在\Term{category}{范畴} $(\mathrm{Sch}/k)_{\mathrm{r\text{é}d}}$ 中自然具有\Term{group-structure}{群结构}，其中 $k$ 是 $A$ 的\Term{residue-field}{剩余域}（参见第 VI$_A$ 讲 0.2）。

先证 1.2 的第一个断言。将 $A$ 换成其\Term{residue-field}{剩余域} $k$ 的\Term{algebraic-closure}{代数闭包}，便可假设 $A$ 是\Term{algebraically-closed-field}{代数闭域} $k$（参见 EGA IV$_2$, 2.3.12）。将 $u$ 换成 $u_{\mathrm{r\text{é}d}}:G_{\mathrm{r\text{é}d}}\to H_{\mathrm{r\text{é}d}}$，还可假设 $G,H$ 都\Term{reduced}{约化}。此时，如刚才所见，$\overline{u(G)}$ 是 $G$ 的一个\Term{reduced}{约化}\Term{group-subscheme}{子群概形}的底\Term{space}{空间}；所以可假设 $u$ 是\Term{dominant}{支配}的。于是 $G(k)$ 在 $H$ 的\Term{connected-component}{连通分支}\Term{set}{集合}上传递地\Term{action}{作用}，只需证明 $u(G)\cap H^0$ 是闭的；这就归结为 $H$ \Term{connected}{连通}的情形。此时 $H$ 是\Term{irreducible}{不可约}且\Term{finite-type}{有限型}的（第 VI$_A$ 讲 2.4）。$u$ \Term{quasi-compact}{拟紧}且\Term{locally-finite-type}{局部有限型}，因而是\Term{finite-type}{有限型}的；$H$ 又是 \Term{noetherian}{诺特}\Term{scheme}{概形}，故 $u(G)$ 是\Term{constructible}{可构造}的（EGA IV$_1$, 1.8.5），从而包含 $H$ 的一个\Term{open-subset}{开集} $V$（EGA 0$_{\mathrm{III}}$, 9.2.2）。由第 VI$_A$ 讲 0.5，
\[
H=V\cdot V\subset u(G)\cdot u(G)=u(G).
\]

再证第二个断言。首先回顾，\Term{kernel}{核}\Term{functor}{函子} $\operatorname{Ker}u$（参见第 I 讲 2.3.6.1）由 $u^{-1}(e)$ \Term{representable}{表示}，其中 $e$ 是 $H$ 的\Term{identity-element}{单位元}。所以，当 $u$ \Term{locally-finite-type}{局部有限型}时，$\operatorname{Ker}u$ 在 $A$ 上也是\Term{locally-finite-type}{局部有限型}的。与上面一样，这次利用 EGA IV$_2$, 4.1.4，可归结为 $A$ 是\Term{algebraically-closed-field}{代数闭域} $k$ 的情形。还可假设 $G,H$ 都\Term{irreducible}{不可约}、\Term{finite-type}{有限型}，且 $u$ \Term{dominant}{支配}：事实上，$k$ \Term{algebraically-closed-field}{代数闭}，所以 $G(k)$ 传递地置换的 $G$ 的各\Term{connected-component}{连通分支}显然都有相同\Term{dimension}{维数}；若 $u^0$ 是 $u$ 在 $G^0$ 上的\Term{restriction}{限制}，则 $(\operatorname{Ker}u)^0\subset\operatorname{Ker}u^0$，且 $\dim\operatorname{Ker}u^0=\dim\operatorname{Ker}u$。同样可见，此时 $u$ 在 $k$ 上是\Term{finite-type}{有限型}的。若 $\eta$ 为 $H$ 的\Term{generic-point}{泛点}，则
\[
\dim u^{-1}(\eta)=\dim G-\dim H
\]
（EGA IV$_3$, 10.6.1 (ii)）。由 EGA IV$_3$, 9.2.3 和 9.2.6，满足 $\dim u^{-1}(y)=\dim u^{-1}(\eta)$ 的 $y\in H$ 组成的\Term{set}{集合}包含一个\Term{nonempty}{非空}\Term{open-subset}{开集} $V$。因 $u$ \Term{dominant}{支配}，$U=u^{-1}(V)$ 是 $G$ 的\Term{nonempty}{非空}\Term{open-subset}{开集}；$G$ 是 Jacobson \Term{scheme}{概形}（EGA IV$_3$, 10.4.7），故 $U$ 含有 $G$ 的一个\Term{closed-point}{闭点} $x$。于是\Term{right-translation}{右平移} $r_x$ 给出 $\operatorname{Ker}u$ 到 $u^{-1}(u(x))$ 的\Term{isomorphism}{同构}，所以
\[
\dim\operatorname{Ker}u=\dim u^{-1}(u(x))=\dim u^{-1}(\eta)=\dim G-\dim H.
\]
\begin{SGAresult}{\Term{lemma}{引理}}{1.2.1}
设 $f:X'\to X$ 和 $g:Y'\to Y$ 是 Artin \Term{local-ring}{局部环} $A$ 上的\Term{scheme}{概形}的两个\Term{quasi-compact}{拟紧}\Term{dominant}{支配态射}。则 $f\times_Ag:X'\times_AY'\to X\times_AY$ 是\Term{dominant}{支配}的（并且\Term{quasi-compact}{拟紧}）。
\end{SGAresult}
\begin{proof}
有 $f\times_Ag=(\operatorname{id}_X\times_Ag)\circ(f\times_A\operatorname{id}_{Y'})$，所以只需证明 $f\times_A\operatorname{id}_{Y'}$ 和 $\operatorname{id}_X\times_Ag$ 都\Term{dominant}{支配}。为此，可将 $A$ 换成其\Term{residue-field}{剩余域} $k$。此时 $X,Y'$ 在 $A=k$ 上\Term{flat}{平坦}，而 $f\times_A\operatorname{id}_{Y'}$（相应地 $\operatorname{id}_X\times_Ag$）由 $f$（相应地 $g$）经\Term{flat}{平坦}\Term{base-change}{基变换} $Y'\to A$（相应地 $X\to A$）得到，故由 EGA IV$_2$, 2.3.7，它是\Term{dominant}{支配}的（并且\Term{quasi-compact}{拟紧}）。
\end{proof}
\SourcePage{337}
\begin{SGAresult}{反例}{1.2.2}
设 $k$ 是\Term{characteristic}{特征}为 $0$ 的域，$G$ 是\Term{constant-group}{常值} $k$-群 $\mathbf Z$，$H$ 是\Term{additive-group}{加法} $k$-群 $\mathbf G_{a,k}$。若 $u:G\to H$ 是非零 $k$-\Term{group-morphism}{群态射}，则 $u(G)$ 在 $H$ 中不闭。
\end{SGAresult}
\begin{SGAresult}{命题}{1.3}
\footnote{\textbf{编者注（5）：}原文中 1.3、1.3.1 和 1.3.1.1 都是在 $A$ 为域时陈述的。为求完整，我们处理了 Artin \Term{local-ring}{局部环}的情形，并为此加入\Term{lemma}{引理} 1.3.0。}设 $A$ 是 Artin \Term{local-ring}{局部环}，$k$ 为其\Term{residue-field}{剩余域}，$G$ 是\Term{locally-finite-type}{局部有限型}且\Term{flat}{平坦}的 $A$-群，$u:G\to H$ 是 $A$-\Term{group-morphism}{群态射}。下列断言\Term{equivalence}{等价}：
\begin{enumerate}
\item[(i)] $u$ 在 $G$ 的某点\Term{flat}{平坦}（相应地\Term{quasi-finite}{拟有限}、\Term{unramified}{非分歧}、\Term{smooth}{光滑}、\Term{etale}{平展}）。
\item[(ii)] $u$ \Term{flat}{平坦}（相应地\Term{quasi-finite}{拟有限}、\Term{unramified}{非分歧}、\Term{smooth}{光滑}、\Term{etale}{平展}）。
\end{enumerate}
\end{SGAresult}
\footnote{\textbf{编者注（6）：}\Term{scheme}{概形}\Term{morphism}{态射} $f:X\to Y$ 在 $x$ 处称为\Term{finite-type}{有限型}（相应地\Term{finite-presentation}{有限呈示}、\Term{quasi-finite}{拟有限}），是指存在分别含 $f(x)$ 和 $x$ 的\Term{affine}{仿射}\Term{open-subset}{开集} $V,U$，使 $f(U)\subset V$，且 $\mathcal O_X(U)$ 是\Term{finite-type}{有限型}的 $\mathcal O_Y(V)$-\Term{algebra}{代数}（相应地是\Term{finite-presentation}{有限呈示}的；相应地还要求对每个 $x'\in U$，\Term{fibre}{纤维} $f^{-1}(f(x'))$ \Term{finite}{有限}；也参见第 V 讲编者注（40））。若 $f$ 在每个点 $x$ 都有相应性质，则称为\Term{locally-quasi-finite}{局部拟有限}（相应地\Term{locally-finite-type}{局部有限型}、\Term{locally-finite-presentation}{局部有限呈示}）。另一方面，称 $f$ 在 $x$ 处\Term{smooth}{光滑}（相应地\Term{unramified}{非分歧}、\Term{etale}{平展}），是指存在 $x$ 的\Term{open-neighborhood}{开邻域} $U$，使 $f|_U:U\to Y$ \Term{smooth}{光滑}（相应地\Term{unramified}{非分歧}、\Term{etale}{平展}）。由定义可见，$f$ 为\Term{finite-presentation}{有限呈示}、\Term{finite-type}{有限型}、\Term{quasi-finite}{拟有限}、\Term{smooth}{光滑}、\Term{unramified}{非分歧}、\Term{etale}{平展}的点集都是 $X$ 的\Term{open-subset}{开集}。此外，\Term{locally-finite-presentation}{局部有限呈示}\Term{morphism}{态射}的\Term{flat-locus}{平坦点集}是\Term{open-subset}{开集}（EGA IV$_3$, 11.3.1），故 $f$ 为\Term{finite-presentation}{有限呈示}且\Term{flat}{平坦}的点集也是 $X$ 的\Term{open-subset}{开集}。下面将反复使用这些事实。}
\begin{proof}
只需证明 (i) \Term{pushout}{推出} (ii)。先用下面的\Term{lemma}{引理}。
\end{proof}
\begin{SGAresult}{\Term{lemma}{引理}}{1.3.0}
设 $A\to B\to C$ 是\Term{commutative-ring}{交换环}的\Term{morphism}{态射}，$\mathfrak n$ 是 $A$ 的\Term{nilpotent-ideal}{幂零理想}。假设 $C/\mathfrak n C$ 是\Term{finite-type}{有限型} $(B/\mathfrak n B)$-\Term{algebra}{代数}。
\begin{enumerate}
\item[(i)] 则 $C$ 是\Term{finite-type}{有限型} $B$-\Term{algebra}{代数}。
\item[(ii)] 此外，若 $C$ 在 $A$ 上\Term{flat}{平坦}，且 $C/\mathfrak n C$ 是\Term{finite-presentation}{有限呈示} $(B/\mathfrak n B)$-\Term{algebra}{代数}，则 $C$ 是\Term{finite-presentation}{有限呈示} $B$-\Term{algebra}{代数}。
\end{enumerate}
\end{SGAresult}
\begin{proof}
取 $x_1,\ldots,x_n\in C$，使其像生成 $(B/\mathfrak n)$-\Term{algebra}{代数} $C/\mathfrak n C$。由\Term{nilpotent}{幂零}\Term{form}{形式}的 \Term{nakayama-lemma}{Nakayama 引理}，$x_i$ 生成 $B$-\Term{algebra}{代数} $C$，这证明了 (i)。令 $\varphi:B[X_1,\ldots,X_n]\to C$ 为所得\Term{surjective}{满射}，$I=\operatorname{Ker}\varphi$。

再假设 $C$ 在 $A$ 上\Term{flat}{平坦}，且 $C/\mathfrak n C$ 在 $B/\mathfrak n B$ 上\Term{finite-presentation}{有限呈示}。一方面，$I/\mathfrak n I$ 可等同于 $\overline\varphi=\varphi\otimes_A(A/\mathfrak n)$ 的\Term{kernel}{核}；另一方面，由 EGA IV$_1$, 1.4.4，$\overline\varphi$ 的\Term{kernel}{核}是\Term{finitely-generated}{有限生成}\Term{ideal}{理想}。取 $P_1,\ldots,P_s\in I$，使其像生成\Term{ideal}{理想} $I/\mathfrak n I$；\Term{nilpotent}{幂零}\Term{form}{形式}的 \Term{nakayama-lemma}{Nakayama 引理}说明它们生成 $I$。这证明了 (ii)。
\end{proof}
\begin{proof}[命题 1.3 的证明（续）]
任取 $x\in G$。因 $G$ 在 $A$ 上\Term{flat}{平坦}，由\Term{fibrewise-flatness-criterion}{逐纤维平坦性判据}的 EGA IV$_3$, 11.3.10.2 \Term{form}{形式}，只要 $u\otimes_Ak$ 在 $x$ 处\Term{flat}{平坦}，$u$ 就在 $x$ 处\Term{flat}{平坦}。同样，由上面的\Term{lemma}{引理}，
\SourcePage{338}
只要 $u\otimes_Ak$ 在 $x$ 处\Term{finite-type}{有限型}（相应地\Term{finite-presentation}{有限呈示}），$u$ 就有相同性质。其他性质此时都可在\Term{fibre}{纤维}上检验（\Term{unramified}{非分歧}、\Term{smooth}{光滑}、\Term{etale}{平展}分别参见 EGA IV$_4$, 17.4.1、17.5.1、17.6.1），故归结为 $A=k$ 的情形。

现设 $x\in G$ 满足 1.3 (i) 中的某一条件。所考虑的性质都可作 (fpqc) \Term{descent}{下降}（参见 EGA IV, 2.5.1、2.7.1、17.7.1），所以将 $k$ 换成 $\kappa(x)$ 的\Term{algebraic-closure}{代数闭包}，可归结为 $k$ \Term{algebraically-closed-field}{代数闭}且 $x\in G(k)$ 的情形。

$G$ 是 Jacobson \Term{scheme}{概形}（参见 EGA IV$_3$, 10.4.7），而 $u$ 为\Term{flat}{平坦}（相应地\Term{quasi-finite}{拟有限}、\Term{unramified}{非分歧}、\Term{smooth}{光滑}、\Term{etale}{平展}）的点集 $W$ 对\Term{generization}{一般化}\Term{stable}{稳定}\footnote{\textbf{编者注（7）：}原文称 $W_{\mathrm{plat}}$ 是\Term{open-subset}{开集}，这一点似乎并非先验显然。}（相应地是\Term{open-subset}{开集}），故只需证明每个 $y\in G(k)$ 都属于 $W$。而对这样的 $y$，\Term{translation}{平移} $r_y\circ r_x^{-1}$ 将 $x$ 送到 $y$，所以 $u$ 在 $y$ 处也有所需性质，即 $y\in W$。
\end{proof}
\begin{SGAresult}{推论}{1.3.1}
设 $A$ 是 Artin \Term{local-ring}{局部环}，$k$ 是其\Term{residue-field}{剩余域}，$G$ 是\Term{flat}{平坦} $A$-群。下列断言\Term{equivalence}{等价}：\footnote{\textbf{编者注（8）：}只假设 $G$ 在 $A$ 上某一点\Term{flat}{平坦}是不够的。例如，设 $A\ne k$，$H$ 为\Term{constant-group}{常值} $A$-群 $(\mathbf Z/2\mathbf Z)_A$，$G$ 是 $H$ 中非\Term{identity-component}{单位分支}\Term{reduced}{约化}的那个闭子 $A$-群；则\Term{structure-morphism}{结构态射} $G\to\operatorname{Spec}A$ 在\Term{identity-element}{单位点} $e$ 处是\Term{local-isomorphism}{局部同构}，但在 $g\ne e$ 处不\Term{flat}{平坦}。}
\begin{enumerate}
\item[(i)] $G$ 在 $A$ 上某一点\Term{locally-quasi-finite}{局部拟有限}（相应地\Term{unramified}{非分歧}、\Term{smooth}{光滑}、\Term{etale}{平展}）。
\item[(ii)] $G$ 在 $A$ 上\Term{locally-quasi-finite}{局部拟有限}（相应地\Term{unramified}{非分歧}、\Term{smooth}{光滑}、\Term{etale}{平展}）。
\end{enumerate}
\end{SGAresult}
\begin{proof}
\footnote{\textbf{编者注（9）：}加入了下面这句话，并简化了\Term{lemma}{引理} 1.3.1.1 的证明，用 EGA IV, 2.7.1 代替同处 17.7.5。}若 $G$ 在某点 $x$ 满足 (i) 中任一条件，则 $x$ 有一个在 $A$ 上\Term{finite-type}{有限型}的\Term{open-neighborhood}{开邻域} $U$。因而由下面的\Term{lemma}{引理}，只需将 1.3 应用于 $H$ 为单位 $A$-群的情形。
\end{proof}
\begin{SGAresult}{\Term{lemma}{引理}}{1.3.1.1}
设 $A$ 是 Artin \Term{local-ring}{局部环}，$G$ 是 $A$-群。若 $G$ 有一个在 $A$ 上\Term{finite-type}{有限型}的\Term{nonempty}{非空}\Term{open-subset}{开集}，则 $G$ 在 $A$ 上\Term{locally-finite-type}{局部有限型}。
\end{SGAresult}
\begin{proof}
由\Term{lemma}{引理} 1.3.0，可假设 $A$ 等于其\Term{residue-field}{剩余域} $k$。再由 (fpqc) \Term{descent}{下降}，可假设 $k$ \Term{algebraically-closed-field}{代数闭}（参见 EGA IV$_2$, 2.7.1）。令 $V$ 为 $G$ 中在 $k$ 上\Term{finite-type}{有限型}的点组成的\Term{open-subset}{开集}；由假设 $V\ne\varnothing$。因 $G$ 是 Jacobson \Term{scheme}{概形}，$V$ 含有一个\Term{closed-point}{闭点} $x$；要证 $V=G$，只需证明 $G$ 的每个\Term{closed-point}{闭点} $y$ 都在 $V$ 中。对这样的 $y$，\Term{translation}{平移} $r_y\circ r_x^{-1}$ 将 $x$ 送到 $y$，故 $y\in V$。
\end{proof}
\begin{SGAresult}{推论}{1.3.2}
设 $A$ 是 Artin \Term{local-ring}{局部环}，$u:G\to H$ 是\Term{locally-finite-type}{局部有限型} $A$-群之间的\Term{group-morphism}{群态射}。下列断言\Term{equivalence}{等价}：
\begin{enumerate}
\item[(i)] $u$ \Term{universally-open}{普遍开}；
\item[(ii)] $u$ 是\Term{open-morphism}{开态射}；
\item[(iii)] $u$ 在 $G$ 的某点是\Term{open-morphism}{开}的；
\item[(iv)] 由 $u$ 诱导的\Term{morphism}{态射} $u^0:G^0\to H^0$ 是\Term{dominant}{支配}的；
\end{enumerate}
\SourcePage{339}
\begin{enumerate}
\item[(iv$'$)] $u^0$ 是\Term{surjective}{满射}；
\item[(v)] 存在 $G$ 的一个\Term{connected-component}{连通分支} $C$，若以 $D$ 表示 $H$ 中包含 $u(C)$ 的\Term{connected-component}{连通分支}，则由 $u$ 诱导的\Term{morphism}{态射} $u':C\to D$ 是\Term{dominant}{支配}的。
\end{enumerate}
\end{SGAresult}
\begin{proof}
(i)$\Rightarrow$(ii)$\Rightarrow$(iii) 和 (iv$'$)$\Rightarrow$(iv)$\Rightarrow$(v) 都显然。$G^0$ 在 $A$ 上\Term{finite-type}{有限型}（第 VI$_A$ 讲 2.4），所以是 \Term{noetherian}{诺特}\Term{scheme}{概形}，故 $u^0$ \Term{quasi-compact}{拟紧}；由 1.2，$u^0(G^0)$ 在 $H^0$ 中闭，故 (iv)$\Rightarrow$(iv$'$)。另一方面，$G^0$（相应地 $C$）在 $G$ 中开（第 VI$_A$ 讲 2.3），而 $H^0$（相应地 $D$）\Term{irreducible}{不可约}（第 VI$_A$ 讲 2.4.1），可见 (ii) \Term{pushout}{推出} (iv)（相应地 (iii) \Term{pushout}{推出} (v)）。因此只需证明 (v) \Term{pushout}{推出} (i)。

赋予 $G$ 的\Term{open-subset}{开集} $C$（相应地 $H$ 的\Term{open-subset}{开集} $D$）以诱导的\Term{scheme}{概形}结构，以 $u'$ 表示由 $u$ 诱导的\Term{morphism}{态射} $C\to D$。令 $k$ 为 $A$ 的\Term{residue-field}{剩余域}。\footnote{\textbf{编者注（10）：}下面将原文展开说明。}由于\Term{base-change}{基变换} $\operatorname{Spec}k\to\operatorname{Spec}A$ 是\Term{universal-homeomorphism}{普遍同胚}，可假设 $A=k$。由假设，$u'$ \Term{dominant}{支配}；又因 $C$ 在 $k$ 上\Term{finite-type}{有限型}（第 VI$_A$ 讲 2.4.1），故是 \Term{noetherian}{诺特}\Term{scheme}{概形}，所以 $u'$ \Term{quasi-compact}{拟紧}。由 EGA IV$_2$, 2.3.7，$u'\otimes_k\overline k$ 仍\Term{quasi-compact}{拟紧}且\Term{dominant}{支配}，其中 $\overline k$ 是 $k$ 的\Term{algebraic-closure}{代数闭包}。$C\otimes_k\overline k$ 是 $G\otimes_k\overline k$ 的若干\Term{connected-component}{连通分支}之并，所以 $u\otimes_k\overline k:G\otimes_k\overline k\to H\otimes_k\overline k$ 满足 (v)。由 EGA IV$_2$, 2.6.4，遂归结为 $A=k$ 是\Term{algebraically-closed-field}{代数闭域}的情形。

此时还可将 $u$ 换成 $u_{\mathrm{r\text{é}d}}$，从而假设 $H$ \Term{reduced}{约化}。令 $\xi,\eta$ 分别为 $C,D$ 的\Term{generic-point}{泛点}。$u'$ \Term{quasi-compact}{拟紧}且\Term{dominant}{支配}，故 $u'(\xi)=\eta$（参见 EGA IV$_1$, 1.1.5）。另一方面，由于 $H$ \Term{reduced}{约化}，\Term{local-ring}{局部环} $\mathcal O_{H,\eta}$ 是域，故 $u'$ 在 $\xi$ 处\Term{flat}{平坦}。\footnote{\textbf{编者注（11）：}原文引用了\Term{generic-flatness}{泛平坦性}\Term{theorem}{定理}（EGA IV$_2$, 6.9.1），这里并不需要。}由 1.3，$u$ \Term{flat}{平坦}；又因 $u$ \Term{locally-finite-type}{局部有限型}，$H$ \Term{locally-noetherian}{局部诺特}，故 $u$ \Term{locally-finite-presentation}{局部有限呈示}。由 EGA IV$_2$, 2.4.6，$u$ \Term{universally-open}{普遍开}。
\end{proof}
\begin{SGAresult}{命题}{1.4}
设 $A$ 是 Artin \Term{local-ring}{局部环}，$u:G\to H$ 是\Term{locally-finite-type}{局部有限型} $A$-群之间的\Term{quasi-compact}{拟紧}\Term{group-morphism}{群态射}。下列断言\Term{equivalence}{等价}：
\begin{enumerate}
\item[(i)] $u$ \Term{proper}{固有}；
\item[(ii)] 存在 $h\in H$，使\Term{fibre}{纤维} $u^{-1}(h)$ \Term{nonempty}{非空}且在 $\kappa(h)$ 上\Term{proper}{固有}；
\item[(iii)] $\operatorname{Ker}u$ 在 $A$ 上\Term{proper}{固有}。
\end{enumerate}
\end{SGAresult}
\begin{proof}
显然 (i) \Term{pushout}{推出} (iii)，而 (iii) \Term{pushout}{推出} (ii)。另一方面，由假设 $u$ \Term{finite-type}{有限型}，且因 $G$ \Term{separated}{分离}（第 VI$_A$ 讲 0.3），$u$ 也\Term{separated}{分离}（EGA I, 5.5.1）。所以只需证明 (ii) \Term{pushout}{推出} $u$ \Term{universally-closed}{普遍闭}；于是可假设 $A$ 等于其\Term{residue-field}{剩余域} $k$。\footnote{\textbf{编者注（12）：}加入了下面这句话。}取 $\kappa(h)$ 的\Term{algebraic-closure}{代数闭包} $k'$，令 $u':G'\to H'$ 是 $u$ 经\Term{base-change}{基变换}所得的\Term{morphism}{态射}，$h'$ 是 $H'$ 中位于 $h$ 上方的一点。则\Term{fibre}{纤维} $u'^{-1}(h')=u^{-1}(h)\times_{\kappa(h)}k'$
\SourcePage{340}
\Term{nonempty}{非空}且\Term{proper}{固有}，而只需证明 $u'$ \Term{proper}{固有}（EGA IV$_2$, 2.6.4）。因此可假设 $k$ \Term{algebraically-closed-field}{代数闭}且 $h\in H(k)$。

由 1.2，$u(G)$ 是 $H$ 的一个\Term{reduced}{约化}\Term{closed-subgroup-scheme}{闭子群概形}的\Term{underlying-set}{底集}合；所有\Term{closed-immersion}{闭浸入}均\Term{proper}{固有}（EGA II, 5.4.2），故可假设 $u$ 是\Term{surjective}{满射}且 $H$ \Term{reduced}{约化}。$u$ 是\Term{surjective}{满射}，所以 $G(k)$ 在 $H$ 的\Term{closed-point}{闭点}集上传递地\Term{action}{作用}；因此对每个\Term{closed-point}{闭点} $y\in H$，$u^{-1}(y)$ 都在 $\kappa(y)$ 上\Term{proper}{固有}。由 EGA IV$_3$, 9.6.1，使 $u^{-1}(y)$ 在 $\kappa(y)$ 上不\Term{proper}{固有}的 $y\in H$ 组成\Term{locally-constructible}{局部可构造}\Term{set}{集合}；它不含\Term{closed-point}{闭点}，所以为空（参见 EGA IV$_3$, 10.3.1、10.4.7）。

\footnote{\textbf{编者注（13）：}下面将原文展开说明。}考虑 $H^0$ 的\Term{generic-point}{泛点} $\eta$。由上，\Term{fibre}{纤维} $u^{-1}(\eta)=G\times_H\operatorname{Spec}\kappa(\eta)$ 在 $\kappa(\eta)$ 上\Term{proper}{固有}。又因 $H$ \Term{reduced}{约化}，$\kappa(\eta)=\mathcal O_{H,\eta}$。当 $V$ 遍历 $H^0$ 的\Term{nonempty}{非空}\Term{open-subset}{开集}时，$\mathcal O_{H,\eta}$ 是各 $\mathcal O_H(V)$ 的\Term{inductive-limit}{归纳极限}；由 EGA IV$_3$, 8.1.2 a) 和 8.10.5 (xii)，存在 $H^0$ 的\Term{nonempty}{非空}\Term{open-subset}{开集} $V$，使 $u$ 在 $V$ 上方的\Term{restriction}{限制}\Term{proper}{固有}。显然，$g\in G(k)$ 所给出的各 $g\cdot V$ 构成 $H$ 的\Term{open-covering}{开覆盖}，且 $u$ 在每个 $g\cdot V$ 上方的\Term{restriction}{限制}都\Term{proper}{固有}。故 $u$ \Term{proper}{固有}（参见 EGA II, 5.4.1）。
\end{proof}
\begin{SGAresult}{推论}{1.4.1}
设 $A$ 是 Artin \Term{local-ring}{局部环}，$u:G\to H$ 是\Term{locally-finite-type}{局部有限型} $A$-群之间的\Term{group-morphism}{群态射}。下列断言\Term{equivalence}{等价}：
\begin{enumerate}
\item[(i)] $u$ \Term{locally-quasi-finite}{局部拟有限}；
\item[(ii)] $u$ 在某点\Term{quasi-finite}{拟有限}；
\item[(iii)] $\operatorname{Ker}u$ \Term{discrete}{离散}；
\item[(iv)] $u$ 在 $G$ 的每个\Term{connected-component}{连通分支}上的\Term{restriction}{限制}都\Term{finite-morphism}{有限}。
\end{enumerate}
若 $u$ 还\Term{quasi-compact}{拟紧}，则这些断言也\Term{equivalence}{等价}于：
\begin{enumerate}\item[(v)] $u$ \Term{finite-morphism}{有限}。\end{enumerate}
\end{SGAresult}
\begin{proof}
显然 (iv) \Term{pushout}{推出} (iii)，(iii) \Term{pushout}{推出} (ii)（EGA I, 6.4.4）；而当 $u$ \Term{quasi-compact}{拟紧}时，(iv) 与 (v) \Term{equivalence}{等价}。1.3 已经证明 (i) 与 (ii) \Term{equivalence}{等价}。

最后证明 (i) \Term{pushout}{推出} (iv)。设 $C$ 是 $G$ 的一个\Term{connected-component}{连通分支}。$C$ 在 $A$ 上\Term{finite-type}{有限型}（第 VI$_A$ 讲 2.4.1），而 $G,H$ \Term{separated}{分离}（第 VI$_A$ 讲 0.2），故由 EGA I, 5.5.1、6.3.4，$u$ 在 $C$ 上的\Term{restriction}{限制} $u'$ \Term{separated}{分离}且\Term{finite-type}{有限型}。\footnote{\textbf{编者注（14）：}下面将原文展开说明。}$u'$ 的\Term{fibre}{纤维}\Term{discrete}{离散}，故 $u'$ \Term{quasi-finite}{拟有限}（参见 EGA II, 6.2.2）。\Term{quasi-finite}{拟有限}且\Term{proper}{固有}的\Term{morphism}{态射}总是\Term{finite-morphism}{有限}（参见 EGA III$_1$, 4.4.2），所以只需证明 $u'$ \Term{universally-closed}{普遍闭}。

为此，可假设 $A$ 等于其\Term{residue-field}{剩余域} $k$。再由 (fpqc) \Term{descent}{下降}（参见 EGA IV$_2$, 2.6.4），只需证明 $u'\otimes_k\overline k$ \Term{universally-closed}{普遍闭}，其中 $\overline k$ 是 $k$ 的\Term{algebraic-closure}{代数闭包}。又因 $C$ 在 $k$ 上\Term{finite-type}{有限型}，$C\otimes_k\overline k$ 是 $G\otimes_k\overline k$ 的\Term{finite}{有限}个\Term{connected-component}{连通分支} $C'_1,\ldots,C'_n$ 的和，
\SourcePage{341}
所以只需逐个 $C'_i$ 证明。这便归结为 $k$ \Term{algebraically-closed-field}{代数闭}的情形。

取 $C$ 的\Term{closed-point}{闭点} $g$。若 $u^0:G^0\to H$ 是 $u$ 在 $G^0$ 上的\Term{restriction}{限制}，则
\[
u'=r_{u(g)}\circ u^0\circ r_g^{-1},
\]
其中 $r_g$ 是以 $g$ 作\Term{right-translation}{右平移}，所以只需证明 $u^0$ \Term{proper}{固有}。由假设 $u$ \Term{locally-quasi-finite}{局部拟有限}，故\Term{fibre}{纤维} $\operatorname{Ker}u$ \Term{discrete}{离散}且\Term{nonempty}{非空}；又已知 $u^0$ \Term{finite-type}{有限型}，所以\Term{fibre}{纤维} $\operatorname{Ker}u^0$ 是\Term{finite}{有限}的（参见 EGA II, 6.2.2），因而\Term{proper}{固有}且\Term{nonempty}{非空}。由 1.4，$u^0$ \Term{proper}{固有}。
\end{proof}
\begin{SGAresult}{推论}{1.4.2}
设 $A$ 是 Artin \Term{local-ring}{局部环}，$u:G\to H$ 是\Term{locally-finite-type}{局部有限型} $A$-群之间的\Term{quasi-compact}{拟紧}\Term{group-morphism}{群态射}。下列断言\Term{equivalence}{等价}：\footnote{\textbf{编者注（15）：}可以去掉 $G,H$ \Term{locally-finite-type}{局部有限型}的假设，因为由 [Per76, 4.2.4]，域 $k$ 上\Term{group-scheme}{群概形}之间的每个\Term{quasi-compact}{拟紧}\Term{monomorphism}{单态射} $u:G\to H$ 都是\Term{closed-immersion}{闭浸入}；若 $G,H$ 都\Term{affine}{仿射}，另见 [DG70, III.3.7.2 b)]（此时每个\Term{morphism}{态射} $G\to H$ 都\Term{affine}{仿射}（EGA II, 1.6.2 (v)），从而\Term{quasi-compact}{拟紧}）。}
\begin{enumerate}
\item[(i)] $u$ 是\Term{closed-immersion}{闭浸入}；
\item[(ii)] $u$ 是\Term{monomorphism}{单态射}；
\item[(iii)] $\operatorname{Ker}u$ \Term{trivial-group}{平凡}，即与单位 $k$-群\Term{isomorphism}{同构}。
\end{enumerate}
特别地，$H$ 的每个\Term{group-subscheme}{子群概形}\footnote{\textbf{编者注（16）：}这一特殊情形亦见第 VI$_A$ 讲 0.5.2，那里不需要\Term{finiteness}{有限性}假设。}都是闭的。
\end{SGAresult}
\begin{proof}
显然 (i) \Term{pushout}{推出} (ii)；考虑 $G,H$ 各自所\Term{representable}{表示}的\Term{functor}{函子}，立刻可知 (ii) 与 (iii) \Term{equivalence}{等价}。最后，若 $\operatorname{Ker}u$ 是单位 $k$-群，则它是一个\Term{proper}{固有}\Term{nonempty}{非空}\Term{fibre}{纤维}，故由 1.4，$u$ 是\Term{proper}{固有}\Term{monomorphism}{单态射}；$H$ \Term{locally-noetherian}{局部诺特}，所以 $u$ 还是\Term{finite-presentation}{有限呈示}的（EGA IV$_1$, 1.6.1），从而是\Term{closed-immersion}{闭浸入}（EGA IV$_3$, 8.11.5）。

最后的断言来自以下事实：$H$ \Term{locally-noetherian}{局部诺特}，所以每个\Term{immersion}{浸入} $G\to H$ 都\Term{quasi-compact}{拟紧}（EGA I, 6.6.4）。
\end{proof}
\begin{SGAresult}{反例}{1.4.3}
设 $k$ 是\Term{characteristic}{特征}为 $0$ 的域，$G$ 是\Term{constant-group}{常值} $k$-群 $\mathbf Z$，$H=\mathbf G_{a,k}$。若 $u:G\to H$ 是非零 $k$-\Term{group-morphism}{群态射}，则 $\operatorname{Ker}u=0$，但 $u$ 不是\Term{closed-immersion}{闭浸入}（参见 1.2.2）。
\end{SGAresult}
下面将用到本应列入第 VI$_A$ 讲的两个结果。
\begin{SGAresult}{\Term{lemma}{引理}}{1.5}
设 $k$ 是域，$G$ 是\Term{locally-finite-type}{局部有限型} $k$-群。则：
\begin{enumerate}
\item[(i)] $G$ 的各\Term{irreducible-component}{不可约分支}\Term{dimension}{维数}都相同。
\item[(ii)] 对每个 $g\in G$，有 $\dim_gG=\dim G$。
\end{enumerate}
\end{SGAresult}
\SourcePage{342}
\begin{proof}
\footnote{\textbf{编者注（17）：}(ii)$\Rightarrow$(i) 是一般事实（参见 EGA 0$_{\mathrm{IV}}$, 14.1.6）；(i)$\Rightarrow$(ii) 则由以下事实\Term{pushout}{推出}：若 $X$ 是域上的\Term{irreducible}{不可约}\Term{finite-type}{有限型}\Term{scheme}{概形}，则对每个\Term{nonempty}{非空}\Term{open-subset}{开集} $U\subset X$，$\dim X=\dim U$。故这里关键是 (i)，它已在第 VI$_A$ 讲 2.4.1 的增补中证明。我们据此修改了\Term{lemma}{引理}的陈述和证明。}断言 (i) 已在第 VI$_A$ 讲 2.4.1 中证明，(ii) 由此而来。事实上，取 $g\in G$，令 $C$ 是包含 $g$ 的\Term{connected-component}{连通分支}。按定义（EGA 0$_{\mathrm{IV}}$, 14.1.2），$\dim_gG$ 是各整数 $\dim U$ 的\Term{infimum}{下确界}，其中 $U$ 遍历 $g$ 的\Term{open-neighborhood}{开邻域}；故 $\dim_gG=\dim U_0$，其中 $U_0$ 是某个这样的\Term{neighborhood}{邻域}，并可假设 $U_0\subset C$（因为 $V\subset U$ 时 $\dim V\leq\dim U$）。$C$ 在 $k$ 上\Term{irreducible}{不可约}且\Term{finite-type}{有限型}（第 VI$_A$ 讲 2.4.1），所以由 EGA IV$_2$, 5.2.1，
\[
\dim U_0=\dim C=\operatorname{trdeg}_k\kappa(\xi),
\]
其中 $\xi$ 是 $C$ 的\Term{generic-point}{泛点}。故 $\dim_gG=\dim C=\dim G$。
\end{proof}
\begin{SGAresult}{命题}{1.6}
\footnote{\textbf{编者注（18）：}将陈述中的“沿\Term{unit-section}{单位截面}”改为“在\Term{unit-section}{单位截面}的每一点”；又在证明末尾明确列出了所引用的 EGA IV$_4$ 结果。}设 $S$ 是\Term{characteristic}{特征}为零的\Term{scheme}{概形}，$G$ 是 $S$-\Term{group-scheme}{群概形}，且在\Term{identity-section}{单位截面} $\varepsilon(S)$ 的每一点都在 $S$ 上\Term{locally-finite-presentation}{局部有限呈示}。则 $G$ 在\Term{identity-section}{单位截面}的每一点都在 $S$ 上\Term{smooth}{光滑}，当且仅当 $\mathcal O_S$-模 $\omega_{G/S}=\varepsilon^*(\Omega^1_{G/S})$（称为 $G$ 的\Term{identity-section}{单位截面}处的\Term{conormal-module}{余法模}）\Term{locally-free}{局部自由}。
\end{SGAresult}
回顾，称\Term{scheme}{概形} $S$ 的\Term{characteristic}{特征}为零，是指对每个\Term{closed-point}{闭点} $x\in S$，域 $\kappa(x)$ 的\Term{characteristic}{特征}都是零。

还要回顾（第 II 讲 4.11），若 $\pi:G\to S$ 为\Term{structure-morphism}{结构态射}，则 $\Omega^1_{G/S}=\pi^*(\omega_{G/S})$，故 $\mathcal O_S$-模 $\omega_{G/S}$ \Term{locally-free}{局部自由}，\Term{equivalence}{等价}于 $\mathcal O_G$-模 $\Omega^1_{G/S}$ \Term{locally-free}{局部自由}。
\begin{proof}
若 $\varepsilon(S)$ 有一个在 $S$ 上\Term{smooth}{光滑}的\Term{open-neighborhood}{开邻域} $U$，则由 EGA IV$_4$, 17.2.3，$\Omega^1_{U/S}$ 是\Term{finite-type}{有限型}\Term{locally-free}{局部自由}模，$\omega_{G/S}=\varepsilon^*(\Omega^1_{U/S})$ 亦然。

反之，若 $\omega_{G/S}$ \Term{locally-free}{局部自由}，则 $\Omega^1_{G/S}=\pi^*(\omega_{G/S})$ 也\Term{locally-free}{局部自由}。$S$ 的\Term{characteristic}{特征}为 $0$，所以由\Term{jacobian-criterion}{Jacobi 判据}（EGA IV$_4$, 16.12.2），$G$ 在 $S$ 上\Term{differentially-smooth}{微分光滑}。由 EGA IV$_4$, 17.12.5，$G$ 在\Term{identity-section}{单位截面}的每一点都在 $S$ 上\Term{smooth}{光滑}。
\end{proof}
\begin{SGAresult}{推论（Cartier）}{1.6.1}
设 $k$ 是\Term{characteristic}{特征}为零的域，则每个\Term{locally-finite-type}{局部有限型} $k$-群都在 $k$ 上\Term{smooth}{光滑}。
\end{SGAresult}
\begin{proof}
此时 $k$-模 $\omega_{G/k}$ 显然\Term{locally-free}{局部自由}，故由 1.6，$G$ 在\Term{identity-element}{单位点} $e$ 处在 $k$ 上\Term{smooth}{光滑}；再由 1.3.1，$G$ 在 $k$ 上\Term{smooth}{光滑}。
\end{proof}

\SGAsection{2}{局部有限表现群及其态射的“开放性质”}
\textbf{2.0。}以下 $S$ 总表示任意\Term{scheme}{概形}；$S$-\Term{group-scheme}{群概形}简称为 $S$-群。对 $S$-群 $G$，以 $\varepsilon$ 表示\Term{identity-section}{单位截面}，$c:G\to G$ 表示\Term{inversion-morphism}{取逆态射}，$\mu$ 表示\Term{multiplication-morphism}{乘法态射}
\SourcePage{343}
$G\times_SG\to G$。对每个 $S$-\Term{scheme}{概形} $X$，以 $\pi$ 或 $\pi_X$ 表示\Term{structure-morphism}{结构态射} $X\to S$。

设 $\mathscr P(u)$ 是 $S$-\Term{scheme}{概形}\Term{morphism}{态射} $u:X\to Y$ 的一种性质。若 $u$ 满足 $\mathscr P(u)$ 时，对任意 $S$-\Term{morphism}{态射} $Y'\to Y$，$u_{(Y')}$ 也满足此性质，就称 $\mathscr P(u)$ 对\Term{base-change}{基变换}\Term{stable}{稳定}。若 $\mathscr P(u)$ 满足以下两个条件，则称它对于\Term{topology}{拓扑} $\mathcal T$ 具有\Term{local-character}{局部性}（参见第 IV 讲第 4、6 节）：
\begin{enumerate}
\item[a)] $\mathscr P(u)$ 对\Term{base-change}{基变换}\Term{stable}{稳定}；
\item[b)] 只要有一个对于\Term{topology}{拓扑} $\mathcal T$ \Term{covering-family}{覆盖}的 $S$-\Term{morphism}{态射}族 $Y_i\to Y$，使每个 $u_{(Y_i)}$ 都满足 $\mathscr P(u)$，则 $u$ 也满足 $\mathscr P(u)$。
\end{enumerate}
\begin{SGAresult}{命题}{2.1}
设 $\mathscr P(u)$ 是 $S$-\Term{scheme}{概形}\Term{morphism}{态射}的一种对于 (fpqc) \Term{topology}{拓扑}具有\Term{local-character}{局部性}的性质，$u:G\to H$ 是 $S$-\Term{group-morphism}{群态射}。假设 $G$ 在 $S$ 上\Term{flat}{平坦}且\Term{universally-open}{普遍开}。

令 $W$ 是 $H$ 中使 $u$ 在其上方满足 $\mathscr P(u)$ 的最大\Term{open-subset}{开集}，$V=u^{-1}(W)$。则 $U=\pi_G(V)$ 是 $S$ 的\Term{open-subset}{开集}，$V$ 是 $G|_U$ 的\Term{open-subgroup-scheme}{开子群概形}。
\end{SGAresult}
\begin{proof}
$\mathscr P(u)$ 对 Zariski \Term{topology}{拓扑}具有\Term{local-character}{局部性}，故所述最大\Term{open-subset}{开集} $W$ 存在。$\pi_G$ \Term{universally-open}{普遍开}，所以 $\pi_G(V)$ 是 $S$ 的\Term{open-subset}{开集} $U$。只需证明 $V$ 是 $G|_U$ 的\Term{group-subscheme}{子群概形}，故可假设 $U=S$。

置 $G'=G\times_SV$、$H'=H\times_SV$、$V'=V\times_SV$、$W'=W\times_SV$、$u'=u_{(V)}$。令 $W'_1$ 是 $H'$ 中使 $u'$ 在其上方满足 $\mathscr P(u)$ 的最大\Term{open-subset}{开集}。$V$ 在 $S$ 上\Term{flat}{平坦}且\Term{universally-open}{普遍开}，所以 $H'$ 在 $H$ 上亦然；下面的\Term{lemma}{引理} 2.1.1 说明 $W'_1=W'$。考虑 $G'$（相应地 $H'$）的 $V$-\Term{scheme}{概形}\Term{automorphism}{自同构} $a$（相应地 $b$），它是以\Term{diagonal-section}{对角截面} $\delta$ 的逆元（相应地 $u(\delta)$ 的逆元）作\Term{right-translation}{右平移}，由下式定义：对任意\Term{morphism}{态射} $T\to S$，$g\in G(T)$、$v\in V(T)$、$h\in H(T)$，
\[
a(g,v)=(g\cdot v^{-1},v),\qquad b(h,v)=(h\cdot u(v^{-1}),v).
\]
显然 $u'\circ a=b\circ u'$，故 $W'$ 对 $b$ \Term{stable}{稳定}，从而 $V'$ 对 $a$ \Term{stable}{稳定}。这说明 $V$ 是 $G$ 的\Term{group-subscheme}{子群概形}。
\end{proof}
\begin{SGAresult}{\Term{lemma}{引理}}{2.1.1}
设 $\mathscr P(u)$ 是 $S$-\Term{morphism}{态射}的一种对于 (fpqc) \Term{topology}{拓扑}具有\Term{local-character}{局部性}的性质。考虑 $S$-\Term{scheme}{概形}\Term{morphism}{态射}的\Term{cartesian-square}{笛卡尔方块}
\[
\begin{tikzcd}
X'\arrow[r,"f'"]\arrow[d]&Y'\arrow[d,"g"]\\
X\arrow[r,"f"]&Y.
\end{tikzcd}
\]
其中 $g$ \Term{flat}{平坦}且为\Term{open-morphism}{开态射}。令 $W$（相应地 $W'_1$）是 $Y$（相应地 $Y'$）中使 $f$（相应地 $f'=f_{(Y')}$）在其上方满足 $\mathscr P(u)$ 的最大\Term{open-subset}{开集}。则 $W'_1=W\times_YY'$。
\end{SGAresult}
\SourcePage{344}
\begin{proof}
置 $W'=W\times_YY'$。因 $\mathscr P(u)$ 对\Term{base-change}{基变换}\Term{stable}{稳定}，$W'\subset W'_1$。因 $g$ 是\Term{open-morphism}{开态射}，$W_1=g(W'_1)$ 是 $Y$ 的\Term{open-subset}{开集}。置 $V_1=f^{-1}(W_1)$、$V'_1=V_1\times_{W_1}W'_1$，显然 $V'_1=f'^{-1}(W'_1)$。$g$ \Term{flat}{平坦}且为\Term{open-morphism}{开态射}，故由它诱导的\Term{morphism}{态射} $W'_1\to W_1=g(W'_1)$ \Term{faithfully-flat}{忠实平坦}且为\Term{open-morphism}{开态射}，因而是 (fpqc) \Term{topology}{拓扑}的\Term{covering}{覆盖}（参见第 IV 讲 6.3.1 (iv)）。由 $f'$ 诱导的\Term{morphism}{态射} $V'_1\to W'_1$ 满足 $\mathscr P(u)$，所以由 $f$ 诱导的\Term{morphism}{态射} $V_1\to W_1$ 也满足此性质。于是 $W_1\subset W$，且
\[
W'_1\subset g^{-1}(W_1)\subset g^{-1}(W)=W',
\]
故 $W'=W'_1$。
\end{proof}
\textbf{注记 2.1.2。}许多\Term{morphism}{态射}性质对于 (fpqc) \Term{topology}{拓扑}具有\Term{local-character}{局部性}，例如\Term{flat}{平坦}、（\Term{universally-open}{普遍}）\Term{open-morphism}{开}、（\Term{locally-finite-type}{局部}）\Term{finite-type}{有限型}、\Term{finite-presentation}{有限呈示}、\Term{quasi-finite}{拟有限}（参见 EGA IV$_2$, 2.5.1、2.6.1、2.7.1），以及\Term{smooth}{光滑}、\Term{etale}{平展}、\Term{unramified}{非分歧}（EGA IV$_4$, 17.7.3）。

2.1 的证明实际上只用了\Term{flat}{平坦}\Term{morphism}{态射}所给出的\Term{base-change}{基变换}；所以，只要一种性质满足 2.0 中对于 (fpqc) \Term{topology}{拓扑}的条件 b)，并对\Term{flat}{平坦}\Term{morphism}{态射}的\Term{base-change}{基变换}\Term{stable}{稳定}，这个命题就适用（例如\Term{quasi-compact}{拟紧}且\Term{dominant}{支配}）。

当然，对于比 Zariski \Term{topology}{拓扑}更细的\Term{topology}{拓扑} $\mathcal T$ 上具有\Term{local-character}{局部性}的性质，也能陈述类似命题；这时应要求 $\pi_G$ \Term{universally-open}{普遍开}，并且对于 $\mathcal T$ 是\Term{covering}{覆盖}。

特别地，若 $G$ 在 $S$ 上\Term{flat}{平坦}且\Term{locally-finite-presentation}{局部有限呈示}，则对于那些在\Term{flat}{平坦}\Term{locally-finite-presentation}{局部有限呈示}\Term{morphism}{态射}的\Term{base-change}{基变换}下\Term{stable}{稳定}、并满足 2.0 中关于 (fpqc) \Term{topology}{拓扑}的条件 b) 的性质，也有类似陈述（例如\Term{regular}{正则}、\Term{reduced}{约化}、Cohen--Macaulay 等，参见 EGA IV$_2$, 6.8）。
\begin{SGAresult}{命题}{2.2}
设 $G,H$ 是两个 $S$-群，$u:G\to H$ 是 $S$-\Term{group-morphism}{群态射}。则：\footnote{\textbf{编者注（19）：}在 (i) 中，\Term{open-subset}{开集} $V$ 由 $G$ 中 $u$ 为\Term{smooth}{光滑}、相应地\Term{etale}{平展}、相应地\Term{finite-presentation}{有限呈示}且\Term{flat}{平坦}的所有点组成，参见编者注（6）。但在 (ii) 中，$V$ 是 $u$ \Term{universally-open}{普遍开}的点集 $E$ 中所含的最大\Term{open-subset}{开集}；$E$ 未必是\Term{open-subset}{开集}，如下例所示（EGA IV$_3$, 14.1.3 (i)）：设 $k$ 是域，$H=S=\operatorname{Spec}A$，其中 $A=k[x]$，$G$ 为 $S$-群 $\operatorname{Spec}A[y]/(xy)$；则 $E$ 就是 $G$ 的\Term{identity-section}{单位截面}，它不是\Term{open-subset}{开集}。}
\begin{enumerate}
\item[(i)] 假设 $G,H$ 中至少一个在 $S$ 上\Term{flat}{平坦}，且至少一个在 $S$ 上\Term{locally-finite-presentation}{局部有限呈示}。令 $V$ 是 $G$ 中使 $u|_V$ \Term{flat}{平坦}且\Term{locally-finite-presentation}{局部有限呈示}（相应地\Term{smooth}{光滑}、\Term{etale}{平展}）的最大\Term{open-subset}{开集}。则 $U=\pi_V(V)$ 是 $S$ 的\Term{open-subset}{开集}，$V$ 是 $G|_U$ 的\Term{open-subgroup-scheme}{开子群概形}。
\item[(ii)] 假设 $G,H$ 中至少一个在 $S$ 上\Term{universally-open}{普遍开}。令 $V$ 是 $G$ 中使 $u|_V$ \Term{universally-open}{普遍开}的最大\Term{open-subset}{开集}。则 $U=\pi_V(V)$ 是 $S$ 的\Term{open-subset}{开集}，$V$ 是 $G|_U$ 的\Term{open-subgroup-scheme}{开子群概形}。
\end{enumerate}
\end{SGAresult}
\begin{proof}
先证 (i)。首先证明 $\pi_G$ 在 $V$ 上的\Term{restriction}{限制} $\pi_V$ \Term{flat}{平坦}且\Term{locally-finite-presentation}{局部有限呈示}：
\begin{enumerate}
\item[a)] 若 $\pi_G$ \Term{flat}{平坦}（相应地\Term{locally-finite-presentation}{局部有限呈示}），则 $\pi_V$ 亦然。
\end{enumerate}
\SourcePage{345}
\begin{enumerate}
\item[b)] 若 $\pi_H$ \Term{flat}{平坦}（相应地\Term{locally-finite-presentation}{局部有限呈示}），则 $\pi_V$ 作为 $u|_V$ 与 $\pi_H$ 的复合，也有此性质。
\end{enumerate}
所以陈述中所考虑的四种情形下，$\pi_V$ 都\Term{flat}{平坦}且\Term{locally-finite-presentation}{局部有限呈示}，故\Term{universally-open}{普遍开}（EGA IV$_2$, 2.4.6）。置 $U=\pi_V(V)$，则 $U$ 是 $S$ 的\Term{open-subset}{开集}。只需证明 $V$ 是 $G|_U$ 的\Term{open-subgroup-scheme}{开子群概形}，故可假设 $U=S$。

置 $G'=G\times_SV$、$H'=H\times_SV$、$V'=V\times_SV$、$u'=u_{(V)}$。$V$ 在 $S$ 上\Term{flat}{平坦}且\Term{locally-finite-presentation}{局部有限呈示}，所以 $H'$ 在 $H$ 上亦然。由 EGA IV$_4$, 17.7.4，$V'$ 是 $G'$ 中使 $u'|_{V'}$ \Term{flat}{平坦}且\Term{locally-finite-presentation}{局部有限呈示}（相应地\Term{smooth}{光滑}、\Term{etale}{平展}）的最大\Term{open-subset}{开集}。将\Term{automorphism}{自同构} $a,b$ 按 2.1 的证明定义，显然 $V'$ 对 $a$ \Term{stable}{稳定}，故 $V$ 是 $G$ 的\Term{group-subscheme}{子群概形}。

再证 (ii)。$\pi_G$ 在 $V$ 上的\Term{restriction}{限制} $\pi_V$ 是\Term{universally-open}{普遍开}的：若 $\pi_G$ \Term{universally-open}{普遍开}，这立即成立；若 $\pi_H$ \Term{universally-open}{普遍开}，则因 $\pi_V$ 是 $u|_V$ 与 $\pi_H$ 的复合而成立。置 $U=\pi_V(V)$，则 $U$ 是 $S$ 的\Term{open-subset}{开集}。只需证明 $V$ 是 $G|_U$ 的\Term{open-subgroup-scheme}{开子群概形}，故可假设 $U=S$。

如上置 $G'=G\times_SV$、$H'=H\times_SV$、$V'=V\times_SV$、$u'=u_{(V)}$。此时 $\pi_V:V\to S$ 是\Term{surjective}{满射}且\Term{universally-open}{普遍开}，$G'\to G$ 亦然。所以由 EGA IV$_3$, 14.3.4 (i)、(ii)，$V'$ 是 $G'$ 中使 $u'|_{V'}$ \Term{universally-open}{普遍开}的最大\Term{open-subset}{开集}。将\Term{automorphism}{自同构} $a,b$ 定义如前，显然 $V'$ 对 $a$ \Term{stable}{稳定}，故 $V$ 是 $G$ 的\Term{group-subscheme}{子群概形}。
\end{proof}
\begin{SGAresult}{推论}{2.3}
设 $G$ 是 $S$-群，$V$ 是 $G$ 中在 $S$ 上\Term{flat}{平坦}且\Term{locally-finite-presentation}{局部有限呈示}（相应地\Term{smooth}{光滑}、\Term{etale}{平展}、\Term{universally-open}{普遍开}）的最大\Term{open-subset}{开集}。则 $U=\pi_V(V)$ 是 $S$ 的\Term{open-subset}{开集}，$V$ 是 $G|_U$ 的\Term{open-subgroup-scheme}{开子群概形}。
\end{SGAresult}
\begin{proof}
将 2.2 应用于 $H$ 为单位 $S$-群、$u$ 为唯一 $S$-\Term{group-morphism}{群态射} $G\to H$ 的情形即可；此时 $\pi_H$ 是\Term{isomorphism}{同构}，而 $\pi_G=\pi_H\circ u$。
\end{proof}
\begin{SGAresult}{推论}{2.4}
设 $G$ 是 $S$-群。若\Term{identity-section}{单位截面}有一个\Term{neighborhood}{邻域} $X$，具有下列性质之一（或若干个）：
\begin{quote}
$X$ 在 $S$ 上\Term{flat}{平坦}且\Term{locally-finite-presentation}{局部有限呈示}（相应地\Term{smooth}{光滑}、\Term{etale}{平展}、\Term{universally-open}{普遍开}），
\end{quote}
则 $G$ 有一个具有相同性质的\Term{open-subgroup-scheme}{开子群概形}。
\end{SGAresult}
\begin{proof}
用 2.3 即可；注意按 2.2 的记号，这里 $\varepsilon(S)\subset V$，所以 $U=S$。
\end{proof}
\begin{SGAresult}{命题}{2.5}
\footnote{\textbf{编者注（20）：}展开说明了陈述和证明，并将 (i) 中对 $H$ 的假设从“\Term{finite-presentation}{有限呈示}”减弱为“\Term{finite-type}{有限型}”。}设 $u:G\to H$ 是 $S$-\Term{group-morphism}{群态射}。
\SourcePage{346}
\begin{enumerate}
\item[(i)] 假设 $G$（相应地 $H$）在其\Term{identity-section}{单位截面}的点处在 $S$ 上\Term{finite-presentation}{有限呈示}且\Term{flat}{平坦}（相应地\Term{finite-type}{有限型}）。则由使 $u_s$ \Term{flat}{平坦}（相应地\Term{smooth}{光滑}、\Term{etale}{平展}）的 $s\in S$ 组成的\Term{set}{集合}
\[
U_{\mathrm{plat}}\supset U_{\mathrm{lisse}}\supset U_{\text{ét}}
\]
都是 $S$ 中的\Term{open-subset}{开集}。

若 $G$（相应地 $H$）还在 $S$ 上\Term{locally-finite-presentation}{局部有限呈示}且\Term{flat}{平坦}（相应地\Term{locally-finite-type}{局部有限型}），则 $G$ 中 $u$ 为\Term{flat}{平坦}（相应地\Term{smooth}{光滑}、\Term{etale}{平展}）的点集 $V_{\mathrm{plat}}$（相应地 $V_{\mathrm{lisse}},V_{\text{ét}}$），等于 $U_{\mathrm{plat}}$（相应地 $U_{\mathrm{lisse}},U_{\text{ét}}$）在 $\pi_G$ 下的\Term{inverse-image}{逆像}。
\item[(ii)] 假设对每个 $s\in S$，$G_s$ 在 $\kappa(s)$ 上\Term{locally-finite-type}{局部有限型}（若 $G$ 在\Term{identity-section}{单位截面}的点处在 $S$ 上\Term{finite-type}{有限型}，则由 1.3.1.1，此条件成立）；又假设 $u$ 在 $G$ 的\Term{identity-section}{单位截面}的点处\Term{locally-finite-type}{局部有限型}（相应地\Term{locally-finite-presentation}{局部有限呈示}）。则由使 $u_s$ \Term{locally-quasi-finite}{局部拟有限}（相应地\Term{unramified}{非分歧}）的 $s\in S$ 组成的\Term{set}{集合}
\[
U_{\mathrm{l.q.f.}}\supset U_{\mathrm{n.r.}}
\]
都是 $S$ 中的\Term{open-subset}{开集}。

若 $u$ 还\Term{locally-finite-type}{局部有限型}（相应地\Term{locally-finite-presentation}{局部有限呈示}），则 $G$ 中 $u$ 为\Term{quasi-finite}{拟有限}（相应地\Term{unramified}{非分歧}）的点集 $V_{\mathrm{l.q.f.}}$（相应地 $V_{\mathrm{n.r.}}$），等于 $U_{\mathrm{l.q.f.}}$（相应地 $U_{\mathrm{n.r.}}$）在 $\pi_G$ 下的\Term{inverse-image}{逆像}。
\end{enumerate}
\end{SGAresult}
\begin{proof}
先证 (i)。首先注意，由 1.3.1.1，对每个 $s\in S$，$G_s$ 在 $\kappa(s)$ 上\Term{locally-finite-type}{局部有限型}。令 $Y$ 是 $H$ 中 $\pi_H$ 为\Term{finite-type}{有限型}的点组成的\Term{open-subset}{开集}，$X_1$ 是 $u^{-1}(Y)$ 中 $\pi_G$ 为\Term{finite-presentation}{有限呈示}的点组成的\Term{open-subset}{开集}。由 EGA IV$_3$, 11.3.1，$X_1$ 中 $\pi_G$ 为\Term{finite-presentation}{有限呈示}且\Term{flat}{平坦}的点集 $X$ 在 $X_1$ 中开，因而在 $G$ 中开。

以 $u_X:X\to Y$ 表示由 $u$ 诱导的\Term{morphism}{态射}。$Y$（相应地 $X$）包含 $H$（相应地 $G$）的\Term{identity-section}{单位截面}，故 $\pi_G$ 在 $X$ 上的\Term{restriction}{限制} $\pi_X$ 是\Term{surjective}{满射}，由 $\varepsilon_G$ 诱导的\Term{morphism}{态射} $S\to X$ 是 $X$ 的一个 $S$-\Term{section}{截面}。

考虑 $X$ 中 $u_X$ 为\Term{flat}{平坦}（相应地\Term{smooth}{光滑}、\Term{etale}{平展}）的点集
\[
W_{\mathrm{plat}}\supset W_{\mathrm{lisse}}\supset W_{\text{ét}}.
\]
显然 $W_{\mathrm{lisse}},W_{\text{ét}}$ 在 $X$ 中开，参见编者注（6）。另一方面，$\pi_Y$ \Term{locally-finite-type}{局部有限型}且 $\pi_X=\pi_Y\circ u_X$ \Term{locally-finite-presentation}{局部有限呈示}，故由 EGA IV$_1$, 1.4.3 (v)，$u_X$ \Term{locally-finite-presentation}{局部有限呈示}。因此 $u_X$ 的\Term{flat-locus}{平坦点集} $W_X$ 在 $X$ 中开（EGA IV$_3$, 11.3.1）。

取 $x\in X$，置 $s=\pi_X(x)$。由 EGA IV$_2$, 11.3.10（再结合 EGA IV$_4$, 17.5.1，相应地 17.6.1），$x$ 属于 $W_{\mathrm{plat}}$（相应地 $W_{\mathrm{lisse}},W_{\text{ét}}$），当且仅当 $u_s$ 在 $x$ 处\Term{flat}{平坦}（相应地\Term{smooth}{光滑}、\Term{etale}{平展}）；由 1.3，这又\Term{equivalence}{等价}于 $u_s$ \Term{flat}{平坦}（相应地\Term{smooth}{光滑}、\Term{etale}{平展}）。

所以，对 $\flat=\mathrm{plat},\mathrm{lisse},\text{ét}$，有
\[
U_\flat=\varepsilon_G^{-1}(W_\flat),\qquad W_\flat=\pi_X^{-1}(U_\flat).
\]
$W_\flat$ 在 $X$ 中开，故在 $G$ 中开；第一个等式说明 $U_\flat$ 在 $S$ 中开。

(i) 的第二个断言也由此得到，因为这时 $Y=H$、$X=G$，且 $V_\flat=W_\flat=\pi_G^{-1}(U_\flat)$。
\SourcePage{347}
再证 (ii)。令 $V_{\mathrm{l.t.f.}}\supset V_{\mathrm{l.q.f.}}$ 是 $G$ 中 $u$ 为\Term{finite-type}{有限型}（相应地\Term{quasi-finite}{拟有限}）的点所成的\Term{open-subset}{开集}。置 $X=V_{\mathrm{l.t.f.}}$，并以 $u_X$ 表示 $u$ 在 $X$ 上的\Term{restriction}{限制}。由假设 $X$ 包含 $\varepsilon(S)$。取 $x\in X$，置 $s=\pi_X(x)$。

若 $u$ 在 $x$ 处\Term{quasi-finite}{拟有限}，则经\Term{base-change}{基变换}后，$u_s$ 在 $x$ 处\Term{quasi-finite}{拟有限}。$G_s$ 按假设\Term{locally-finite-type}{局部有限型}，所以由 1.3.1，$u_s$ \Term{locally-quasi-finite}{局部拟有限}。

反之，若 $u_s$ \Term{locally-quasi-finite}{局部拟有限}，则 $u_X^{-1}(u_X(x))\subset u_s^{-1}(u_s(x))$ 是\Term{finite}{有限}\Term{set}{集合}。$u_X:X\to H$ \Term{locally-finite-type}{局部有限型}，故由 Chevalley \Term{semicontinuity}{半连续性}\Term{theorem}{定理}（EGA IV$_3$, 13.1.3），$x$ 在 $X$ 中有一个\Term{open-neighborhood}{开邻域} $W$，使对每个 $x'\in W$，\Term{fibre}{纤维} $u_X^{-1}(u_X(x'))$ 都\Term{discrete}{离散}。所以由 EGA II, 6.2.2（以及 EGA III$_2$, Err$_{\mathrm{III}}$ 20），$u_X$ 在 $x$ 处\Term{quasi-finite}{拟有限}。

因此
\[
U_{\mathrm{l.q.f.}}=\varepsilon_G^{-1}(V_{\mathrm{l.q.f.}}),\qquad
V_{\mathrm{l.q.f.}}=\pi_X^{-1}(U_{\mathrm{l.q.f.}}),
\]
第一个等式说明 $U_{\mathrm{l.q.f.}}$ 在 $S$ 中开。若 $G$ 还在 $S$ 上\Term{locally-finite-type}{局部有限型}，则 $X=G$，故 $V_{\mathrm{l.q.f.}}$ 是 $U_{\mathrm{l.q.f.}}$ 在 $\pi_G$ 下的\Term{inverse-image}{逆像}。这证明了 (ii) 的前半部分。

现在考虑由 $u$ 为\Term{finite-presentation}{有限呈示}（相应地\Term{unramified}{非分歧}）的点组成的\Term{open-subset}{开集} $V_{\mathrm{l.p.f.}}\supset V_{\mathrm{n.r.}}$，并假设 $Z=V_{\mathrm{l.p.f.}}$ 包含 $G$ 的\Term{identity-section}{单位截面}。取 $z\in Z$，置 $s=\pi_Z(z)$、$h=u_Z(z)$。

若 $u$ 在 $z$ 处\Term{unramified}{非分歧}，则经\Term{base-change}{基变换}后，$u_s$ 在 $z$ 处\Term{unramified}{非分歧}；$G_s$ 按假设\Term{locally-finite-type}{局部有限型}，故由 1.3.1，$u_s$ \Term{unramified}{非分歧}。

反之，若 $u_s$ 在 $z$ 处\Term{unramified}{非分歧}，则\Term{fibre}{纤维} $u_s^{-1}(h)=u^{-1}(h)$ 在 $z$ 处在 $\kappa(h)$ 上\Term{unramified}{非分歧}。$Z$ 是 $G$ 的\Term{open-subset}{开集}，故\Term{fibre}{纤维} $u_Z^{-1}(h)=u_s^{-1}(h)\cap Z$ 在 $z$ 处在 $\kappa(h)$ 上\Term{unramified}{非分歧}。$u_Z$ \Term{locally-finite-presentation}{局部有限呈示}，故由 EGA IV$_4$, 17.4.1，$u_Z$ 在 $z$ 处\Term{unramified}{非分歧}。

因此
\[
U_{\mathrm{n.r.}}=\varepsilon_G^{-1}(V_{\mathrm{n.r.}}),\qquad
V_{\mathrm{n.r.}}=\pi_Z^{-1}(U_{\mathrm{n.r.}}),
\]
第一个等式说明 $U_{\mathrm{n.r.}}$ 在 $S$ 中开。若 $G$ 还在 $S$ 上\Term{locally-finite-presentation}{局部有限呈示}，则 $Z=G$，故 $V_{\mathrm{n.r.}}$ 是 $U_{\mathrm{n.r.}}$ 在 $\pi_G$ 下的\Term{inverse-image}{逆像}。这就证明了 (ii)。
\end{proof}
\begin{SGAresult}{推论}{2.6}
设 $u:G\to H$ 是一个\Term{radicial}{根性}的 $S$-\Term{group-morphism}{群态射}（若 $u$ 是\Term{monomorphism}{单态射}，则此条件成立，见 EGA I, 3.5.4）。假设 $G$（相应地 $H$）在 $S$ 上\Term{locally-finite-presentation}{局部有限呈示}且\Term{flat}{平坦}（相应地\Term{locally-finite-type}{局部有限型}）。则使 $u_s$ 为\Term{open-immersion}{开浸入}的 $s\in S$ 所成的\Term{set}{集合} $U$ 在 $S$ 中开，且 $u$ 在 $U$ 上方的\Term{restriction}{限制}是\Term{open-immersion}{开浸入}。
\end{SGAresult}
\begin{proof}
由 2.5 (i)，使 $u_s$ 为\Term{etale}{平展}的 $s\in S$ 所成的\Term{set}{集合} $U_0$ 在 $S$ 中开。$u$ \Term{radicial}{根性}，故每个 $u_s$ 也\Term{radicial}{根性}；由 EGA IV$_4$, 17.9.1，$U=U_0$，故 $U$ 开。最后，由 2.5 (i)，$u$ 在 $U$ 上方的\Term{restriction}{限制}\Term{etale}{平展}；$u$ 又\Term{radicial}{根性}，所以此\Term{restriction}{限制}是\Term{open-immersion}{开浸入}（EGA IV$_4$, 17.9.1）。
\end{proof}
\begin{SGAresult}{命题}{2.7}
设 $G$ 是 $S$-群。下列条件\Term{equivalence}{等价}：
\begin{enumerate}
\item[(i)] $G$ 在\Term{identity-section}{单位截面}的各点处在 $S$ 上\Term{unramified}{非分歧}；
\SourcePage{348}
\item[(ii)] \Term{identity-section}{单位截面}是\Term{open-immersion}{开浸入}；
\item[(iii)] $G$ 在\Term{identity-section}{单位截面}的各点处在 $S$ 上\Term{finite-presentation}{有限呈示}，且对每个 $s\in S$，$G_s$ 在 $\kappa(s)$ 上\Term{unramified}{非分歧}。
\end{enumerate}
若还假设 $G$ 在 $S$ 上\Term{locally-finite-presentation}{局部有限呈示}，则上述三个条件各自\Term{equivalence}{等价}于：
\begin{enumerate}\item[(iv)] $G$ 在 $S$ 上\Term{unramified}{非分歧}。\end{enumerate}
\end{SGAresult}
\begin{proof}
(i) 与 (ii) 的\Term{equivalence}{等价}性来自下面更一般的\Term{lemma}{引理} 2.7.1。注意，由 1.3.1.1，条件 (i) 或 (iii) 都\Term{pushout}{推出}：对每个 $s\in S$，$G_s$ 在 $\kappa(s)$ 上\Term{locally-finite-type}{局部有限型}。由 EGA IV$_4$, 17.4.1，断言 (i) \Term{equivalence}{等价}于对每个 $s\in S$，$G_s$ 在其\Term{identity-element}{单位点} $e_s$ 处在 $\kappa(s)$ 上\Term{unramified}{非分歧}；由 1.3.1，这又\Term{equivalence}{等价}于 $G_s$ 在 $\kappa(s)$ 上\Term{unramified}{非分歧}。所以 (i) 与 (iii) \Term{equivalence}{等价}。最后，若 $G$ 在 $S$ 上\Term{locally-finite-presentation}{局部有限呈示}，则由 EGA IV$_4$, 17.4.1，(iii) 与 (iv) \Term{equivalence}{等价}。
\end{proof}
\begin{SGAresult}{\Term{lemma}{引理}}{2.7.1}
设 $G$ 是带有\Term{section}{截面} $\varepsilon$ 的 $S$-\Term{scheme}{概形}。则 $G$ 在此\Term{section}{截面}的各点处在 $S$ 上\Term{unramified}{非分歧}，当且仅当 $\varepsilon$ 是\Term{open-immersion}{开浸入}。
\end{SGAresult}
\begin{proof}
必要性由 EGA IV$_4$, 17.4.1 a)$\Rightarrow$b$''$ 得到。反之，若 $\varepsilon$ 是\Term{open-immersion}{开浸入}，则\Term{structure-morphism}{结构态射} $G\to S$ 在 $\varepsilon(S)$ 上的\Term{restriction}{限制}是\Term{isomorphism}{同构}，故 $G$ 在 $\varepsilon(S)$ 的各点处在 $S$ 上\Term{unramified}{非分歧}。
\end{proof}
\begin{SGAresult}{推论}{2.8}
设 $u:G\to H$ 是 $S$-\Term{group-morphism}{群态射}。假设对每个 $s\in S$，$G_s$ 在 $\kappa(s)$ 上\Term{locally-finite-type}{局部有限型}。\footnote{\textbf{编者注（21）：}调整了表述，将“\Term{quasi-finite}{拟有限}”情形和“\Term{unramified}{非分歧}”情形的断言分开。}

 a) 若 $u$ \Term{locally-finite-type}{局部有限型}，则下列条件\Term{equivalence}{等价}：
\begin{enumerate}
\item[(i)] $u$ \Term{locally-quasi-finite}{局部拟有限}；
\item[(ii)] 对每个 $s\in S$，$u_s:G_s\to H_s$ \Term{locally-quasi-finite}{局部拟有限}；
\item[(iii)] $\operatorname{Ker}u$ 在 $S$ 上\Term{locally-quasi-finite}{局部拟有限}；
\item[(iv)] $\operatorname{Ker}u$ 的\Term{fibre}{纤维}都\Term{discrete}{离散}。
\end{enumerate}
 b) 若 $u$ \Term{locally-finite-presentation}{局部有限呈示}，则下列条件\Term{equivalence}{等价}：
\begin{enumerate}
\item[(v)] $u$ \Term{unramified}{非分歧}；
\item[(vi)] 对每个 $s\in S$，$u_s:G_s\to H_s$ \Term{unramified}{非分歧}；
\item[(vii)] $\operatorname{Ker}u$ \Term{unramified}{非分歧}；
\item[(viii)] \Term{identity-section}{单位截面} $S\to\operatorname{Ker}u$ 是\Term{open-immersion}{开浸入}。
\end{enumerate}
\end{SGAresult}
\begin{proof}
(i)$\Leftrightarrow$(ii) 和 (v)$\Leftrightarrow$(vi) 来自 2.5 (ii)；下面的回顾 2.8.1 说明 (i)$\Rightarrow$(iii) 和 (v)$\Rightarrow$(vii)。
\SourcePage{349}
对每个 $s\in S$，以 $e_s$ 表示 $H_s$ 的\Term{identity-element}{单位元}。则 (iii)（相应地 (vii)）\Term{pushout}{推出}，对每个 $s\in S$，
\[
(\operatorname{Ker}u)_s=\operatorname{Ker}u_s=u_s^{-1}(e_s)
\]
在 $\kappa(s)=\kappa(e_s)$ 上\Term{locally-quasi-finite}{局部拟有限}（相应地\Term{unramified}{非分歧}）；所以 $u_s$ 在 $G_s$ 的\Term{identity-element}{单位点}处\Term{quasi-finite}{拟有限}（相应地\Term{unramified}{非分歧}）；由 1.3，$u_s$ \Term{locally-quasi-finite}{局部拟有限}（相应地\Term{unramified}{非分歧}）。故 (iii)$\Rightarrow$(ii)，(vii)$\Rightarrow$(vi)。最后，(ii)$\Leftrightarrow$(iv) 由 1.4.1 得到，(vii)$\Leftrightarrow$(viii) 由 2.7 得到。
\end{proof}
\textbf{回顾 2.8.1。}回顾（参见第 I 讲 2.3.6.1），对给定的 $S$-\Term{group-morphism}{群态射} $u:G\to H$，$u$ 的\Term{kernel}{核}记作 $\operatorname{Ker}u$，是 $G$ 的如下\Term{group-subfunctor}{子群函子}：对每个\Term{morphism}{态射} $f:T\to S$，置
\[
(\operatorname{Ker}u)(T)=\{a\in G(T)\mid u\circ a=\varepsilon_H\circ f\}.
\]
由同处（或 EGA I, 4.4.1），此\Term{functor}{函子}由 $S$-群 $G\times_HS=u^{-1}(\varepsilon_H(S))$ \Term{representable}{表示}，后者也简记为 $\operatorname{Ker}u$。特别地，\Term{structure-morphism}{结构态射} $\operatorname{Ker}u\to S$ 是 $u$ 的\Term{base-change}{基变换}。
\begin{SGAresult}{\Term{lemma}{引理}}{2.9}
设 $\pi:X\to S$ 是有一个 $S$-\Term{section}{截面} $\varepsilon$ 的\Term{morphism}{态射}。
\begin{enumerate}
\item[(i)] 若 $\pi$ 为\Term{injectivity}{单射}，则 $\pi$ \Term{integral-morphism}{整}。\footnote{这也是 EGA IV$_4$, 18.12.11 的特殊情形，因为 $\pi$ 显然是\Term{universal-homeomorphism}{普遍同胚}。}
\item[(ii)] 若 $\pi$ \Term{locally-finite-type}{局部有限型}，且对每个 $s\in S$，$\pi_s$ 都是\Term{isomorphism}{同构}，则 $\pi$ 是\Term{isomorphism}{同构}。\footnote{这也是 EGA IV$_4$, 18.12.6 的直接推论。}
\end{enumerate}
\end{SGAresult}
\begin{proof}
首先，由下面的\Term{lemma}{引理} 2.9.1，$\pi$ 是\Term{homeomorphism}{同胚}，故为\Term{affine}{仿射}\Term{morphism}{态射}。

若 $\pi$ 为\Term{injectivity}{单射}，则 $\varepsilon$ 为\Term{surjective}{满射}。$\varepsilon$ 是\Term{surjective}{满}\Term{immersion}{浸入}，所以 $\varepsilon(S)$ 与 $X$ 中由 $\mathcal O_X$ 的\Term{nilideal}{幂零元理想} $I$ 定义的\Term{closed-subscheme}{闭子概形}\Term{isomorphism}{同构}。$\varepsilon$ 是 $\pi$ 的 $S$-\Term{section}{截面}，故有 $\mathcal O_S$-模的\Term{direct-sum}{直和}\Term{resolution}{分解} $\mathcal O_X=\mathcal O_S\oplus I$。$I$ 是 $\mathcal O_X$ 的\Term{nilideal}{幂零元理想}，显然在 $\mathcal O_S$ 上整，因此 $\mathcal O_X$ 在 $\mathcal O_S$ 上整，故 $\pi$ \Term{integral-morphism}{整}。

再假设 $\pi$ \Term{locally-finite-type}{局部有限型}。$\pi\circ\varepsilon=\operatorname{id}_S$，故 $\varepsilon$ \Term{locally-finite-presentation}{局部有限呈示}（参见 EGA IV$_1$, 1.4.3 (v)），因而 $I$ 是 $\mathcal O_X$ 的\Term{finitely-generated}{有限生成}\Term{ideal}{理想}（EGA IV$_1$, 1.4.1）。对每个 $s\in S$，有
\[
\mathcal O_{X_s}=\kappa(s)\oplus I\otimes_{\mathcal O_S}\kappa(s).
\]
由假设 $\varepsilon_s$ 是\Term{isomorphism}{同构}，故对所有 $s\in S$，$I\otimes_{\mathcal O_S}\kappa(s)=0$；从而对所有 $x\in X$，$I\otimes_{\mathcal O_X}\kappa(x)=0$。由 \Term{nakayama-lemma}{Nakayama 引理}，$I=0$，所以 $\pi$ 是\Term{isomorphism}{同构}。
\end{proof}
\begin{SGAresult}{\Term{lemma}{引理}}{2.9.1}
若\Term{scheme}{概形}\Term{morphism}{态射} $f:X\to Y$ 是\Term{homeomorphism}{同胚}，则 $f$ 是\Term{affine}{仿射}\Term{morphism}{态射}。\footnote{稍微更一般且证明相同的结果，参见 EGA IV$_4$, 18.12.7.1。}
\end{SGAresult}
\begin{proof}
只需证明每个 $y\in Y$ 都有一个\Term{open-neighborhood}{开邻域} $W$，使 $f$ 在 $W$ 上方的\Term{restriction}{限制}是\Term{affine}{仿射}\Term{morphism}{态射}。取 $y\in Y$，
\SourcePage{350}
令 $V$ 是 $y$ 在 $Y$ 中的一个\Term{affine}{仿射}\Term{open-neighborhood}{开邻域}。置 $V'=f^{-1}(V)$，则 $V'$ 是 $x=f^{-1}(y)$ 在 $X$ 中的\Term{open-neighborhood}{开邻域}。可取 $x$ 的\Term{affine}{仿射}\Term{open-neighborhood}{开邻域} $W'\subset V'$。置 $W=f(W')$。$W$ 是 $y$ 在 $Y$ 中的\Term{open-neighborhood}{开邻域}，且包含于\Term{affine}{仿射}\Term{scheme}{概形} $V$，所以 $W$ \Term{separated}{分离}。因 $W'$ 为\Term{affine}{仿射}\Term{scheme}{概形}，$f$ 在 $W$ 上方的\Term{restriction}{限制}就是\Term{affine}{仿射}\Term{morphism}{态射}（EGA II, 1.2.3）。
\end{proof}
\begin{SGAresult}{推论}{2.10}
设 $G$ 是\Term{locally-finite-type}{局部有限型} $S$-群。若对每个 $s\in S$，$G_s$ 都是单位 $\kappa(s)$-群，则 $G$ 是单位 $S$-群。
\end{SGAresult}
更一般地：
\begin{SGAresult}{推论}{2.11}
设 $f:X\to Y$ 是\Term{locally-finite-type}{局部有限型} $S$-\Term{morphism}{态射}。则 $f$ 是\Term{monomorphism}{单态射}，当且仅当对每个 $s\in S$，$f_s$ 是\Term{monomorphism}{单态射}。
\end{SGAresult}
\begin{proof}
必要性显然；证明充分性。若每个 $f_s$ 是\Term{monomorphism}{单态射}，则每个 $y\in Y$ 对应的 $f_y$ 更是\Term{monomorphism}{单态射}；故可假设 $Y=S$。

由 EGA I, 5.3.8，要证明 $f$ 是\Term{monomorphism}{单态射}，只需证明 $\Delta_f:X\to X\times_SX$ 是\Term{isomorphism}{同构}，\Term{equivalence}{等价}地，只需证明第一\Term{projection}{投影} $p:X\times_SX\to X$ 是\Term{isomorphism}{同构}。若 $f_s$ 是\Term{monomorphism}{单态射}，同样由 EGA I, 5.3.8，第一\Term{projection}{投影} $X_s\times_{\kappa(s)}X_s\to X_s$（即 $p_s$）是\Term{isomorphism}{同构}。而 $p$ 有 $S$-\Term{section}{截面} $\Delta_f$，所以\Term{lemma}{引理} 2.9 说明：若每个 $p_s$ 是\Term{isomorphism}{同构}，则 $p$ 亦然。
\end{proof}
\begin{SGAresult}{推论}{2.12}
设 $G$ 是有一个 $S$-\Term{section}{截面} $\varepsilon$ 的 $S$-\Term{scheme}{概形}，且\Term{structure-morphism}{结构态射} $\pi:G\to S$ 是\Term{closed-morphism}{闭态射}。设 $s\in S$，$\pi$ 在 $\varepsilon(s)$ 处\Term{finite-presentation}{有限呈示}，且 $\varepsilon_s:\operatorname{Spec}\kappa(s)\to G_s$ 是\Term{isomorphism}{同构}（\Term{equivalence}{等价}地，$\pi_s:G_s\to\operatorname{Spec}\kappa(s)$ 是\Term{isomorphism}{同构}）。则 $s$ 在 $S$ 中有一个\Term{open-neighborhood}{开邻域} $U$，使 $\varepsilon|_U:U\to G\times_SU$ 是\Term{isomorphism}{同构}。
\end{SGAresult}
\begin{proof}
令 $V$ 是 $G$ 中 $\pi$ 为\Term{unramified}{非分歧}的点集。已知 $V$ 是\Term{open-subset}{开集}（参见编者注（6）），且包含 $\varepsilon(s)$。故 $W=\varepsilon^{-1}(V)$ 是 $S$ 中包含 $s$ 的\Term{open-subset}{开集}，且对每个 $t\in W$，$\pi$ 在 $\varepsilon(t)$ 处\Term{unramified}{非分歧}。$\pi$ 是\Term{closed-morphism}{闭态射}，所以 $\pi|_W$ 亦然，故可假设 $W=S$。

由 2.7.1，$X=G-\varepsilon(S)$ 是 $G$ 的\Term{closed-subset}{闭子集}，且不与 $G_s$ 相交。$\pi$ 是\Term{closed-morphism}{闭态射}，所以 $F=\pi(X)$ 是 $S$ 的\Term{closed-subset}{闭子集}，且不包含 $s$。置 $U=S-F$，则 $U$ 是 $S$ 的\Term{open-subset}{开集}，$\varepsilon|_U$ 是 $U$ 到 $G|_U$ 的\Term{isomorphism}{同构}。
\end{proof}

\SGAsection{3}{局部有限表现群的单位分支}
\textbf{3.0。}设 $A$（相应地 $B$）是 $S$-\Term{scheme}{概形} $X$（相应地 $Y$）的\Term{subset}{子集}。滥用记号，以 $A\times_SB$ 表示 $X\times_SY$ 中第一\Term{projection}{投影}属于 $A$、第二\Term{projection}{投影}属于 $B$ 的点组成的\Term{subset}{子集}。
\SourcePage{351}
设 $A$ 是 $S$-群 $G$ 的\Term{subset}{子集}。若 $c(A)\subset A$ 且 $\mu(A\times_SA)\subset A$，则称 $A$ 对 $G$ 的\Term{group-law}{群运算}\Term{stable}{稳定}。
\begin{SGAresult}{定义}{3.1}
设 $G$ 是满足下列条件的 $S$-\Term{group-functor}{群函子}：
\[
(+)\qquad\text{对每个 }s\in S,\quad G_s=G\otimes_S\kappa(s)\text{ 都可表示。}
\]
以 $\underline G_s^0$ 表示 $\kappa(s)$-群 $G_s$ 中\Term{identity-element}{单位元}所在的\Term{connected-component}{连通分支}。\footnote{\textbf{编者注（22）：}这个记号略有滥用，但若这个\Term{connected-component}{连通分支}是 $G$ 的某个\Term{open-subgroup-scheme}{开子群概形} $G^0$ 的底\Term{topological-space}{拓扑空间}，它与第 VI$_A$ 讲 2.3 的记号相容，参见第 VI$_A$ 讲 2.2 bis。}定义 $G$ 的一个子 $S$-\Term{group-functor}{群函子} $G^0$，称为 $G$ 的\Term{identity-component}{单位分支}，其定义为：对每个\Term{morphism}{态射} $T\to S$，
\[
G^0(T)=\{u\in G(T)\mid\forall s\in S,\ u_s(T_s)\subset\underline G_s^0\}.
\]
这样定义了从 $(\widehat{\mathrm{Sch}/S})\text{-}\mathrm{gr}$ 到自身的\Term{functor}{函子} $G\mapsto G^0$。
\end{SGAresult}
\textbf{注记 3.2。}(i) 若 $S$-\Term{group-functor}{群函子} $G$ 满足 $(+)$，则由第 II 讲，
\[
\underline{\operatorname{Lie}}(G^0/S)=\underline{\operatorname{Lie}}(G/S).
\]
\footnote{\textbf{编者注（23）：}加入了下面的说明。}事实上，对每个 $S$-\Term{scheme}{概形} $T$，以 $I_T$ 表示 $T$ 上的\Term{dual-numbers}{对偶数}\Term{scheme}{概形}，以 $\tau:T\to I_T$ 表示\Term{zero-section}{零截面}（参见第 II 讲 2.1）。按定义，
\[
\underline{\operatorname{Lie}}(G/S)(T)=\{u\in G(I_T)\mid u\circ\tau=e\},
\]
其中 $e:T\to G$ 是 $T\to S$ 与\Term{identity-section}{单位截面} $S\to G$ 的复合；同样，
\begin{align*}
\underline{\operatorname{Lie}}(G^0/S)(T)
&=\{u\in G^0(I_T)\mid u\circ\tau=e\}\\
&=\{u\in G(I_T)\mid u\circ\tau=e\text{ 且 }u_s((I_T)_s)\subset\underline G_s^0,\ \forall s\in S\}.
\end{align*}
对每个 $s\in S$，$T_s$ 与 $(I_T)_s=I_{T_s}$ 有相同的\Term{underlying-set}{底集}合，所以若 $u\in\underline{\operatorname{Lie}}(G/S)(T)$，则 $u_s(I_{T_s})=e_s$，其中 $e_s$ 是 $G_s$ 的\Term{identity-element}{单位点}。故 $u\in\underline{\operatorname{Lie}}(G^0/S)(T)$。因此对每个 $T$，包含关系 $\underline{\operatorname{Lie}}(G^0/S)(T)\subset\underline{\operatorname{Lie}}(G/S)(T)$ 都是等式，从而得到所述等式。

(ii) 设 $G,G'$ 是两个满足 $(+)$ 的 $S$-\Term{group-functor}{群函子}，则：
\begin{enumerate}
\item[a)] 若 $G\subset G'$，则 $G^0\subset G'^0$；
\item[b)] 若 $G\subset G'$ 且 $G'^0\subset G$，则 $G^0=G'^0$；
\item[c)] 若对每个 $s\in S$，$G_s$ 在 $\kappa(s)$ 上\Term{locally-finite-type}{局部有限型}，则由第 VI$_A$ 讲 2.3，$G^0$ 满足 $(+)$，且 $(G^0)^0=G^0$。
\end{enumerate}
\begin{SGAresult}{命题}{3.3}
设 $G$ 是满足 $(+)$ 的 $S$-\Term{group-functor}{群函子}，$S'$ 是 $S$-\Term{scheme}{概形}。则 $(G\times_SS')^0=G^0\times_SS'$；换言之，\Term{functor}{函子} $G\mapsto G^0$
\SourcePage{352}
与\Term{base-change}{基变换}交换，即下图交换：
\[
\begin{tikzcd}[column sep=large]
(\widehat{\mathrm{Sch}/S})\text{-}\mathrm{gr}\arrow[r,"{G\mapsto G^0}"]\arrow[d]
 &(\widehat{\mathrm{Sch}/S})\text{-}\mathrm{gr}\arrow[d]\\
(\widehat{\mathrm{Sch}/S'})\text{-}\mathrm{gr}\arrow[r]
 &(\widehat{\mathrm{Sch}/S'})\text{-}\mathrm{gr}.
\end{tikzcd}
\]
\end{SGAresult}
\begin{proof}
只需对每个 $s'\in S'$ 检验下式，其中 $s$ 是 $s'$ 在 $S$ 中的像：
\[
\bigl((G\times_SS')\otimes_{S'}\kappa(s')\bigr)^0
=(G_s\otimes_{\kappa(s)}\kappa(s'))^0
=\underline G_s^0\otimes_{\kappa(s)}\kappa(s').
\]
这由第 VI$_A$ 讲 2.1.2 得到。为后面 4.2 的使用，注意这里并未用到 $G_s$ 的\Term{group-structure}{群结构}，只用了 $\underline G_s^0$ 有一个\Term{rational-point}{有理点} $\varepsilon(s)$，因而\Term{geometrically-connected}{几何连通}这一事实（亦见 EGA IV$_2$, 4.5.14）。
\end{proof}
\textbf{特殊情形 3.4。}设 $G$ 是 $S$-\Term{group-scheme}{群概形}，以 $\underline G^0$ 表示当 $s$ 遍历 $S$ 时各 $\underline G_s^0$ 的并。则 $\underline G^0$ 是 $G$ 中对\Term{group-law}{群运算}\Term{stable}{稳定}的\Term{subset}{子集}（参见 3.0），且对任意\Term{morphism}{态射} $S'\to S$，
\[
G^0(S')=\{u\in G(S')\mid u(S')\subset\underline G^0\}.
\]
若 $\underline G^0$ 是 $G$ 的\Term{open-subset}{开子集}，则 $G^0$ 显然由 $G$ 在\Term{open-subset}{开集} $\underline G^0$ 上诱导的\Term{subscheme}{子概形}\Term{representable}{表示}。
\begin{SGAresult}{命题}{3.5}
设 $S$ 是\Term{quasi-compact}{拟紧}\Term{quasi-separated}{拟分离}\Term{scheme}{概形}，$G$ 是\Term{fibre}{纤维}\Term{locally-finite-type}{局部有限型}的 $S$-群。则 $G$ 有一个包含 $\underline G^0$ 的\Term{quasi-compact}{拟紧}\Term{open-subset}{开集} $U$。
\end{SGAresult}
\begin{proof}
\Term{identity-section}{单位截面} $\varepsilon$ 是\Term{immersion}{浸入}，故 $\varepsilon(S)$ 是 $G$ 的\Term{quasi-compact}{拟紧}子\Term{space}{空间}；所以 $G$ 有一个包含 $\varepsilon(S)$ 的\Term{quasi-compact}{拟紧}\Term{open-subset}{开集} $V$。$S$ \Term{quasi-separated}{拟分离}且 $V$ \Term{quasi-compact}{拟紧}，所以 $V$ 在 $S$ 上\Term{quasi-compact}{拟紧}（EGA IV$_1$, 1.2.4），故 $V\times_SV$ 在 $S$ 上\Term{quasi-compact}{拟紧}，从而自身\Term{quasi-compact}{拟紧}。于是 $V\cdot V=\mu(V\times_SV)$ \Term{quasi-compact}{拟紧}。置 $V_s=V\cap G_s$、$V_s^0=V\cap\underline G_s^0$。$V_s^0$ 是 $\underline G_s^0$ 的\Term{open-subset}{开集}，而 $\underline G_s^0$ \Term{irreducible}{不可约}（第 VI$_A$ 讲 2.4），故 $V_s^0$ 在其中\Term{dense}{稠密}，因而 $V_s^0\cdot V_s^0=\underline G_s^0$（第 VI$_A$ 讲 0.5）。所以 $V_s\cdot V_s\supset\underline G_s^0$，从而 $V\cdot V\supset\underline G^0$。最后，$V\cdot V$ \Term{quasi-compact}{拟紧}，所以 $G$ 有一个包含 $V\cdot V$，因而也包含 $\underline G^0$ 的\Term{quasi-compact}{拟紧}\Term{open-subset}{开集} $U$。
\end{proof}
\begin{SGAresult}{推论}{3.6}
若 $S$-群 $G$ 的\Term{fibre}{纤维}都\Term{locally-finite-type}{局部有限型}且\Term{connected}{连通}，则 $G$ 在 $S$ 上\Term{quasi-compact}{拟紧}。
\end{SGAresult}
\begin{proof}
此断言在 $S$ 上是\Term{local-scheme}{局部}的（EGA I, 6.6.1），故可假设 $S$ \Term{affine}{仿射}。由 3.5，$G$ 有一个\Term{quasi-compact}{拟紧}\Term{open-subset}{开集} $U$，包含 $\underline G^0=G$；故 $G$ \Term{quasi-compact}{拟紧}，从而在\Term{affine}{仿射}\Term{scheme}{概形} $S$ 上\Term{quasi-compact}{拟紧}（EGA I, 6.6.4 (v)）。
\end{proof}
\SourcePage{353}
\begin{SGAresult}{命题}{3.7}
\footnote{\textbf{编者注（24）：}回顾下列定义与结果（参见 EGA 0$_{\mathrm{III}}$, \S9.1，EGA IV$_1$, \S\S1.8、1.9）。设 $X$ 是\Term{topological-space}{拓扑空间}。(i) 若包含映射 $U\hookrightarrow X$ \Term{quasi-compact}{拟紧}，即对每个\Term{quasi-compact}{拟紧}\Term{open-subset}{开集} $V\subset X$，$U\cap V$ 都\Term{quasi-compact}{拟紧}，则称\Term{open-subset}{开集} $U$ \Term{retrocompact}{回紧}。(ii) 若\Term{subset}{子集} $C$ 是\Term{finite}{有限}个形如 $U-V$ 的\Term{set}{集合}之并，其中 $U,V$ 是 $X$ 的\Term{retrocompact}{回紧}\Term{open-subset}{开集}，则称 $C$ \Term{constructible}{可构造}。(iii) 若对每个 $x\in X$，有 $x$ 的\Term{open-neighborhood}{开邻域} $U$ 使 $L\cap U$ 在 $U$ 中\Term{constructible}{可构造}，则称\Term{subset}{子集} $L$ \Term{locally-constructible}{局部可构造}。（注意，若 $X$ \Term{quasi-compact}{拟紧}且\Term{quasi-separated}{拟分离}，则此时 $L$ \Term{constructible}{可构造}。）(iv) 若对每个 $x\in X$，有 $x$ 的\Term{open-neighborhood}{开邻域} $U$ 使 $I\cap U$ 是 $U$ 中若干\Term{locally-constructible}{局部可构造}\Term{subset}{子集}之并，则称\Term{subset}{子集} $I$ \Term{ind-constructible}{归纳可构造}。现设 $f:X\to Y$ 是\Term{scheme}{概形}\Term{morphism}{态射}。由 Chevalley \Term{constructibility-theorem}{可构造性定理}（参见 EGA IV$_1$, 1.8.4、1.9.5 (viii)），若 $f$ \Term{finite-presentation}{有限呈示}（相应地\Term{locally-finite-presentation}{局部有限呈示}），则每个\Term{locally-constructible}{局部可构造}\Term{subset}{子集}在 $f$ 下的像都是\Term{locally-constructible}{局部可构造}的（相应地\Term{ind-constructible}{归纳可构造}的）。}设 $G$ 是\Term{locally-finite-presentation}{局部有限呈示} $S$-群。则：
\begin{enumerate}
\item[(i)] $\underline G^0$ 在 $G$ 中\Term{ind-constructible}{归纳可构造}。
\item[(ii)] 若 $G$ 还在 $S$ 上\Term{quasi-separated}{拟分离}，且 $S$ \Term{quasi-compact}{拟紧}\Term{quasi-separated}{拟分离}，则 $\underline G^0$ \Term{constructible}{可构造}。
\item[(iii)] 因而，若 $G$ 在 $S$ 上\Term{quasi-separated}{拟分离}，则 $\underline G^0$ \Term{locally-constructible}{局部可构造}。
\end{enumerate}
\end{SGAresult}
\begin{proof}
先证第一个断言。$\pi:G\to S$ \Term{locally-finite-presentation}{局部有限呈示}，故对给定的 $s\in S$，可取 $G$ 中包含 $\varepsilon(s)$ 的\Term{open-subset}{开集} $U$，以及 $S$ 中包含 $s$ 的\Term{open-subset}{开集} $V$，使 $\pi(U)\subset V$，且由 $\pi$ 诱导的\Term{morphism}{态射} $\pi':U\to V$ \Term{finite-presentation}{有限呈示}。则 $T=\varepsilon^{-1}(U)$ 是 $S$ 的\Term{open-subset}{开集}；置 $W=\pi'^{-1}(T)$、$\pi''=\pi'|_W$，于是 $\pi'':W\to T$ \Term{finite-presentation}{有限呈示}，且由 $\varepsilon$ 诱导的\Term{morphism}{态射} $\varepsilon'':T\to W$ 是其\Term{section}{截面}。

对每个 $t\in T$，$\underline G_t^0$ \Term{irreducible}{不可约}（第 VI$_A$ 讲 2.4），所以 $W\cap\underline G_t^0$ 在 $\underline G_t^0$ 中\Term{dense}{稠密}，故\Term{irreducible}{不可约}，从而\Term{connected}{连通}。因此，它正是 $\pi''^{-1}(t)$ 中包含 $\varepsilon''(t)$ 的\Term{connected-component}{连通分支}。由 EGA IV$_3$, 9.7.12，当 $t\in T$ 时各 $W\cap\underline G_t^0$ 的并 $W^0$ 在 $W$ 中\Term{locally-constructible}{局部可构造}。另一方面，由第 VI$_A$ 讲 0.5，
\[
\underline G^0\cap\pi^{-1}(T)=W^0\cdot W^0,
\]
即 $\underline G^0\cap\pi^{-1}(T)$ 是 $W^0\times_TW^0$ 在由 $\mu$ 诱导的\Term{morphism}{态射} $\mu'':W\times_TW\to\pi^{-1}(T)$ 下的像。\footnote{\textbf{编者注（25）：}下面将原文展开说明。}$W\times_TW$（相应地 $\pi^{-1}(T)$）在 $T$ 上\Term{finite-presentation}{有限呈示}（相应地\Term{locally-finite-presentation}{局部有限呈示}），所以由 EGA IV$_1$, 1.4.3、1.2.2，$\mu''$ \Term{locally-finite-presentation}{局部有限呈示}且\Term{quasi-separated}{拟分离}；若 $\pi^{-1}(T)$ 还在 $T$ 上\Term{quasi-separated}{拟分离}，则 $\mu''$ \Term{quasi-compact}{拟紧}（同处 1.2.4），故\Term{finite-presentation}{有限呈示}。$W^0$ 在 $W$ 中\Term{locally-constructible}{局部可构造}，因而 $W^0\times_TW^0$ 在 $W\times_TW$ 中\Term{locally-constructible}{局部可构造}；由 Chevalley \Term{constructibility-theorem}{可构造性定理}（同处 1.8.4、1.9.5 (viii)），$\underline G^0\cap\pi^{-1}(T)$ 在 $\pi^{-1}(T)$ 中\Term{ind-constructible}{归纳可构造}；若 $G$（从而 $\pi^{-1}(T)$）在 $T$ 上\Term{quasi-separated}{拟分离}，它还在 $\pi^{-1}(T)$ 中\Term{locally-constructible}{局部可构造}。这证明了 (i)、(iii)。

再假设 $G$ 在 $S$ 上\Term{quasi-separated}{拟分离}，且 $S$ \Term{quasi-compact}{拟紧}\Term{quasi-separated}{拟分离}。由 3.5，$G$ 有一个包含 $\underline G^0$ 的\Term{quasi-compact}{拟紧}\Term{open-subset}{开集} $U$。$G$ 在 $S$ 上\Term{quasi-separated}{拟分离}，所以 $G$ 自身\Term{quasi-separated}{拟分离}；故\Term{open-subset}{开集} $U$ \Term{retrocompact}{回紧}（EGA IV$_1$, 1.2.7），只需证明 $\underline G^0$ 在 $U$ 中\Term{constructible}{可构造}（EGA 0$_{\mathrm{III}}$, 9.1.8）。
\SourcePage{354}
此外，$U$ \Term{quasi-compact}{拟紧}，所以在 $S$ 上\Term{quasi-compact}{拟紧}（EGA IV$_1$, 1.2.4）；它还在 $S$ 上\Term{quasi-separated}{拟分离}，因此 $\pi|_U$ \Term{finite-presentation}{有限呈示}。由 EGA IV$_3$, 9.7.12，$\underline G^0$ 在 $U$ 中\Term{locally-constructible}{局部可构造}；$U$ \Term{quasi-compact}{拟紧}\Term{quasi-separated}{拟分离}，故 $\underline G^0$ 在 $U$ 中\Term{constructible}{可构造}（EGA IV$_1$, 1.8.1）。这证明了 (ii)，并说明对 $S$ 的每个\Term{quasi-compact}{拟紧}\Term{quasi-separated}{拟分离}\Term{open-subset}{开集} $T$（例如每个\Term{affine}{仿射}\Term{open-subset}{开集}），$\underline G^0\cap\pi^{-1}(T)$ 都\Term{constructible}{可构造}。
\end{proof}
\begin{SGAresult}{推论}{3.8}
设 $S_0$ 是\Term{quasi-compact}{拟紧}\Term{quasi-separated}{拟分离}\Term{scheme}{概形}，$I$ 是\Term{filtered}{滤过递增}\Term{preordered-set}{预序集}，$(\mathcal A_i)_{i\in I}$ 是交换\Term{quasi-coherent}{拟凝聚} $\mathcal O_{S_0}$-\Term{algebra}{代数}的\Term{inductive-system}{归纳系}。置
\[
\mathcal A=\varinjlim\mathcal A_i,\quad S_i=\operatorname{Spec}\mathcal A_i\ (i\in I),\quad S=\operatorname{Spec}\mathcal A
\]
（参见 EGA II, 1.3.1）。设 $G$ 是\Term{locally-finite-presentation}{局部有限呈示} $S_0$-\Term{group-scheme}{群概形}。则\Term{canonical}{典范}映射 $\varinjlim G^0(S_i)\to G^0(S)$ 是\Term{bijection}{双射}。
\end{SGAresult}
\begin{proof}
$G$ 在 $S$ 上\Term{locally-finite-presentation}{局部有限呈示}，故由 EGA IV$_2$, 8.14.2 c)，\Term{canonical}{典范}映射 $\varinjlim G(S_i)\to G(S)$ 是\Term{bijection}{双射}。这立即说明 $\varinjlim G^0(S_i)\to G^0(S)$ 为\Term{injectivity}{单射}。再证\Term{surjective}{满射}。取 $g\in G^0(S)\subset G(S)$，则存在 $i\in I$，使 $g$ 通过 $S\to S_i$ 及某个 $g_i\in G(S_i)$ \Term{resolution}{分解}；由假设 $g^{-1}(\underline G^0)=S$。由 3.7，$\underline G^0$ 在 $G$ 中\Term{ind-constructible}{归纳可构造}，故 $g_i^{-1}(\underline G^0)$ 在 $S_i$ 中亦然。由 EGA IV$_2$, 8.3.4，存在 $j\geq i$，使 $g_j^{-1}(\underline G^0)=S_j$，其中 $g_j$ 是 $g_i$ 经\Term{base-change}{基变换} $S_j\to S_i$ 所得的映射。因此 $g_j\in G^0(S_j)$，所以 $g$ 来自 $\varinjlim G^0(S_i)$ 中的一个元素。
\end{proof}
\begin{SGAresult}{命题}{3.9}
设 $G$ 是\Term{locally-finite-presentation}{局部有限呈示} $S$-群。若 $G^0$ \Term{representable}{可表示}，则\Term{canonical}{典范}\Term{morphism}{态射} $i:G^0\to G$ 是\Term{open-immersion}{开浸入}；此外，$G^0$ 在 $S$ 上\Term{quasi-compact}{拟紧}。
\end{SGAresult}
\begin{proof}
$G^0$ 是 $G$ 的\Term{subfunctor}{子函子}，所以 $i$ 是\Term{monomorphism}{单态射}，因而\Term{radicial}{根性}。由\Term{functor}{函子} $G^0$ 的定义，直接按 EGA IV$_4$, 17.1.1 的定义检验可知 $i$ \Term{formally-etale}{形式平展}（注意 $\underline G^0$ 是 $i$ 在 $G$ 中的像）。最后，由 $S$-\Term{scheme}{概形}\Term{locally-finite-presentation}{局部有限呈示}的\Term{functor}{函子}刻画（EGA IV$_3$, 8.14.2 c)）以及 3.8，$G$ 在 $S$ 上\Term{locally-finite-presentation}{局部有限呈示}，所以 $G^0$ 亦然。因此 $i$ \Term{locally-finite-presentation}{局部有限呈示}（参见 EGA IV$_1$, 1.4.3），从而\Term{radicial}{根性}且\Term{etale}{平展}，故为\Term{open-immersion}{开浸入}（EGA IV$_4$, 17.9.1）。

最后一个断言由 3.6 得到。
\end{proof}
\begin{SGAresult}{\Term{theorem}{定理}}{3.10}
设 $G$ 是 $S$-群。下列条件\Term{equivalence}{等价}：
\begin{enumerate}
\item[(i)] $G$ 在\Term{identity-section}{单位截面}的各点处在 $S$ 上\Term{smooth}{光滑}。
\item[(ii)] $G$ 在\Term{identity-section}{单位截面}的各点处在 $S$ 上\Term{flat}{平坦}且\Term{locally-finite-presentation}{局部有限呈示}，并且对每个 $s\in S$，$G_s$ 在 $\kappa(s)$ 上\Term{smooth}{光滑}。
\item[(iii)] $G$ 有一个在 $S$ 上\Term{smooth}{光滑}的\Term{open-subgroup-scheme}{开子群概形} $G'$。
\item[(iv)] $G^0$ 由 $G$ 的一个在 $S$ 上\Term{smooth}{光滑}的\Term{open-subscheme}{开子概形}\Term{representable}{表示}。
\end{enumerate}
\end{SGAresult}
\SourcePage{355}
\begin{proof}
显然 (iv)$\Rightarrow$(iii)$\Rightarrow$(i)；由 1.3.1、2.4，(i) \Term{pushout}{推出} (ii)、(iii)。此外，由 EGA IV$_4$, 17.5.1，(ii)$\Rightarrow$(i)。

最后证明 (iii) \Term{pushout}{推出} (iv)。下面的\Term{lemma}{引理} 3.10.1 说明 $G'$ 包含 $\underline G^0$，并且 $G'^0=G^0$。所以只需证明 $\underline G^0$ 在 $G$ 中开；因为由 3.4，$G^0$ 将由 $G'$ 在\Term{open-subset}{开集} $\underline G'^0=\underline G^0$ 上诱导的\Term{smooth}{光滑}\Term{group-subscheme}{子群概形}\Term{representable}{表示}。因而可假设 $G'=G$。

要证 $\underline G^0$ 开，只需证明每个 $s\in S$ 都有一个\Term{neighborhood}{邻域} $T\subset S$，使 $\underline G^0\cap\pi^{-1}(T)$ 在 $\pi^{-1}(T)$ 中开。取 $s\in S$。$G=G'$，所以 $\pi$ \Term{locally-finite-presentation}{局部有限呈示}；仿照 3.7 的证明，\Term{constructible}{可构造} $S$ 中包含 $s$ 的\Term{open-subset}{开集} $T$、$G$ 中包含 $\varepsilon(s)$ 的\Term{open-subset}{开集} $W$，使由 $\pi$ 诱导的\Term{morphism}{态射} $\pi'':W\to T$ \Term{finite-presentation}{有限呈示}，并以由 $\varepsilon$ 诱导的\Term{morphism}{态射} $\varepsilon'':T\to W$ 为\Term{section}{截面}。对每个 $t\in T$，$W\cap\underline G_t^0$ 正是 $\pi''^{-1}(t)$ 中包含 $\varepsilon''(t)$ 的\Term{connected-component}{连通分支}。$\pi$ \Term{smooth}{光滑}，故 $\pi''$ \Term{smooth}{光滑}且\Term{finite-presentation}{有限呈示}。由 EGA IV$_3$, 15.6.5，$t\in T$ 时各 $W\cap\underline G_t^0$ 的并 $W^0$ 在 $W$ 中开。

另一方面，由第 VI$_A$ 讲 0.5，$W^0\cdot W^0=\underline G^0\cap\pi^{-1}(T)$，须证它在 $\pi^{-1}(T)$ 中开。故从现在起可假设 $T=S$，只需证明 $W^0\cdot W^0$ 在 $G$ 中开。$\pi$ \Term{universally-open}{普遍开}，故 $\mu$ 也\Term{universally-open}{普遍开}（第 VI$_A$ 讲 0.1）。所以 $W^0$ 在 $G$ 中开，便\Term{pushout}{推出} $W^0\cdot W^0=\mu(W^0\times_SW^0)$ 在 $G$ 中开。

4.4 将改进这一结果。
\end{proof}
\begin{SGAresult}{\Term{lemma}{引理}}{3.10.1}
设 $S$-群 $G$ 的\Term{fibre}{纤维}\Term{locally-finite-type}{局部有限型}。则 $G$ 的每个\Term{open-subgroup-scheme}{开子群概形} $H$ 都包含 $\underline G^0$，并满足 $G^0=H^0$。
\end{SGAresult}
\begin{proof}
取 $s\in S$，置 $G_s'^0=H_s\cap\underline G_s^0$。它是 $\underline G_s^0$ 的\Term{open-subset}{开集}；$\underline G_s^0$ \Term{irreducible}{不可约}（第 VI$_A$ 讲 2.4），所以它在 $\underline G_s^0$ 中\Term{dense}{稠密}，故 $G_s'^0\cdot G_s'^0=\underline G_s^0$（第 VI$_A$ 讲 0.5）。因此
\[
\underline G_s^0=G_s'^0\cdot G_s'^0\subset H_s\cdot H_s=H_s.
\]
故对每个 $s$，$\underline G_s^0=\underline H_s^0$，从而 $\underline G^0\subset H$，并且 $G^0=H^0$。
\end{proof}
\begin{SGAresult}{命题}{3.11}
设 $u:G\to H$ 是\Term{locally-finite-presentation}{局部有限呈示} $S$-群之间的\Term{group-morphism}{群态射}。若 $u$ \Term{flat}{平坦}，则由 $u$ 诱导的映射 $u^0:\underline G^0\to\underline H^0$ 为\Term{surjective}{满射}。若 $G$ 在 $S$ 上\Term{flat}{平坦}且 $H$ 的\Term{fibre}{纤维}\Term{reduced}{约化}，则逆命题也成立。\footnote{\textbf{编者注（26）：}与第 VI$_A$ 讲 5.6 比较。}
\end{SGAresult}
\begin{proof}
设 $u$ \Term{flat}{平坦}。则每个 $u_s$ 都\Term{flat}{平坦}且\Term{locally-finite-presentation}{局部有限呈示}，故为\Term{open-morphism}{开态射}（EGA IV$_2$, 2.4.6）；由 1.3.2，$u_s^0:\underline G_s^0\to\underline H_s^0$ 是\Term{surjective}{满射}。\footnote{\textbf{编者注（27）：}此处引用 1.3.2，缩短了原文；又在下面删除了对\Term{generic-flatness}{泛平坦性}\Term{theorem}{定理}（EGA IV$_2$, 6.9.1）的非必要引用，从而简化了原文。}故映射 $u^0:\underline G^0\to\underline H^0$ 为\Term{surjective}{满射}。

反之，假设此映射为\Term{surjective}{满射}，$G$ 在 $S$ 上\Term{flat}{平坦}，$H$ 的\Term{fibre}{纤维}\Term{reduced}{约化}。则对每个 $s\in S$，$u_s^0:\underline G_s^0\to\underline H_s^0$ 为\Term{surjective}{满射}，从而在 $\underline H_s^0$ 的\Term{generic-point}{泛点} $\eta$ 上方每一点都\Term{flat}{平坦}（因为 $\mathcal O_{H_s^0,\eta}$ 是
\SourcePage{356}
域）。由 1.3，$u_s$ \Term{flat}{平坦}；再由\Term{fibrewise-flatness-criterion}{逐纤维平坦性判据}（EGA IV$_3$, 11.3.11），$u$ \Term{flat}{平坦}。
\end{proof}

\SGAsection{4}{局部有限表现群的纤维维数}
\begin{SGAresult}{命题}{4.1}
设 $G$ 是\Term{locally-finite-type}{局部有限型} $S$-\Term{scheme}{概形}，带有 $S$-\Term{section}{截面} $\varepsilon$，并且对每个 $s\in S$，$\dim G_s=\dim_{\varepsilon(s)}G_s$（若 $G$ 为 $S$-群，则由 1.5 此条件成立）。
\begin{enumerate}
\item[(i)] 函数 $s\mapsto\dim G_s$ 在 $S$ 上\Term{upper-semicontinuous}{上半连续}。
\item[(ii)] 若 $G$ 还在 $S$ 上\Term{locally-finite-presentation}{局部有限呈示}，则此函数\Term{locally-constructible}{局部可构造}。
\end{enumerate}
\end{SGAresult}
\begin{proof}
令 $\pi:G\to S$ 为\Term{structure-morphism}{结构态射}。Chevalley \Term{semicontinuity}{半连续性}\Term{theorem}{定理}（EGA IV$_3$, 13.1.3）断言，函数 $x\mapsto\dim_x\pi^{-1}(\pi(x))$ 在 $G$ 上\Term{upper-semicontinuous}{上半连续}。而对每个 $s\in S$，
\[
\dim G_s=\dim\pi^{-1}(s)=\dim_{\varepsilon(s)}\pi^{-1}(\pi(\varepsilon(s))).
\]
函数 $s\mapsto\varepsilon(s)$ 在 $S$ 上\Term{continuous}{连续}，故复合函数 $s\mapsto\dim G_s$ 在 $S$ 上\Term{upper-semicontinuous}{上半连续}。

若 $G$ 在 $S$ 上\Term{locally-finite-presentation}{局部有限呈示}，则如上推理，要证 $s\mapsto\dim G_s$ \Term{locally-constructible}{局部可构造}，只需证明 $x\mapsto\dim_x\pi^{-1}(\pi(x))$ 在 $G$ 上\Term{locally-constructible}{局部可构造}；这由 EGA IV$_3$, 9.9.1 得到。
\end{proof}
\begin{SGAresult}{命题}{4.2}
设 $\pi:G\to S$ 是\Term{locally-finite-presentation}{局部有限呈示} $S$-\Term{scheme}{概形}，带有 $S$-\Term{section}{截面} $\varepsilon$，且满足下列两个条件（若 $G$ 是 $S$-群，则由 1.5 和第 VI$_A$ 讲 2.4.1，这些条件成立）：
\begin{enumerate}
\item[a)] 对每个 $s\in S$ 及每个 $x\in G_s$，$\dim G_s=\dim_xG_s$（\Term{equivalence}{等价}地，由 1.5，对每个 $s\in S$，$G_s$ 的各\Term{irreducible-component}{不可约分支}\Term{dimension}{维数}相同）。
\item[b)] 对每个 $s\in S$，若以 $\underline G_s^0$ 表示 $G_s$ 中包含 $\varepsilon(s)$ 的\Term{connected-component}{连通分支}，则 $\underline G_s^0$ \Term{geometrically-irreducible}{几何不可约}。
\end{enumerate}
取 $s\in S$。下列条件\Term{equivalence}{等价}：
\begin{enumerate}
\item[(i)] $G$ 在 $\underline G_s^0$ 的各点处在 $S$ 上\Term{universally-open}{普遍开}。
\item[(i bis)] $G$ 在 $\underline G_s^0$ 中 $\varepsilon(s)$ 的某个\Term{neighborhood}{邻域}的每一点处都在 $S$ 上\Term{universally-open}{普遍开}。\footnote{\textbf{编者注（28）：}原文写作“在 $\varepsilon(s)$ 的点处”；仅假设 $G$ 在 $\varepsilon(s)$ 处在 $S$ 上\Term{universally-open}{普遍开}是不够的，编者注（19）的例子说明了这一点。已相应修正证明。}
\item[(ii)] 函数 $t\mapsto\dim G_t$ 在 $s$ 的某个\Term{neighborhood}{邻域}中为常数。
\item[(iii)] $\underline G^0$ “在 $\underline G_s^0$ 的各点处在 $S$ 上\Term{universally-open}{普遍开}”：即给定 $S'\to S$、$s'\in S'_s$、$g\in\underline G_{s'}^0$ 及 $g$ 在 $G'=G_{S'}$ 中的\Term{open-neighborhood}{开邻域} $V$，则 $\pi(V\cap\underline G'^0)$ 是 $s'$ 在 $S'$ 中的\Term{open-neighborhood}{开邻域}。\footnote{\textbf{编者注（29）：}加入了 O. Gabber 指出的条件 (iii)；后面（5.7）将用到它。}
\end{enumerate}
\end{SGAresult}
\SourcePage{357}
\begin{proof}
显然 (i)$\Rightarrow$(i bis)。由 EGA IV$_3$, 14.3.3.1 (ii)，$\underline G_s^0$ 中 $\pi_G$ 为\Term{universally-open}{普遍开}的点集在 $\underline G_s^0$ 中闭。$\underline G_s^0$ \Term{irreducible}{不可约}，故 (i bis)$\Rightarrow$(i)。

证明 (i) \Term{pushout}{推出} (ii)。令 $T$ 是使 $\dim G_t=\dim G_s$ 的 $t\in S$ 组成的\Term{set}{集合}。由 4.1，$T$ \Term{locally-constructible}{局部可构造}，故由 EGA IV$_1$, 1.10.1，要证 $T$ 是 $s$ 的\Term{neighborhood}{邻域}，只需证明 $s$ 的每个\Term{generization}{一般化} $s'$ 都属于 $T$。

令 $x$ 是 $\underline G_s^0$ 的\Term{generic-point}{泛点}，$U$ 是 $G$ 中包含 $x$ 的\Term{open-subset}{开集}。$\pi_G$ 在 $\xi$ 处\Term{universally-open}{普遍开}，故由 EGA IV$_3$, 14.3.13，对 $s$ 的每个\Term{generization}{一般化} $s'$，有
\[
\dim(U\cap G_{s'})\geq\dim_x(U\cap G_s).
\]
结合假设 a)，得 $\dim G_{s'}\geq\dim G_s$。而由 4.1，$s\mapsto\dim G_s$ \Term{upper-semicontinuous}{上半连续}，故又有 $\dim G_{s'}\leq\dim G_s$，所以 $s'\in T$。这证明了 (i)$\Rightarrow$(ii)。

显然 (iii)$\Rightarrow$(i)；再证 (ii)$\Rightarrow$(iii)。\Term{fibre}{纤维}的\Term{dimension}{维数}在\Term{group-extension}{扩张}\Term{base-field}{基域}时\Term{invariant}{不变}，而 $\underline G^0$ 的形成与\Term{base-change}{基变换}交换（参见 3.3 的证明），故可假设 $S'=S$、$s'=s$。又因为 $G$ 中任何与 $\underline G_s^0$ 相交的\Term{open-subset}{开集} $V$ 都包含 $\underline G_s^0$ 的\Term{generic-point}{泛点} $\eta$，故可假设 $g=\eta$。

还可假设 $S$ \Term{affine}{仿射}。取 $\varepsilon(s)$ 的一个\Term{affine}{仿射}\Term{open-neighborhood}{开邻域} $U$，则它在 $S$ 上\Term{finite-presentation}{有限呈示}。将 $S$ 换成 $\varepsilon^{-1}(U)$，再将 $U$ 换成 $U\cap\pi^{-1}(S)$，就归结为 $\pi:U\to S$ \Term{finite-presentation}{有限呈示}、且有\Term{section}{截面} $\varepsilon:S\to U$ 的情形。由 EGA IV$_3$, 9.7.12，$\underline G^0\cap U$ 在 $U$ 中\Term{constructible}{可构造}。再将 $V$ 换成 $V\cap U$ 中的一个\Term{affine}{仿射}\Term{open-subset}{开集}，则 $\underline G^0\cap V$ 在 $V$ 中\Term{constructible}{可构造}。$\pi:V\to S$ \Term{finite-presentation}{有限呈示}，故由 Chevalley \Term{constructibility-theorem}{可构造性定理}（参见 EGA IV$_1$, 1.8.4），$\pi(\underline G^0\cap V)$ 在 $S$ 中\Term{locally-constructible}{局部可构造}。

所以由同处 1.10.1，要证 $\pi(\underline G^0\cap V)$ 是 $s$ 的\Term{open-neighborhood}{开邻域}，只需证明：对 $s$ 的每个\Term{generization}{一般化} $t$，都存在 $\eta$ 的一个\Term{generization}{一般化} $\xi$ 属于 $\underline G^0$（从而属于 $\underline G^0\cap V$）。$\underline G_t^0$ 的\Term{generic-point}{泛点} $\xi$ 正是 $\eta$ 的\Term{generization}{一般化}。事实上，$\varepsilon(s)$ 属于 $\{\xi\}$ 在 $G$ 中的\Term{closure}{闭包} $X$，故由 Chevalley \Term{semicontinuity}{半连续性}\Term{theorem}{定理}（参见 EGA IV$_3$, 13.1.3），
\[
\dim_{\varepsilon(s)}X_s\geq\dim_\xi X_t.
\]
另一方面，由假设 (ii)，$\dim_\xi\underline G_t^0=\dim_{\varepsilon(s)}G_s$。因此，$X_s$ 中包含 $\varepsilon(s)$ 的某个\Term{irreducible-component}{不可约分支}等于 $\underline G_s^0$，故 $\eta\in\overline{\{\xi\}}$。这证明了 (ii)$\Rightarrow$(iii)，从而证明命题。

\footnote{\textbf{编者注（30）：}上面的证明由 O. Gabber 提供；我们也保留了原文对 (ii)$\Rightarrow$(i) 的证明。}还可如下证明 (ii)$\Rightarrow$(i)。$\pi$ \Term{locally-finite-presentation}{局部有限呈示}，故存在 $G$ 中包含 $\varepsilon(s)$ 的\Term{open-subset}{开集} $U$，及 $S$ 中包含 $s$ 的\Term{open-subset}{开集} $V$，使 $\pi(U)\subset V$，且由 $\pi$ 诱导的\Term{morphism}{态射} $\pi':U\to V$ \Term{finite-presentation}{有限呈示}。置 $T=\varepsilon^{-1}(U)$、$W=\pi'^{-1}(T)=U\cap\pi^{-1}(T)$。则由 $\pi'$ 诱导的\Term{morphism}{态射} $\pi'':W\to T$ \Term{finite-presentation}{有限呈示}，且由 $\varepsilon$ 诱导的\Term{morphism}{态射} $\varepsilon'':T\to W$ 是其\Term{section}{截面}。对每个 $t\in T$，$\underline G_t^0$
\SourcePage{358}
\Term{irreducible}{不可约}，故 $W\cap\underline G_t^0$ 在其中\Term{dense}{稠密}，从而\Term{irreducible}{不可约}，因而\Term{connected}{连通}；它正是 $\pi''^{-1}(t)$ 中包含 $\varepsilon''(t)$ 的\Term{connected-component}{连通分支}。

由于 $W\cap\underline G_t^0$ 是 $\underline G_t^0$ 的\Term{dense-open-set}{稠密开集}，由 1.5 和 EGA IV$_2$, 5.2.1，
\[
\dim(W\cap\underline G_t^0)=\dim\underline G_t^0=\dim G_t.
\]
所以函数 $t\mapsto\dim(W\cap\underline G_t^0)$ 在 $s$ 于 $T$ 中的某个\Term{neighborhood}{邻域}上为常数。最后证明，对每个 $t\in T$，$W\cap\underline G_t^0$ \Term{geometrically-irreducible}{几何不可约}。若 $K$ 是 $\kappa(t)$ 的\Term{field-extension}{扩域}，则
\[
(W\cap\underline G_t^0)\otimes_{\kappa(t)}K
=(W\otimes_{\kappa(t)}K)\cap(\underline G_t^0\otimes_{\kappa(t)}K)
\]
是 $\underline G_t^0\otimes_{\kappa(t)}K$ 的\Term{nonempty}{非空}\Term{open-subset}{开集}；后者\Term{irreducible}{不可约}，故前者也\Term{irreducible}{不可约}。

于是可应用 EGA IV$_3$, 15.6.6 (ii)，得 $\pi''$（从而 $\pi$）在 $W\cap\underline G_s^0$ 的各点处\Term{universally-open}{普遍开}。但由 EGA IV$_3$, 14.3.3.1 (ii)，$G_s$ 中 $\pi$ 为\Term{universally-open}{普遍开}的点集在 $G_s$ 中闭；它既包含 $W\cap\underline G_s^0$，也就包含其\Term{closure}{闭包} $\underline G_s^0$。
\end{proof}
\begin{SGAresult}{推论}{4.3}
若 $G$ 是\Term{flat}{平坦}且\Term{locally-finite-presentation}{局部有限呈示}的 $S$-群，则函数 $s\mapsto\dim G_s$ 在 $S$ 上\Term{locally-constant}{局部常值}。
\end{SGAresult}
\begin{proof}
这直接来自 4.2，因为\Term{flat}{平坦}且\Term{locally-finite-presentation}{局部有限呈示}的\Term{morphism}{态射}总是\Term{universally-open}{普遍开}（EGA IV$_2$, 2.4.6）。
\end{proof}
\begin{SGAresult}{推论}{4.4}
设 $G$ 是 $S$-群，且在\Term{identity-section}{单位截面}的各点处在 $S$ 上\Term{locally-finite-presentation}{局部有限呈示}。考虑以下条件：
\begin{enumerate}
\item[(i)] $G$ 在\Term{identity-section}{单位截面}的各点处在 $S$ 上\Term{smooth}{光滑}（参见 3.10）。
\item[(ii)] 对每个 $s\in S$，$G_s$ 在 $\kappa(s)$ 上\Term{smooth}{光滑}，且 $s\mapsto\dim G_s$ 在 $S$ 上\Term{locally-constant}{局部常值}。
\item[(iii)] 对每个 $s\in S$，$G_s$ 在 $\kappa(s)$ 上\Term{smooth}{光滑}；并存在\Term{identity-section}{单位截面}的\Term{neighborhood}{邻域} $V$，使 $\pi_V:V\to S$ \Term{universally-open}{普遍开}。
\item[(iv)] 对每个 $s\in S$，$G_s^0$ 在 $\kappa(s)$ 上\Term{smooth}{光滑}，且 $G^0$ 由 $G$ 的一个在 $S$ 上\Term{universally-open}{普遍开}的\Term{open-subgroup-scheme}{开子群概形}\Term{representable}{表示}。
\end{enumerate}
则有 (i)$\Rightarrow$(ii)$\Leftrightarrow$(iii)$\Leftrightarrow$(iv)。若还假设 $S$ \Term{reduced}{约化}，则 (i)--(iv) \Term{equivalence}{等价}，并\Term{pushout}{推出} $G^0$ 在 $S$ 上\Term{smooth}{光滑}。\footnote{\textbf{编者注（31）：}这里当然不能去掉 $S$ \Term{reduced}{约化}的假设：若 $k$ 是域，$S=\operatorname{Spec}k[\delta]$，其中 $\delta^2=0$，则\Term{trivial-group}{平凡} $k$-群 $G=\operatorname{Spec}k$ 是一个满足 (ii)--(iv) 的 $S$-群，但在 $A$ 上不\Term{flat}{平坦}，因而不\Term{smooth}{光滑}。}
\end{SGAresult}
\begin{proof}
先证 (i) \Term{pushout}{推出} (ii)。由 1.5，对每个 $x\in G$，$\dim_x\pi^{-1}(\pi(x))=\dim\pi^{-1}(\pi(x))$。所以由 EGA IV$_4$, 17.10.2，函数
\[
x\longmapsto\dim_x\pi^{-1}(\pi(x))=\dim\pi^{-1}(\pi(x))
\]
在\Term{identity-section}{单位截面}附近\Term{continuous}{连续}；故 $s\mapsto\dim G_s$ 在 $S$ 上\Term{continuous}{连续}，从而\Term{locally-constant}{局部常值}。另一方面，由 1.3.1，每个 $G_s$ 都在 $\kappa(s)$ 上\Term{smooth}{光滑}。

再证 (ii) \Term{pushout}{推出} (iv)。只需证明 $\underline G^0$ 在 $G$ 中开；因为由 3.4，$G^0$ 就由 $\underline G^0$ 上诱导的\Term{group-subscheme}{子群概形}\Term{representable}{表示}，
\SourcePage{359}
而陈述中 $G^0$ 的性质由 2.4、4.2 得到。给定 $s\in S$，按 3.10 的证明构造 $W,T,\pi'',\varepsilon'',W^0$。由 EGA IV$_3$, 15.6.7，$W^0$ 在 $W$ 中开。另一方面，在假设 4.4 (ii) 下，4.2 说明 $\pi$ 在 $W^0$ 的每一点\Term{universally-open}{普遍开}，所以 $\mu$ 在 $W^0\times_SW^0$ 的每一点也\Term{universally-open}{普遍开}（第 VI$_A$ 讲 0.1）。这说明 $W^0\cdot W^0$ 开，然后仿照\Term{theorem}{定理} 3.10 的证明即可完成。

显然 (iv)$\Rightarrow$(iii)，而将 4.2 用于 $V$ 即得 (iii)$\Rightarrow$(ii)。

最后，假设 (ii)--(iv) 成立，且 $S$ \Term{reduced}{约化}，证明 $G^0$ 在 $S$ 上\Term{smooth}{光滑}。\footnote{\textbf{编者注（32）：}下面将原文展开说明。}可假设 $G=G^0$。由下面的 5.5，$G$ 在 $S$ 上\Term{finite-presentation}{有限呈示}。故 $\pi_G$ \Term{finite-presentation}{有限呈示}，\Term{fibre}{纤维}\Term{geometrically-integral}{几何整}，\Term{dimension}{维数}在 $S$ 上\Term{locally-constant}{局部常值}。由 EGA IV$_3$, 15.6.7，由 $\pi_G$ 诱导的\Term{morphism}{态射} $G\times_SS_{\mathrm{r\text{é}d}}\to S_{\mathrm{r\text{é}d}}$ \Term{flat}{平坦}，所以 $S$ \Term{reduced}{约化}时 $\pi_G$ \Term{flat}{平坦}。由\Term{theorem}{定理} 3.10，此时 $G=G^0$ 在 $S$ 上\Term{smooth}{光滑}。
\end{proof}

\SGAsection{5}{群与齐性空间的分离性}
\begin{SGAresult}{命题}{5.1}
$S$-群 $G$ \Term{separated}{分离}，当且仅当它的\Term{identity-section}{单位截面}是\Term{closed-immersion}{闭浸入}。
\end{SGAresult}
\begin{proof}
必要性见 EGA I, 5.4.6；充分性来自下列\Term{cartesian-square}{笛卡尔方块}（参见第 VI$_A$ 讲 0.3）：
\[
\begin{tikzcd}
G\times_SG\arrow[r,"{\mu\circ(\operatorname{id}_G\times c)}"]&G\\
G\arrow[u,"\Delta_{G/S}"]\arrow[r,"\pi"]&S\arrow[u,"\varepsilon"'].
\end{tikzcd}
\]
\end{proof}
\begin{SGAresult}{命题}{5.2}
若 $S$ \Term{discrete}{离散}，则每个 $S$-群都\Term{separated}{分离}。
\end{SGAresult}
\begin{proof}
此时 $S=\coprod_{s\in S}\operatorname{Spec}\mathcal O_{S,s}$。由 EGA I, 5.5.5，只需证明对每个 $s\in S$，$G\times_S\operatorname{Spec}\mathcal O_{S,s}$ \Term{separated}{分离}；$\mathcal O_{S,s}$ 是零维\Term{local-ring}{局部环}，故这来自第 VI$_A$ 讲 0.3。\footnote{\textbf{编者注（33）：}将“Artin”改成了“零维”（并在第 VI$_A$ 讲 0.3 中提及这一推广）。}
\end{proof}
\begin{SGAresult}{\Term{theorem}{定理}}{5.3}
设 $S$ 是\Term{scheme}{概形}，$G$ 是\Term{locally-finite-presentation}{局部有限呈示} $S$-\Term{group-scheme}{群概形}，且在\Term{identity-section}{单位截面}附近在 $S$ 上\Term{universally-open}{普遍开}。设 $X$ 是带 $G$-\Term{action}{作用}的 $S$-\Term{scheme}{概形}，使\Term{morphism}{态射}
\[
\Phi:G\times_SX\longrightarrow X\times_SX,\qquad(g,x)\longmapsto(gx,x)
\]
为\Term{surjective}{满射}。还假设对每个 $s\in S$：
\begin{enumerate}
\item[(i)] 存在 $X$ 的\Term{open-subscheme}{开子概形} $U$，它在 $S$ 上\Term{separated}{分离}，且 $U_s$ 在 $X_s$ 中\Term{dense}{稠密}；
\SourcePage{360}
\item[(ii)] \Term{fibre}{纤维} $X_s$ 在 $\kappa(s)$ 上\Term{locally-finite-type}{局部有限型}。\footnote{\textbf{编者注（34）：}我们修正了原文，假设 $G$ 在\Term{identity-section}{单位截面}附近在 $S$ 上\Term{universally-open}{普遍开}，并加入了假设 (ii)；下面 O. Gabber 的例子说明，若无附加假设，原文 5.3、5.4 的陈述是不成立的。此外，\Term{theorem}{定理} 5.3 是下面\Term{theorem}{定理} 5.3A 的改写版本；后者见于 1965 年版 SGAD，由 M. Raynaud 给出，参见第 X 讲 8.5、8.8 的注（$*$）。\par\textbf{\Term{theorem}{定理} 5.3A（Raynaud）。}设 $G$ 是\Term{locally-finite-presentation}{局部有限呈示} $S$-群，在 $S$ 上\Term{universally-open}{普遍开}，且\Term{fibre}{纤维}\Term{connected}{连通}。则 $G$ 在 $S$ 上\Term{separated}{分离}。更一般地，若 $S$-\Term{scheme}{概形} $X$ 在 $S$ 上\Term{locally-finite-presentation}{局部有限呈示}，且带有一个 $G$-\Term{action}{作用}，使\Term{morphism}{态射} $\Phi:G\times_SX\to X\times_SX$、$(g,x)\mapsto(gx,x)$ 为\Term{surjective}{满射}，则 $X$ 在 $S$ 上\Term{separated}{分离}。}
\end{enumerate}
则 $X$ 在 $S$ 上\Term{separated}{分离}。
\end{SGAresult}
\begin{SGAresult}{推论}{5.4}
设 $S,G,X$ 如 5.3 的预备假设。还假设 $X$ 的\Term{fibre}{纤维}\Term{locally-finite-type}{局部有限型}且\Term{connected}{连通}。则：
\begin{enumerate}
\item[(i)] $X$ 在 $S$ 上\Term{separated}{分离}。
\item[(ii)] 若 $X$ 有一个在 $S$ 上\Term{quasi-compact}{拟紧}的\Term{open-subset}{开集} $V$，与每个\Term{nonempty}{非空}\Term{fibre}{纤维} $X_s$ 相交，\footnote{\textbf{编者注（35）：}若无此假设，有 O. Gabber 指出的下列反例：设 $G=S$ 是\Term{closed-point}{闭点}为 $s$ 的\Term{local-scheme}{局部概形}，且 $S^*=S-\{s\}$ 不\Term{quasi-compact}{拟紧}；取 $X$ 为 $\{s\}$ 与 $S^*$ 的\Term{disjoint-union}{不交并}。}则 $X$ 在 $S$ 上\Term{quasi-compact}{拟紧}。
\end{enumerate}
\end{SGAresult}
\begin{proof}
(i) 取满足 $X_s\ne\varnothing$ 的 $s\in S$。由 $\Phi$ 经\Term{base-change}{基变换}所得的\Term{morphism}{态射} $\Phi_s:G_s\times_{\kappa(s)}X_s\to X_s\times_{\kappa(s)}X_s$ 为\Term{surjective}{满射}，且 $X_s$ \Term{connected}{连通}，故由第 VI$_A$ 讲 2.5.4，$X_s$ \Term{irreducible}{不可约}。若 $U$ 是 $X$ 中满足 $U_s\ne\varnothing$ 的\Term{affine}{仿射}\Term{open-subset}{开集}，则 $U_s$ 在 $X_s$ 中\Term{dense}{稠密}，故可应用\Term{theorem}{定理}。

证明 (ii) 时可假设 $S$ \Term{affine}{仿射}。则 $V$ \Term{quasi-compact}{拟紧}；由 3.5，$G$ 有一个包含 $\underline G^0$ 的\Term{quasi-compact}{拟紧}\Term{open-subset}{开集} $U$。取满足 $X_s\ne\varnothing$ 的 $s\in S$。$X_s$ \Term{irreducible}{不可约}（第 VI$_A$ 讲 2.6.6），且 $U_s$ 包含 $\underline G_s^0$，所以\Term{morphism}{态射} $U_s\times_{\kappa(s)}V_s\to X_s$、$(g,x)\mapsto gx$ 是\Term{surjective}{满射}（第 VI$_A$ 讲 2.6.4）。因此\Term{morphism}{态射} $U\times_SV\to X$ 是\Term{surjective}{满射}，$U,V$ 又\Term{quasi-compact}{拟紧}，故 $X$ \Term{quasi-compact}{拟紧}。$S$ 已假设\Term{affine}{仿射}，故\Term{separated}{分离}，从而 $X$ 在 $S$ 上\Term{quasi-compact}{拟紧}（参见 EGA I, 6.6.4 (v)）。
\end{proof}
\textbf{注记 5.4.1。}\footnote{\textbf{编者注（36）：}加入了这条注记，参见编者注（34）。}证明中将看到，仅作如下假设，\Term{theorem}{定理} 5.3 的结论仍然成立：
\begin{quote}
(i$'$) 存在 $X$ 的\Term{open-subscheme}{开子概形} $U$，在 $S$ 上\Term{separated}{分离}，且 $U_s$ 在 $X_s$ 的每个\Term{irreducible-component}{不可约分支}中\Term{dense}{稠密}。
\end{quote}
（若 $X_s$ \Term{local-scheme}{局部}只有\Term{finite}{有限}个\Term{irreducible-component}{不可约分支}，则 (i) \Term{pushout}{推出}这一条件；假设 (ii) 下正是如此。）另一方面，下面将看到，若每个\Term{fibre}{纤维} $X_s$ \Term{quasi-separated}{拟分离}且\Term{connected}{连通}，5.4 也成立。
\SourcePage{361}
\begin{SGAresult}{推论}{5.5}
设 $S$ 是\Term{scheme}{概形}，$G$ 是\Term{locally-finite-presentation}{局部有限呈示} $S$-\Term{group-scheme}{群概形}，其\Term{fibre}{纤维}\Term{connected}{连通}，且在 $S$ 上\Term{universally-open}{普遍开}。则 $G$ 在 $S$ 上\Term{separated}{分离}且\Term{finite-presentation}{有限呈示}。\footnote{\textbf{编者注（37）：}这里指出如下结果（[Ray70a], VI 2.5）：若 $S$ \Term{normal}{正规}，$G$ \Term{smooth}{光滑}且\Term{fibre}{纤维}\Term{connected}{连通}，$X$ 是在 $S$ 上\Term{locally-finite-type}{局部有限型}的 (fppf) $G$-\Term{homogeneous-space}{齐性空间}（即 $X\to S$ 和 $G\times_SX\to X\times_SX$ 是 (fppf) \Term{topology}{拓扑}的\Term{covering}{覆盖}），则 $X$ 在 $S$ 上\Term{locally-quasi-projective}{局部拟射影}。特别地，$G$ 在 $S$ 上\Term{quasi-projective}{拟射影}。亦见第 VI$_A$ 讲编者注（35）。}
\end{SGAresult}
\begin{proof}
由 3.6、5.4，$G$ 在 $S$ 上\Term{quasi-compact}{拟紧}且\Term{separated}{分离}；它还在 $S$ 上\Term{locally-finite-presentation}{局部有限呈示}，故在 $S$ 上\Term{finite-presentation}{有限呈示}。
\end{proof}
\textbf{5.6。\Term{theorem}{定理} 5.3 的证明。}证明 5.3 之前，先证明几个\Term{lemma}{引理}。
\begin{SGAresult}{\Term{lemma}{引理}}{5.6.0}
\footnote{\textbf{编者注（38）：}加入了 O. Gabber 提供的这个\Term{lemma}{引理}，它改进了 EGA IV$_3$, 14.4.1.2，并在不改变假设的前提下修正了同处 14.4.1.3 的证明（与 EGA IV$_4$ 中的勘误 Err$_{\mathrm{IV}}$, 38 比较）。}
\begin{enumerate}
\item[(i)] 设 $A\subset B$ 是两个\Term{integral-domain}{整环}，$B$ 在 $A$ 上整；设 $\mathfrak q\in\operatorname{Spec}B$，并假设 $A$ 在 $\mathfrak p=\mathfrak q\cap A$ 处\Term{unibranch}{单支}。则\Term{morphism}{态射} $\pi:\operatorname{Spec}B\to\operatorname{Spec}A$ 在 $\mathfrak q$ 处是\Term{open-morphism}{开}的。
\item[(ii)] 设 $X,Y$ 是两个\Term{irreducible}{不可约}\Term{prescheme}{预概形}，$f:X\to Y$ 是\Term{dominant}{支配态射}，$x\in X$；假设 $f$ 在 $x$ 处\Term{quasi-finite}{拟有限}，且 $y=f(x)$ 是 $Y$ 的\Term{unibranch}{单支点}。则 $f$ 在 $x$ 处是\Term{open-morphism}{开}的。特别地，若 $Y=\operatorname{Spec}\mathcal O_{Y,y}$ 是\Term{closed-point}{闭点}为 $y$ 的\Term{local-prescheme}{局部预概形}，则对 $x$ 的每个\Term{neighborhood}{邻域} $U$，$f(U)=Y$。
\end{enumerate}
\end{SGAresult}
\begin{proof}
(i) 令 $K,L$ 分别为 $A,B$ 的\Term{fraction-field}{分式域}，$A'$ 为 $A$ 的\Term{normalization}{正规化}，$B'$ 为 $L$ 中由 $A'$ 与 $B$ 生成的\Term{subring}{子环}。则 $B'$ 在 $A'$ 上整。置
\[
Y=\operatorname{Spec}A,\quad X=\operatorname{Spec}B,\quad
Y'=\operatorname{Spec}A',\quad X'=\operatorname{Spec}B',
\]
从而有\Term{commutative-diagram}{交换图}
\[
\begin{tikzcd}
X'\arrow[r,"\pi'"]\arrow[d]&Y'\arrow[d,"\varphi"]\\
X\arrow[r,"\pi"]&Y,
\end{tikzcd}
\]
其中所有\Term{morphism}{态射}都\Term{integral-morphism}{整}且为\Term{surjective}{满射}。

$A$ 在 $\mathfrak p$ 处\Term{unibranch}{单支}，所以 $Y'$ 在 $\mathfrak p$ 上方只有一个点 $\mathfrak p'$。故若 $U$ 是 $\mathfrak p'$ 在 $Y'$ 中的\Term{open-neighborhood}{开邻域}，则 $\varphi(Y'-U)$ 是不含 $\mathfrak p$ 的闭集，其补\Term{open-subset}{开集}包含在 $\varphi(U)$ 中。这说明 $\varphi:Y'\to Y$ 在 $\mathfrak p'$ 处是\Term{open-morphism}{开}的，因此只需证明 $\pi'$ 是\Term{open-morphism}{开态射}。于是归结为 $A=A'$ \Term{normal}{正规}的情形。

取 $K$ 的一个包含 $L$ 的\Term{quasi-galois}{拟 Galois 扩张} $N$。令 $\widetilde X=\operatorname{Spec}\widetilde B$，其中 $\widetilde B$ 是 $B$ 在 $N$ 中的\Term{integral-closure}{整闭包}，并令 $G=\operatorname{Aut}(N/K)$。$\widetilde X\to X$ 是\Term{surjective}{满射}，所以只需证明 $\widetilde\pi:\widetilde X\to Y$ 是\Term{open-morphism}{开态射}。取 $\widetilde X$ 的\Term{open-subset}{开集} $U$，置 $U'=\bigcup_{g\in G}gU$。$G$ 在 $\widetilde X\to Y$ 的\Term{fibre}{纤维}上传递地\Term{action}{作用}（参见 [BAC], 第 V 章 \S2.3，命题 6），所以 $\widetilde\pi(U)=\widetilde\pi(U')$，而后者就是闭集 $\widetilde\pi(\widetilde X-U')$ 的补\Term{open-subset}{开集}。这证明了 (i)。
\SourcePage{362}
(ii) 可假设 $Y,X$ \Term{reduced}{约化}。由 \Term{zariski-main-theorem}{Zariski 主定理}（参见 EGA IV$_3$, 8.12.9），存在 $x$ 的\Term{affine}{仿射}\Term{open-neighborhood}{开邻域} $U$ 和 $y$ 的\Term{affine}{仿射}\Term{open-neighborhood}{开邻域} $V=\operatorname{Spec}A$，使 $f(U)\subset V$，并有\Term{resolution}{分解}
\[
\begin{tikzcd}
U\arrow[r,hook,"j"]\arrow[dr,"f"']&V'\arrow[d,"u"]\\
&V,
\end{tikzcd}
\]
其中 $j$ 是\Term{open-immersion}{开浸入}，$u$ 是\Term{finite-morphism}{有限态射}。将 $V'$ 换成 $j(U)$ 的\Term{closure}{闭包}，可假设 $V'$ \Term{irreducible}{不可约}，所以 $V'=\operatorname{Spec}B$，其中 $B$ 是在 $A$ 上\Term{finite-algebra}{有限}的\Term{integral-domain}{整} $A$-\Term{algebra}{代数}。此外，$f$ \Term{dominant}{支配}，故 $u$ 也\Term{dominant}{支配}，所以环\Term{morphism}{态射} $A\to B$ 为\Term{injectivity}{单射}。$A$ 按假设在 $y$ 处\Term{unibranch}{单支}，故由 (i)，$u$ 在 $j(x)$ 处是\Term{open-morphism}{开}的，从而 $f=u\circ j$ 在 $x$ 处是\Term{open-morphism}{开}的。这证明 (ii) 的第一个断言。第二个断言随之而来，因为若 $Y$ 是\Term{closed-point}{闭点}为 $y$ 的\Term{local-prescheme}{局部预概形}，则包含 $y$ 的\Term{open-subset}{开集}必等于 $Y$。（若 $Y=\operatorname{Spec}\mathcal O_{Y,y}$，还可用 \Term{zariski-main-theorem}{Zariski 主定理}的\Term{local-scheme}{局部}\Term{form}{形式}代替 EGA IV$_3$, 8.12.9；例如见 [Pes66]，或 [Ray70b] 第 IV 章\Term{theorem}{定理} 1。）
\end{proof}
\begin{SGAresult}{\Term{lemma}{引理}}{5.6.1.0}
\footnote{\textbf{编者注（39）：}明确写出了\Term{lemma}{引理} 5.6.1 的证明所用的这个\Term{lemma}{引理}。}设 $f:X\to S$ 是\Term{locally-finite-presentation}{局部有限呈示}\Term{morphism}{态射}，$s\in S$、$x\in X_s$。置 $n=\dim_xX_s$，令 $q:\mathbf A_S^n\to S$ 为\Term{structure-morphism}{结构态射}。

假设给定\Term{quasi-finite}{拟有限}\Term{morphism}{态射} $u:X\to\mathbf A_S^n$，使 $f=q\circ u$；并假设 $f$ 在 $X_s$ 的某个包含 $x$、\Term{dimension}{维数}为 $n$ 的\Term{irreducible-component}{不可约分支} $\mathfrak z$ 的\Term{generic-point}{泛点} $z$ 处\Term{universally-open}{普遍开}。则 $u$ 在 $x\in X$ 处\Term{universally-open}{普遍开}。
\end{SGAresult}
\begin{proof}
首先注意：($\dagger$) 这些假设在任何\Term{covering}{覆盖} $s$ 的\Term{base-change}{基变换} $\pi:S'\to S$ 下保持成立（即 $\pi^{-1}(s)\ne\varnothing$）。事实上，取这样的 $S'\to S$，及经\Term{base-change}{基变换}所得的图
\[
\begin{tikzcd}
X'\arrow[r,"u'"]\arrow[dr,"f'"']&\mathbf A_{S'}^n\arrow[d,"q'"]\\
&S',
\end{tikzcd}
\]
取 $s'\in S'$ 位于 $s$ 上方，$x'\in X'_{s'}$ 位于 $x$ 上方。$X'_{s'}=X_s\otimes_{\kappa(s)}\kappa(s')$，故由\Term{generization}{一般化}的\Term{lifting}{提升}及\Term{dimension}{维数}在\Term{field-extension}{扩域}下的\Term{invariant}{不变}性（EGA IV$_2$, 2.3.4 (i)、4.1.4），$x'$ 包含于 $X'_{s'}$ 的一个\Term{irreducible-component}{不可约分支} $\mathfrak z'$，其\Term{generic-point}{泛点} $z'$ 位于 $z$ 上方，且
\[
n\leq\dim\mathfrak z'\leq\dim_{x'}X'_{s'}\leq n,
\]
故 $\dim\mathfrak z'=n=\dim_{x'}X'_{s'}$。$f$ 在 $x$ 处\Term{universally-open}{普遍开}，故 $f'$ 在 $x'$ 处也\Term{universally-open}{普遍开}。

置 $Y=\mathbf A_S^n$。由 EGA IV$_3$, 14.3.3.1 (i)，要证 $u$ 在 $x$ 处\Term{universally-open}{普遍开}，只需对每个整数 $r\geq0$ 及 $X'=X[T_1,\ldots,T_r]$ 中
\SourcePage{363}
位于 $x$ 上方的每个点 $x'$，证明 $u':X'\to Y'=\mathbf A_S^n[T_1,\ldots,T_r]$ 在 $x'$ 处是\Term{open-morphism}{开}的。以 $S'=S[T_1,\ldots,T_r]$，$q'$ 为\Term{projection}{投影} $Y'\simeq\mathbf A_{S'}^n\to S'$，则这正是\Term{base-change}{基变换} $S'\to S$ 所得的情形。由上，归结为证明 $u$ 在 $x$ 处是\Term{open-morphism}{开}的。

置 $y=u(x)$。$f,q$ 都\Term{locally-finite-presentation}{局部有限呈示}，故 $u$ 亦然（参见 EGA IV$_1$, 1.4.3）；由 EGA IV$_1$, 1.10.3，只需证明
\[
u(\operatorname{Spec}\mathcal O_{X,x})=\operatorname{Spec}\mathcal O_{Y,y}.
\]
为此，可假设 $S$ \Term{affine}{仿射}且\Term{integral-scheme}{整}。令 $\pi:S'\to S$ 为其\Term{normalization}{正规化}，以 $u':X'\to Y'$ 表示 $u$ 经\Term{base-change}{基变换}所得的\Term{morphism}{态射}。$\pi_Y:Y'\to Y$ \Term{integral-morphism}{整}且为\Term{surjective}{满射}，故
\[
\operatorname{Spec}\mathcal O_{Y,y}
=\bigcup_{y'}\pi_Y(\operatorname{Spec}\mathcal O_{Y',y'}),
\]
其中并取遍 $Y'$ 中位于 $y$ 上方的所有点。因此只需证明，对每个这样的 $y'$，及 $X'_{y'}$ 中位于 $x$ 上方的每个点 $x'$，有
\[
u'(\operatorname{Spec}\mathcal O_{X',x'})=\operatorname{Spec}\mathcal O_{Y',y'}.
\]
这些假设在\Term{base-change}{基变换} $\pi:S'\to S$ 下保持成立，故归结为 $S'$ 的情形，即可假设 $S$ 是\Term{integral-scheme}{整}\Term{normal}{正规}\Term{scheme}{概形}。

由 EGA IV$_3$, 14.3.13，对 $f$ 的假设说明：$X$ 有一个\Term{irreducible-component}{不可约分支} $Z$ 包含 $z$（因而包含 $x$），\Term{dominant}{支配} $S$，且
\[
\dim_z Z_s=n=\dim Z_\eta,
\]
其中 $\eta$ 是 $S$ 的\Term{generic-point}{泛点}。令 $\xi$ 是 $Z$ 的\Term{generic-point}{泛点}。$u$ 因而 $u_\eta:Z_\eta\to\mathbf A_\eta^n$ 都\Term{quasi-finite}{拟有限}，所以 $u_\eta(\xi)=u(\xi)$ 在 $\mathbf A_\eta^n$ 中的\Term{closure}{闭包}\Term{dimension}{维数}为 $n$。故 $u(\xi)$ 是 $\mathbf A_\eta^n$ 的\Term{generic-point}{泛点}，也是 $Y=\mathbf A_S^n$ 的\Term{generic-point}{泛点}。因此，以 $g$ 表示 $u$ 在 $Z$ 上的\Term{restriction}{限制}，则 $g:Z\to Y$ \Term{quasi-finite}{拟有限}且\Term{dominant}{支配}。$Y$ \Term{normal}{正规}，由\Term{lemma}{引理} 5.6.0，$g$ 是\Term{open-morphism}{开态射}。于是 $u$ 在 $Z$ 的每一点都是\Term{open-morphism}{开}的，特别在 $x$ 处如此。故 $u(\operatorname{Spec}\mathcal O_{X,x})=\operatorname{Spec}\mathcal O_{Y,y}$，\Term{lemma}{引理} 5.6.1.0 得证。
\end{proof}
下面的\Term{lemma}{引理}有利地替代了 EGA IV$_3$, 14.5.10，\footnote{\textbf{编者注（40）：}原文误作 19.5.10，已修正。}因为它不依赖 Noether 性假设。
\begin{SGAresult}{\Term{lemma}{引理}}{5.6.1}
设 $S$ 是\Term{scheme}{概形}，$f:X\to S$ 是\Term{locally-finite-presentation}{局部有限呈示} $S$-\Term{scheme}{概形}，$s\in S$，$x$ 是 $X_s$ 的\Term{closed-point}{闭点}。假设 $f$ 在 $X_s$ 的某个包含 $x$、且满足 $\dim_xX_s=\dim\mathfrak z$ 的\Term{irreducible-component}{不可约分支} $\mathfrak z$ 的\Term{generic-point}{泛点} $z$ 处\Term{universally-open}{普遍开}。则存在\Term{commutative-diagram}{交换图}
\[
\begin{tikzcd}
X\arrow[r,"f"]&S\\
S''\arrow[u,"h"]\arrow[ur,"\varphi"]\arrow[r,"\pi"']&S'\arrow[u,"w"'],
\end{tikzcd}
\]
\SourcePage{364}
其中 $S'$ 是\Term{affine}{仿射}\Term{scheme}{概形}，$w$ 是\Term{etale}{平展}\Term{morphism}{态射}，$\pi$ 是\Term{finite-presentation}{有限呈示}的\Term{finite-morphism}{有限}\Term{surjective}{满态射}，且 $\varphi^{-1}(s)$ 只含一个点 $s''$，满足 $h(s'')=x$、$\kappa(s'')=\kappa(x)$。\footnote{\textbf{编者注（41）：}因此 $h$ 诱导 $X\times_SS''\to S''$ 的一个\Term{section}{截面} $\sigma$，使 $\sigma(s'')$ 位于 $x$ 上方；与 EGA IV$_3$, 14.5.10 比较。}
\end{SGAresult}
\begin{proof}
首先可假设 $S=\operatorname{Spec}A$、$X=\operatorname{Spec}B$，其中 $B$ 是\Term{finite-presentation}{有限呈示} $A$-\Term{algebra}{代数}。令 $\mathfrak p,\mathfrak q$ 是 $A,B$ 中分别对应于 $s,x$ 的\Term{prime-ideal}{素理想}，于是 $\mathfrak p=A\cap\mathfrak q$。置 $n=\dim_xX_s$，取 $t_1,\ldots,t_n\in B$，使其在 $\mathcal O_{X_s,x}=B_{\mathfrak q}/\mathfrak pB_{\mathfrak q}$ 中的像构成\Term{system-of-parameters}{参数系}。则 $\mathcal O_{X_s,x}/(t_1,\ldots,t_n)$ 在 $\kappa(x)$ 上\Term{dimension}{维数}\Term{finite}{有限}，因 $x$ 是\Term{algebraic-scheme}{代数} $\kappa(s)$-\Term{scheme}{概形} $X_s$ 的\Term{closed-point}{闭点}，故它在 $\kappa(s)$ 上的\Term{dimension}{维数}也\Term{finite}{有限}。因此，由 $T_i\mapsto t_i$ 定义的 $S$-\Term{morphism}{态射}
\[
u:X\longrightarrow\mathbf A_S^n=\operatorname{Spec}A[T_1,\ldots,T_n]
\]
\Term{finite-presentation}{有限呈示}，$x$ 在其\Term{fibre}{纤维} $u^{-1}(u(x))$ 中孤立，且 $u(x)=\tau_0(s)$；这里 $\tau_0:S\to\mathbf A_S^n$ 是 $\mathbf A_S^n\to S$ 的“\Term{zero-section}{零截面}”，对应于将各 $T_i$ 送到 $0$ 的 $A$-\Term{algebra}{代数}\Term{morphism}{态射} $A[T_1,\ldots,T_n]\to A$。$X$ 中在 $\mathbf A_S^n$ 上方的自身\Term{fibre}{纤维}中孤立的点集是\Term{open-subset}{开集}（EGA IV$_3$, 13.1.4），所以缩小 $X$ 后，可假设 $u$ \Term{quasi-finite}{拟有限}且 $u^{-1}(u(x))=\{x\}$。由\Term{lemma}{引理} 5.6.1.0，$u$ 在 $x$ 处\Term{universally-open}{普遍开}。

令 $B_0=B/(t_1,\ldots,t_n)$、$X_0=X\times_{\mathbf A_S^n}\tau_0(S)=\operatorname{Spec}B_0$，$u_0:X_0\to S$ 为 $u$ 经\Term{base-change}{基变换} $\tau_0:S\hookrightarrow\mathbf A_S^n$ 所得的\Term{morphism}{态射}。则 $u_0$ \Term{quasi-finite}{拟有限}且\Term{finite-presentation}{有限呈示}，在 $x$ 处\Term{universally-open}{普遍开}，而 $x$ 是 $X_0$ 中位于 $s$ 上方的唯一一点。

令 $A'$ 是\Term{local-ring}{局部环} $A_{\mathfrak p}=\mathcal O_{S,s}$ 的\Term{henselization}{Hensel 化}，置 $S'=\operatorname{Spec}A'$、$B'_0=B_0\otimes_AA'$、$X'_0=X_0\times_SS'=\operatorname{Spec}B'_0$。则 $S'$ 的\Term{closed-point}{闭点} $s'$ 是其中位于 $s$ 上方的唯一一点，$\kappa(s')=\kappa(s)$，而 $X'_0$ 在 $s'$ 上方只有一个点 $x'$；有 $\kappa(x')=\kappa(x)$，且 $x'$ 也是 $X'_0$ 中位于 $s$ 和 $x$ 上方的唯一一点。$A'$ 是\Term{henselian}{Hensel 环}，故由 EGA IV$_4$, 18.5.11，$X'_0$ 是两个既开又闭部分的\Term{disjoint-union}{不交并}：
\[
X'_0=V\sqcup W,\qquad V=\operatorname{Spec}\mathcal O_{X'_0,x'}.
\tag{$*$}
\]
而\Term{local-ring}{局部环} $\mathcal O_{X'_0,x'}$ 在 $A'$ 上\Term{finite-algebra}{有限}且\Term{finite-presentation}{有限呈示}。所以 $u'_0$ 在 $V$ 上的\Term{restriction}{限制} $\pi$ \Term{finite-morphism}{有限}且\Term{finite-presentation}{有限呈示}。又因 $u'_0:X'_0\to S'$ 在 $x'$ 处是\Term{open-morphism}{开}的，$u'_0(V)=\operatorname{Spec}\mathcal O_{S',s'}=S'$，故 $\pi$ 为\Term{surjective}{满射}。

这证明了 $S=S'$ 时的所需结果。一般地，$A'$ 是各个在 $A$ 上\Term{etale}{平展}的\Term{subalgebra}{子代数} $A_i$ 的\Term{filtered}{滤过}\Term{inductive-limit}{归纳极限}，且 $S_i=\operatorname{Spec}A_i$ 在 $s$ 上方只有一个点 $s_i$（并有 $\kappa(s_i)=\kappa(s)$）。于是 $B'_0=\varinjlim B_i$，其中 $B_i=B_0\otimes_AA_i$。置 $X_i=\operatorname{Spec}B_i=X_0\times_SS_i$，$C=\mathcal O_{X'_0,x'}$。由上面的 ($*$)，对某个\Term{idempotent}{幂等元} $f\in B'_0$，有 $C\simeq B'_0/fB'_0$；存在指标 $i$ 及 $f_i\in B_i$，使 $f$ 是 $f_i$ 在 $B'_0$ 中的像。置 $C_i=B_i/f_iB_i$、$V_i=\operatorname{Spec}C_i$。
\SourcePage{365}
则 $C=C_i\otimes_{A_i}A'$，故 $V=V_i\times_{S_i}S'$；而 $C_i$ 是\Term{finite-presentation}{有限呈示} $A_i$-\Term{algebra}{代数}（因为 $B_i$ 如此）。因此，\Term{morphism}{态射}图
\[
\begin{tikzcd}
X'_0\arrow[dr,"u'_0"]&\\
V\arrow[u,"h'"]\arrow[r,"\pi"']&S'
\end{tikzcd}
\]
来自\Term{finite-presentation}{有限呈示}\Term{morphism}{态射}图
\[
\begin{tikzcd}
X_i\arrow[dr,"u_i"]&\\
V_i\arrow[u,"h_i"]\arrow[r,"\pi_i"']&S_i
\end{tikzcd}
\]
经 $S'\to S_i$ 的\Term{base-change}{基变换}。对每个 $j\geq i$，置 $V_j=V_i\times_{S_i}S_j$，令 $\pi_j:V_j\to S_j$ 为由 $\pi_i$ 经\Term{base-change}{基变换}所得的（\Term{finite-presentation}{有限呈示}）\Term{morphism}{态射}。$\pi:V\to S'$ \Term{finite-morphism}{有限}且为\Term{surjective}{满射}，故由 EGA IV$_2$, 8.10.5，存在指标 $j$，使 $\pi_j:V_j\to S_j$ \Term{finite-morphism}{有限}且为\Term{surjective}{满射}。此时 $w:S_j\to S$ \Term{etale}{平展}且\Term{affine}{仿射}，$S_j$ 在 $s$ 上方只有一个点 $s_j$；$V_j$ 在 $s_j$ 上方只有一个点 $x_j$（因为 $x'$ 是 $V=V_j\times_{S_j}S'$ 中位于 $s'$ 上方的唯一一点）：
\[
\begin{tikzcd}[column sep=large]
X_j\arrow[r]\arrow[dr,"u_j"]&X_0\arrow[r]\arrow[dr,"u_0"]&X\arrow[d,"f"]\arrow[dr,"u"]&\\
V_j\arrow[u,"h_j"]\arrow[r,"\pi_j"']&S_j\arrow[r,"w"']&S\arrow[r,"\tau_0"']&\mathbf A_S^n.
\end{tikzcd}
\]
所以 $x_j$ 是 $V_j$ 中位于 $s$ 上方的唯一一点，它在\Term{morphism}{态射} $V_j\to X_j\to X$ 下的像为 $x$，且 $\kappa(x_j)=\kappa(x)$。
\end{proof}
\begin{SGAresult}{\Term{lemma}{引理}}{5.6.2.0}
\footnote{\textbf{编者注（42）：}加入了 O. Gabber 提供的这个\Term{lemma}{引理}。它使 5.6.2 的证明得以简化，并使\Term{theorem}{定理} 5.3 及 5.4 得到更一般的\Term{form}{形式}，见下面 5.7、5.8。}设 $k$ 是域，$G$ 是\Term{nonempty}{非空}\Term{finite-type}{有限型} $k$-\Term{scheme}{概形}。
\begin{enumerate}
\item[(i)] 设 $K$ 是 $k$ 的\Term{field-extension}{扩域}，$W$ 是 $G_K$ 的\Term{dense-open-set}{稠密开集}。以 $\pi$ 表示\Term{projection}{投影} $G_K\to G$，则
\[
U=\{g\in G\mid W\text{ 包含 }\pi^{-1}(g)\text{ 的每个极大点}\}
\]
是 $G$ 的\Term{dense-open-set}{稠密开集}。（注意，若 $g$ 是 $G$ 中属于 $U$ 的\Term{closed-point}{闭点}，则 $\pi^{-1}(g)\subset W$。）
\item[(ii)] 再假设 $G$ \Term{geometrically-irreducible}{几何不可约}。设 $\mu:G\times X\to Y$ 是 $k$-\Term{scheme}{概形}\Term{morphism}{态射}，$x\in X$，$\Omega$ 是 $Y$ 的\Term{open-subset}{开集}，满足 $\mu_x^{-1}(\Omega_{\kappa(x)})\ne\varnothing$；这里 $\mu_x$ 表示\Term{morphism}{态射} $G_{\kappa(x)}\to Y_{\kappa(x)}$、$g\mapsto\mu(g,x)$。则
\[
U=\{g\in G\mid\mu\text{ 将 }g\times x\text{ 的每个极大点映入 }\Omega\}
\]
\SourcePage{366}
是 $G$ 的\Term{dense-open-set}{稠密开集}；对每个属于 $U$ 的\Term{closed-point}{闭点} $g$，有 $\mu(g\times x)\subset\Omega$，因而对 $G_{\kappa(x)}$ 中位于 $g$ 上方的每个点 $g'$（相应地 $X_{\kappa(g)}$ 中位于 $x$ 上方的每个点 $x'$），都有 $\mu(g',x)\in\Omega_{\kappa(x)}$（相应地 $\mu(g,x')\in\Omega_{\kappa(g)}$）。
\item[(iii)] 假设 $G=H^0$，其中 $H$ 是\Term{locally-finite-type}{局部有限型} $k$-\Term{group-scheme}{群概形}，\Term{action}{作用}于\Term{nonempty}{非空} $k$-\Term{scheme}{概形} $X$，使\Term{morphism}{态射} $H\times_SX\to X\times_SX$、$(h,x)\mapsto(hx,x)$ 为\Term{surjective}{满射}。设 $U$ 是 $X$ 的\Term{open-subset}{开集}。则：
\begin{enumerate}
\item[(a)] $G\cdot U$ 是 $X$ 的\Term{open-subset}{开集}，等于 $X$ 中\Term{generic-point}{泛点}属于 $U$ 的那些\Term{irreducible-component}{不可约分支}之并。
\item[(b)] 若 $U$ 在 $X$ 的每个\Term{irreducible-component}{不可约分支}中\Term{dense}{稠密}，则对 $X$ 的每个\Term{finite}{有限}\Term{subset}{子集} $A$，存在\Term{closed-point}{闭点} $g\in G$，使 $g\cdot A'\subset U_{\kappa(g)}$，其中 $A'$ 是 $A$ 在 $X_{\kappa(g)}$ 中的\Term{inverse-image}{逆像}。特别地，\Term{morphism}{态射} $G\times U\to X$ 为\Term{surjective}{满射}。
\end{enumerate}
\end{enumerate}
\end{SGAresult}
\begin{proof}
(i) 令 $\overline k$ 是 $k$ 的\Term{separable-closure}{可分闭包}，$L$ 是 $k$ 的一个包含 $\overline k$ 和 $K$ 的副本的\Term{field-extension}{扩域}。以 $\pi_L$（相应地 $\varphi$）表示\Term{projection}{投影} $G_L\to G$（相应地 $G_L\to G_K$）。对每个 $g\in G$，$\varphi$ 将 $\pi_L^{-1}(g)$ 的\Term{maximal-point}{极大点}满地映到 $\pi^{-1}(g)$ 的\Term{maximal-point}{极大点}，所以
\[
U=\{g\in G\mid W_L\text{ 包含 }\pi_L^{-1}(g)\text{ 的每个极大点}\}.
\]
因此将 $K$ 换成 $L$，可假设 $K$ 包含 $\overline k$。

\Term{projection}{投影} $p:G_K\to G_{\overline k}$ 是\Term{surjective}{满}\Term{open-morphism}{开态射}，所以 $V=p(W)$ 是 $G_{\overline k}$ 的\Term{dense-open-set}{稠密开集}；又对每个 $g\in G_{\overline k}$，$p^{-1}(g)=\operatorname{Spec}(\kappa(g)\otimes_{\overline k}K)$ \Term{irreducible}{不可约}（参见 EGA IV$_2$, 4.3.3），故对每个 $g\in V$，$p^{-1}(g)$ 的\Term{generic-point}{泛点}属于 $W$。另一方面，令 $\mathscr G=\operatorname{Gal}(\overline k/k)$。\Term{projection}{投影} $q:G_{\overline k}\to G$ 是\Term{surjective}{满射}，$\mathscr G$ 在 $q$ 的\Term{fibre}{纤维}上传递地\Term{action}{作用}，所以 $U=q(V')$，其中 $V'$ 是 $V$ 的所有 $\mathscr G$-\Term{conjugacy}{共轭}的交。

令 $Z$ 是闭集 $G_{\overline k}-V$，赋予\Term{reduced}{约化}\Term{closed-subscheme}{闭子概形}结构。$G_{\overline k}$ 从而 $Z$ 都在 $\overline k$ 上\Term{finite-presentation}{有限呈示}，所以存在 $k$ 的某个\Term{finite-galois-extension}{有限 Galois 扩张} $k_1$，使 $Z$ 是 $G\otimes_kk_1$ 的某个\Term{reduced}{约化}\Term{closed-subscheme}{闭子概形} $Z_1$ 的\Term{base-change}{基变换}。故 $Z$ 的 $\mathscr G$-\Term{conjugacy}{共轭}只有\Term{finite}{有限}个，它们的并仍是 $G_{\overline k}$ 的一个\Term{nowhere-dense}{无处稠密}闭集 $F$，而 $V'=G_{\overline k}-F$ 是 $G_{\overline k}$ 的\Term{dense-open-set}{稠密开集}。$q:G_{\overline k}\to G$ 是\Term{surjective}{满}\Term{open-morphism}{开态射}，所以 $U=q(V')$ 是 $G$ 的\Term{dense-open-set}{稠密开集}。此外，对 $G$ 的每个\Term{closed-point}{闭点} $g$，\Term{fibre}{纤维} $\pi^{-1}(g)$ 由 $G_K$ 的\Term{finite}{有限}个\Term{closed-point}{闭点}组成；故若 $g\in U$，则 $\pi^{-1}(g)\subset W$。这证明了 (i)，而 (ii) 随之而来。

另一方面，(iii)(a) 已在第 VI$_A$ 讲 2.6.4 中证明。最后，若 $U$ 在 $X$ 的每个\Term{irreducible-component}{不可约分支}中\Term{dense}{稠密}，则 $X=G\cdot U$，所以对每个 $x\in X$，$\mu_x^{-1}(U_{\kappa(x)})$ 是 $G_{\kappa(x)}$ 的\Term{nonempty}{非空}\Term{open-subset}{开集}；于是 (iii)(b) 由 (ii) 得到。
\end{proof}
\begin{SGAresult}{推论}{5.6.2}
设 $S,G,X$ 如 5.3 的预备假设，$U$ 是 $X$ 的\Term{open-subset}{开集}，$s\in S$，$A$ 是 $X_s$ 的\Term{finite}{有限}\Term{subset}{子集}。假设 $U_s$ 在 $X_s$ 中\Term{dense}{稠密}，且 $X_s$ 在 $\kappa(s)$ 上\Term{locally-finite-type}{局部有限型}。\footnote{\textbf{编者注（43）：}加入了对 $X_s$ 的假设，并利用增补 5.6.2.0 简化（且修正）了证明。此外，证明说明，仅假设 $U_s$ 在 $X_s$ 的每个\Term{irreducible-component}{不可约分支}中\Term{dense}{稠密}，结论就成立。}
\SourcePage{367}
则存在\Term{morphism}{态射} $f:S''\to S$，它是\Term{finite-morphism}{有限}\Term{surjective}{满态射} $S''\to S'$ 与\Term{etale}{平展}\Term{morphism}{态射} $S'\to S$ 的复合；还存在\Term{morphism}{态射} $h:S''\to G$，使 $A$ 在 $X\times_SS''$ 中（即在 $X_s\times_{\operatorname{Spec}\kappa(s)}S''_s$ 中）的\Term{inverse-image}{逆像} $A''$ 包含于 $\ell_h^{-1}(U_{S''})$，其中 $\ell_h$ 是由 $h\in G(S'')$ 定义的 $X_{S''}=X\times_SS''$ 上的\Term{translation}{平移}。
\end{SGAresult}
\begin{proof}
$X_s$ 在 $\kappa(s)$ 上\Term{locally-finite-type}{局部有限型}，故其\Term{connected-component}{连通分支}都是开且\Term{irreducible}{不可约}的（参见第 VI$_A$ 讲 2.5.4），所以 $U_s$ 在 $X_s$ 的每个\Term{irreducible-component}{不可约分支}中\Term{dense}{稠密}。

由\Term{lemma}{引理} 5.6.2.0，存在\Term{closed-point}{闭点} $g\in\underline G_s^0$，使对 $X_{\kappa(g)}$ 中位于 $A$ 的点上方的每个点 $a'$，$g\cdot a'\in U_s\otimes_{\kappa(s)}\kappa(g)$。

由\Term{lemma}{引理} 5.6.1，存在\Term{morphism}{态射} $f:S''\to S$，为\Term{finite-morphism}{有限}\Term{surjective}{满态射} $S''\to S'$ 与\Term{etale}{平展}\Term{morphism}{态射} $S'\to S$ 的复合；并存在\Term{morphism}{态射} $h:S''\to G$，使 $f^{-1}(s)$ 只含一点 $s_0$，且 $h(s_0)=g$、$\kappa(s_0)=\kappa(g)$。于是，以 $A''$ 表示 $A$ 在
\[
X_s\times_{\operatorname{Spec}\kappa(s)}S''_s
=X_s\otimes_{\kappa(s)}\kappa(s_0)
=X_s\otimes_{\kappa(s)}\kappa(g)
\]
中的\Term{inverse-image}{逆像}，则 $X_{S''}$ 上的\Term{translation}{平移} $\ell_h$ 将 $A''$ 送入 $U_s\otimes_{\kappa(s)}\kappa(s_0)$。
\end{proof}
\begin{SGAresult}{\Term{lemma}{引理}}{5.6.3}
设 $X$ 是 $S$-\Term{scheme}{概形}。下列条件\Term{equivalence}{等价}：
\begin{enumerate}
\item[(i)] $X$ 在 $S$ 上\Term{separated}{分离}。
\item[(ii)] 对每个 $S$-\Term{scheme}{概形} $T$，每个\Term{section}{截面} $\sigma:T\to X_T$ 都是\Term{closed-immersion}{闭浸入}。
\item[(iii)] 对每个\Term{reduced}{约化} $S$-\Term{scheme}{概形} $T$，在 $T$ 的一个\Term{dense-open-set}{稠密开集} $U$ 上相同的两个 $S$-\Term{morphism}{态射} $f_1,f_2:T\to X$ 必相等。
\item[(iv)] 对每个 $s\in S$ 及每对点 $x_1,x_2\in X_s$，存在\Term{morphism}{态射} $\varphi:S''\to S'\to S$，及 $X_{S''}$ 中在 $S''$ 上\Term{separated}{分离}的\Term{open-subscheme}{开子概形} $V$，满足：
\begin{enumerate}
\item[a)] $S'\to S$ 是\Term{open-morphism}{开态射}，$S''\to S'$ 是\Term{closed-morphism}{闭}\Term{surjective}{满态射}，且 $\varphi^{-1}(s)\ne\varnothing$；
\item[b)] $\{x_1,x_2\}$ 在 $X_{S''}$ 中的\Term{inverse-image}{逆像}包含于 $V$。
\end{enumerate}
\item[(iv$'$)] 对每个 $s\in S$，$X_s$ 的每对点 $x_1,x_2$ 都包含在 $X$ 的某个在 $S$ 上\Term{separated}{分离}的\Term{open-subset}{开集} $V$ 中。\footnote{\textbf{编者注（44）：}简化了条件 (iv) 的表述，并加入条件 (iv$'$)。又补上了 (i)$\Rightarrow$(iii) 的证明，(iv)$\Rightarrow$(iii) 的证明将用到它。另外注意，若 $T=\operatorname{Spec}k[\varepsilon,x]/(\varepsilon^2,\varepsilon x)$（$k$ 为域），$X=T_{\mathrm{r\text{é}d}}=\operatorname{Spec}k[x]$，则由 $x\mapsto x+\lambda\varepsilon$（$\lambda\in k$）定义的\Term{morphism}{态射} $\varphi_\lambda:T\to X$ 在\Term{dense-open-set}{稠密开集} $\operatorname{Spec}B_x$ 上相同，但彼此不相等。}
\end{enumerate}
\end{SGAresult}
\begin{proof}
(ii)$\Rightarrow$(i) 显然（取 $T=X$，$\sigma$ 为\Term{diagonal-section}{对角截面}），(i)$\Rightarrow$(iv$'$)$\Rightarrow$(iv) 也显然。另一方面，由 EGA I, 5.4.6，(i)$\Rightarrow$(ii)。

(iii)$\Rightarrow$(ii)。设 $\sigma:T\to X_T$ 是 $p:X_T\to T$ 的\Term{section}{截面}。由 EGA I, 5.3.13，$\sigma$ 是\Term{immersion}{浸入}，即 $T$ 与 $X_T$ 的某个\Term{locally-closed-subscheme}{局部闭子概形} $E$ 之间的\Term{isomorphism}{同构}。要证 $E$ 闭，可假设 $T,E$ \Term{reduced}{约化}。令 $\overline E$ 为 $X_T$ 中以 $E$ 的\Term{closure}{闭包}为底\Term{space}{空间}的\Term{reduced}{约化}\Term{closed-subscheme}{闭子概形}，则 $E$ 是 $\overline E$ 的\Term{dense}{稠密}\Term{open-subscheme}{开子概形}。\Term{immersion}{浸入} $i:\overline E\hookrightarrow X_T$ 与 $\sigma\circ p\circ i$ 在 $E$ 上相同，由 (iii)，它们在 $\overline E$ 上也相同。所以 $\overline E$ 的每个点都属于 $\sigma(T)=E$，故 $\overline E=E$。
\SourcePage{368}
(i)$\Rightarrow$(iii)。假设 $X$ 在 $S$ 上\Term{separated}{分离}，$T$ 为\Term{reduced}{约化} $S$-\Term{scheme}{概形}，$f_1,f_2:T\to X$ 是在\Term{dense-open-set}{稠密开集} $U$ 上相同的两个 $S$-\Term{morphism}{态射}。令 $g:T\to X\times_SX$ 为\Term{set-component}{分量} $f_1,f_2$ 所给出的\Term{morphism}{态射}。$D=\Delta_{X/S}(X)$ 闭，故其在 $g$ 下的\Term{inverse-image}{逆像}为 $T$ 中包含 $U$ 的闭集，从而等于 $T$。$T$ \Term{reduced}{约化}，故 $g$ 通过 $D$ \Term{resolution}{分解}（参见 EGA I, 5.2.2）；因此 $f_1=p_1\circ g=p_2\circ g=f_2$。

(iv)$\Rightarrow$(iii)。设 $T$ 是\Term{reduced}{约化} $S$-\Term{scheme}{概形}，$f_1,f_2:T\to X$ 是在\Term{dense-open-set}{稠密开集} $U$ 上相同的两个 $S$-\Term{morphism}{态射}。$T$ \Term{reduced}{约化}，故要证 $f_1=f_2$，只需证明二者在\Term{set}{集合}意义下相同。事实上，若此事成立，取 $t\in T$，取 $X$ 中包含 $f_1(t)=f_2(t)$ 的\Term{affine}{仿射}\Term{open-subset}{开集} $V$，并令 $W$ 是 $f_1,f_2$ 的共同底\Term{continuous-map}{连续映射}下 $V$ 的\Term{inverse-image}{逆像}，这是 $t$ 的\Term{open-neighborhood}{开邻域}。则 $f_i|_W:W\to V$ 在 $W$ 的\Term{dense-open-set}{稠密开集} $U\cap W$ 上相同。$V$ \Term{separated}{分离}、$W$ \Term{reduced}{约化}，故 $f_1|_W=f_2|_W$，从而 $f_1=f_2$。

于是，取 $t\in T$，以 $s$ 表示它在 $S$ 中的像，证明 $X_s$ 中两点 $x_1=f_1(t)$、$x_2=f_2(t)$ 相等。取 (iv) 所给的 $\varphi:S''\to S'\to S$ 及 $X\times_SS''$ 的\Term{open-subset}{开集} $V$；置 $T'=T\times_SS'$、$T''=T\times_SS''$，以 $g:T''\to T$ 和 $f_i'':T''\to X_{S''}$（$i=1,2$）表示经\Term{base-change}{基变换}所得的\Term{morphism}{态射}。

$U$ 在 $T$ 中\Term{dense}{稠密}，$T'\to T$ 为\Term{open-morphism}{开态射}，故 $U$ 在 $T'$ 中的\Term{inverse-image}{逆像} $U'$ 在 $T'$ 中\Term{dense}{稠密}。令 $U''$ 是 $U'$ 在 $T''$ 中的\Term{inverse-image}{逆像}，$F$ 是 $T''$ 中以 $\overline{U''}$ 为底\Term{space}{空间}的\Term{reduced}{约化}\Term{subscheme}{子概形}。$T''\to T'$ 是\Term{surjective}{满}\Term{closed-morphism}{闭态射}，故 $F$ 的像是包含 $U'$ 的闭集，因而等于 $T'$。所以 $F\cap g^{-1}(t)$ 包含一点 $u$。

对 $i=1,2$，令 $h_i=f_i''|_F$。则 $h_i(u)$ 是 $X_{S''}$ 中位于 $f_i(t)=x_i$ 上方的一点，因 $V$ 包含 $\{x_1,x_2\}$ 在 $X_{S''}$ 中的\Term{inverse-image}{逆像}，故 $h_i(u)\in V$。

于是 $W=h_1^{-1}(V)\cap h_2^{-1}(V)$ 是 $F$ 中包含 $u$ 的\Term{open-subset}{开集}，两个 $S''$-\Term{morphism}{态射} $h_i|_W:W\to V$ 在 $W$ 的\Term{dense-open-set}{稠密开集} $U''\cap W$ 上相同。$V$ 在 $S''$ 上\Term{separated}{分离}，故 $h_1|_W=h_2|_W$，于是 $h_1(u)=h_2(u)$，从而 $f_1(t)=f_2(t)$。这证明了 (iv)$\Rightarrow$(iii)。

\Term{theorem}{定理} 5.3 现在由 5.6.2 及 5.6.3 的 (iv)$\Rightarrow$(i) 得到。
\end{proof}
\begin{SGAresult}{反例}{5.6.4}
并非每个 $S$-群 $G$ 都\Term{separated}{分离}。设 $S$ 是有一个非孤立\Term{closed-point}{闭点} $s$ 的\Term{scheme}{概形}，$G$ 是将两份 $S$ 沿\Term{open-subset}{开集} $S-\{s\}$ \Term{gluing}{粘合}所得的\Term{scheme}{概形}。容易看出，$G$ 在 $S$ 上不\Term{separated}{分离}，且自然带有 $S$-\Term{group-structure}{群结构}，除 $G_s$ 外各\Term{fibre}{纤维}均为\Term{unit-group}{单位群}，而 $G_s$ 与二元群 $\mathbf Z/2\mathbf Z$ \Term{isomorphism}{同构}。\footnote{\textbf{编者注（45）：}这个例子也见于第 VI$_A$ 讲 0.3 编者注（5）。}
\end{SGAresult}
\begin{SGAresult}{\Term{theorem}{定理}}{5.7}
\footnote{\textbf{编者注（46）：}一方面，删去了重复 5.5 的推论 5.6.5。另一方面，原文在注记 5.7 中陈述了下面的推论 5.7.1，并将证明指向 4.7；Lect. Notes 151 中并无此编号（但 1965 年版 SGAD 中有，其 4.5、4.6 已变为 5.6.1、5.6.2）。我们加入了 O. Gabber 提供的\Term{theorem}{定理} 5.7，它使 SGAD 的上述陈述更为精确。}设 $S$ 是\Term{scheme}{概形}，$G$ 是\Term{locally-finite-presentation}{局部有限呈示} $S$-群，函数 $s\mapsto\dim G_s$ 在 $S$ 上\Term{locally-constant}{局部常值}。设 $X$ 是
\SourcePage{369}
带 $G$-\Term{action}{作用}的 $S$-\Term{scheme}{概形}，$U$ 是 $X$ 中在 $S$ 上\Term{separated}{分离}的\Term{open-subset}{开集}。则 $\underline G^0\cdot U$ 是 $X$ 中在 $S$ 上\Term{separated}{分离}的\Term{open-subset}{开集}。
\end{SGAresult}
\begin{proof}
令 $\mu:G\times_SX\to X$ 表示 $G$ 在 $X$ 上的\Term{action}{作用}；它是 $G\times_SX$ 的\Term{automorphism}{自同构} $(g,x)\mapsto(g,gx)$ 与到 $X$ 的\Term{projection}{投影}的复合。$G\times_SU$ 是 $G\times_SX$ 的\Term{open-subset}{开集}，由 4.2 (iii)，$V=\underline G^0\cdot U$ 在 $X$ 中开。

然后仿照 5.6.2 的证明，由 5.6.3 的 (iv)$\Rightarrow$(i) 得 $V$ 在 $S$ 上\Term{separated}{分离}。
\end{proof}
\begin{SGAresult}{推论}{5.7.1}
设 $S,G$ 如 5.7，$\sigma,\tau$ 是 $G^0$ 的两个 $S$-\Term{section}{截面}（即 $\sigma,\tau\in G^0(S)$）。则 $\sigma,\tau$ 的\Term{coincidence-subscheme}{重合子概形} $S(\sigma,\tau)\subset S$（即 $G\times_SG$ 的\Term{diagonal}{对角线}在\Term{morphism}{态射} $(\sigma,\tau)$ 下的\Term{inverse-image}{逆像}）是闭的。
\end{SGAresult}
\begin{proof}
对每个 $s\in S$，取 $G$ 中包含 $\varepsilon(s)$ 的\Term{affine}{仿射}\Term{open-subset}{开集} $U$。则 $V=\varepsilon^{-1}(U)$ 是 $s$ 在 $S$ 中的\Term{open-neighborhood}{开邻域}；$\underline G^0U$ \Term{separated}{分离}，所以 $S(\sigma,\tau)\cap V$ 在 $V$ 中闭。
\end{proof}
\textbf{注记 5.7.2。}Gabber 告知我们，可证明：若 $S$ 是\Term{closed-point}{闭点}为 $s$ 的\Term{henselian-local-scheme}{Hensel 局部概形}，则当 $U$ 遍历 $\varepsilon(s)$ 的\Term{affine}{仿射}\Term{open-neighborhood}{开邻域}时，各\Term{open-subset}{开集} $\underline G^0U$ 的交是 $G$ 的一个在 $S$ 上\Term{separated}{分离}的\Term{open-subgroup-scheme}{开子群概形}。

\textbf{5.8。补充。}\footnote{\textbf{编者注（47）：}加入了这一段补充和反例，全部由 O. Gabber 提供。}先回顾第 VI$_A$ 讲命题 2.6.6。
\begin{SGAresult}{命题}{5.8.1}
设 $k$ 是域，$G$ 是\Term{locally-finite-type}{局部有限型} $k$-群，\Term{action}{作用}于 $k$-\Term{scheme}{概形} $X$，使\Term{morphism}{态射} $G\times X\to X\times X$、$(g,x)\mapsto(gx,x)$ 为\Term{surjective}{满射}。假设 $X$ \Term{quasi-separated}{拟分离}。则 $X$ 的\Term{connected-component}{连通分支}都\Term{irreducible}{不可约}。
\end{SGAresult}
\begin{SGAresult}{推论}{5.8.2}
设 $S,G,X$ 如 5.3 的预备假设。还假设每个\Term{fibre}{纤维} $X_s$ 都\Term{quasi-separated}{拟分离}且\Term{connected}{连通}。则 $X$ 在 $S$ 上\Term{separated}{分离}且\Term{quasi-compact}{拟紧}。
\end{SGAresult}
\begin{proof}
由上一个命题，每个\Term{fibre}{纤维} $X_s$ 都\Term{irreducible}{不可约}，余下证明与 5.4 相同。
\end{proof}
\begin{SGAresult}{例}{5.8.3}
固定\Term{algebraically-closed-field}{代数闭域} $k$。首先回顾，每个“\Term{locally-boolean}{局部布尔}”\Term{topological-space}{拓扑空间} $X$（即有一个由紧\Term{open-subset}{开集}组成的基），赋予取值于 $k$ 的\Term{locally-constant}{局部常值}函数\Term{sheaf}{层}后，都是一个 $k$-\Term{scheme}{概形}（参见 [DG70], I \S1, 2.12）。这时称 $X$ 为\Term{locally-boolean}{局部布尔} $k$-\Term{scheme}{概形}。

另一方面，每个具有\Term{hausdorff}{分离}\Term{open-subset}{开集}基的\Term{topological-space}{拓扑空间} $E$（以及每个\Term{scheme}{概形} $X$）都有\Term{hausdorff}{分离}的\Term{dense-open-set}{稠密开集}。事实上，\Term{hausdorff}{分离}\Term{open-subset}{开集}的任何递增并仍是\Term{hausdorff}{分离}\Term{open-subset}{开集}（对于\Term{scheme}{概形} $X$，这由 5.6.3 得到），所以存在一个极大的这样的\Term{open-subset}{开集} $U$。$U$ 必\Term{dense}{稠密}：若有\Term{nonempty}{非空}\Term{open-subset}{开集} $V$ 使 $U\cap V=\varnothing$，则 $V$ 包含
\SourcePage{370}
一个\Term{nonempty}{非空}\Term{hausdorff}{分离}\Term{open-subset}{开集} $W$，而 $U\cup W$ 仍\Term{hausdorff}{分离}，与 $U$ 的极大性矛盾。

现在，令 $C$ 为 Cantor \Term{space}{空间}，可将它看作赋有\Term{product-topology}{积拓扑}的群 $(\mathbf Z/2\mathbf Z)^{\mathbf N}$ 的底\Term{space}{空间}。对每个 $p\in C$，取另一份 $C$ 记作 $C(p)$，将各 $C(p)$ 沿 $C-\{p\}$ 与 $C$ \Term{gluing}{粘合}所得的\Term{space}{空间}记作 $X$。则 $X$ 是非\Term{separated}{分离}的\Term{locally-boolean}{局部布尔} $k$-\Term{scheme}{概形}，而 $C$ 是 $X$ 的\Term{dense-open-set}{稠密开集}。

令 $G$ 为 $k$-\Term{scheme}{概形} $X$ 的\Term{automorphism}{自同构}群（即 $X$ 的\Term{homeomorphism}{同胚}群）。则 $G$ 在 $X$ 上传递地\Term{action}{作用}。事实上，$X$ 由 $C$ 以及对每个 $p\in C$ 增加的第二个点 $p'$ 组成；$p'$ 的\Term{characteristic}{特征}是 $X-\{p\}$ 中使 $(C-\{p\})\cup\{x\}$ \Term{hausdorff}{分离}的唯一点 $x$。由此，$C$ 的每个\Term{automorphism}{自同构} $\varphi$ 都延拓为 $X$ 的\Term{automorphism}{自同构} $\varphi_X$，使每个 $p$ 都满足 $\varphi_X(p')=\varphi_X(p)'$。另一方面，交换 $p,p'$ 而固定其他各点的映射 $\theta_p:X\to X$ 是 $X$ 的\Term{automorphism}{自同构}。最后，$C$ 的\Term{automorphism}{自同构}群在 $C$ 上传递地\Term{action}{作用}：例如利用 $C$ 的\Term{group-structure}{群结构}，考虑各\Term{translation}{平移}即可。

所以\Term{discrete}{离散} $k$-群 $G$（因而\Term{locally-finite-type}{局部有限型}）\Term{action}{作用}于 $k$-\Term{scheme}{概形} $X$，使\Term{morphism}{态射} $G\times X\to X\times X$、$(g,x)\mapsto(gx,x)$ 为\Term{surjective}{满射}，但 $X$ 不\Term{separated}{分离}（尽管 $C$ 是\Term{separated}{分离}的\Term{dense-open-set}{稠密开集}）。
\end{SGAresult}
\begin{SGAresult}{例}{5.8.4}
沿用上例记号。将 $C$ 描述为 $\mathbf Z/2\mathbf Z$ 中序列 $(u_n)_{n\in\mathbf N}$ 的\Term{set}{集合}，可见 $C$ 去掉任一点 $p$ 后，与可数个 $C$ 的\Term{disjoint-union}{不交并}\Term{homeomorphism}{同胚}。事实上，例如利用 $C$ 的\Term{group-structure}{群结构}，通过\Term{translation}{平移}可归结为 $p$ 为\Term{identity-element}{单位元} $0$、即零序列的情形。此时 $C-\{0\}$ 是各子\Term{space}{空间}
\[
C_i=\{(u_n)_{n\in\mathbf N}\mid u_i=1,\ u_j=0\text{ 对所有 }j<i\}\qquad(i\in\mathbf N)
\]
的\Term{disjoint-union}{不交并}，每个 $C_i$ 都与 $C$ \Term{homeomorphism}{同胚}。对 $C$ 的每个\Term{nonempty}{非空}\Term{finite}{有限}\Term{subset}{子集} $F$，按 $|F|$ 归纳，得 $C-F$ 与可数个 $C$ 的\Term{disjoint-union}{不交并}\Term{homeomorphism}{同胚}，从而与 $C$ 去掉一点所得的\Term{space}{空间}\Term{homeomorphism}{同胚}。

对每个二元\Term{subset}{子集} $F$，取另一份 $C$ 记作 $C(F)$，以 $q_F$ 表示 $C(F)$ 中的点 $0$，选定\Term{homeomorphism}{同胚} $\varphi_F:C(F)-\{q_F\}\xrightarrow{\sim}C-F$。将各 $C(F)$ 通过 $\varphi_F$ 与 $C$ \Term{gluing}{粘合}所得的\Term{space}{空间}记作 $X'$。则 $X'$ 是非\Term{separated}{分离}的\Term{locally-boolean}{局部布尔} $k$-\Term{scheme}{概形}。由构造，$X'$ 上每个\Term{locally-constant}{局部常值}函数 $f:X'\to k$ 都是常函数。事实上，若 $x,y\in C$、$F=\{x,y\}$，则 $q_F$ 的每个\Term{neighborhood}{邻域}都与 $x$ 或 $y$ 的每个\Term{neighborhood}{邻域}相交，所以若 $f:X\to k$ \Term{locally-constant}{局部常值}，则 $f(x)=f(q_F)=f(y)$；若 $F'=\{z,t\}$，其中 $z,t\in C$，则同样有 $f(q_{F'})=f(z)=f(q_{\{z,x\}})=f(x)$。因此 $X'$ \Term{connected}{连通}。

此外，每个 $x\in X'$ 都有一个带点\Term{neighborhood}{邻域}与 $(C,0)$ \Term{homeomorphism}{同胚}。更确切地说，对每个 $F$，固定带点\Term{topological-space}{拓扑空间}的\Term{homeomorphism}{同胚}
\[
h_F:(C,0)\xrightarrow{\sim}(C(F),q_F),
\]
并以 $T$ 表示 $C$ 的\Term{translation}{平移}群。若 $x=q_F$，则有所述\Term{homeomorphism}{同胚} $q_F$；若 $x\in C$，则\Term{translation}{平移} $t_x\in T$ 是\Term{homeomorphism}{同胚} $(C,0)\xrightarrow{\sim}(C,x)$（并且是 $T$ 中具有此性质的唯一元素）。
\SourcePage{371}
令 $L$ 为各 $h_F$ 生成的\Term{free-group}{自由群}，令 $G=T*L$ 为 $T,L$ 的“\Term{free-product}{自由积}”（即\Term{coproduct}{余积}）。对每个 $h\in\{h_F\}_{|F|=2}\cup T$，以 $\sigma(h)$ 表示 $G$ 中对应于 $h$ 的\Term{generator}{生成元}，以 $S(h)$（相应地 $B(h)$）表示 $h$ 的源（相应地靶）。还约定 $\sigma(h_F^{-1})=\sigma(h_F)^{-1}$ 及 $S(h_F^{-1})=B(h_F)$（相应地 $B(h_F^{-1})=S(h_F)$）；以 $E$ 表示由 $T$ 及各 $h_F,h_F^{-1}$ 组成的\Term{set}{集合}。

在积\Term{space}{空间} $P=G\times X'$ 上（$G$ 赋予\Term{discrete-topology}{离散拓扑}），考虑由下列关系生成的\Term{equivalence-relation}{等价关系}：对所有 $h\in E$，
\[
(g\sigma(h),x)\sim(g,h(x))\qquad\text{当 }x\in S(h).
\]
令 $Z$ 为\Term{quotient-space}{商空间}。$Z$ 是从\Term{disjoint-union}{不交并} $\coprod_{g\in G}\{g\}\times X'$ 通过\Term{open-subset}{开集}的\Term{gluing}{粘合}得到的；因而对 $P$ 的每个\Term{open-subset}{开集} $\Omega$，其\Term{saturation}{饱和集} $\widetilde\Omega$ 是\Term{open-subset}{开集}（参见 [BTop], I \S5.1，例 2）。具体说，$P$ 的每个\Term{open-subset}{开集}都是它与各“切片” $\{g\}\times X'$ 的交之并，所以只需考虑形如 $\{1\}\times W$ 的\Term{open-subset}{开集}，其中 $W$ 是 $X'$ 的\Term{open-subset}{开集}。此时\Term{saturation}{饱和集}是 $\{1\}\times W$ 以及
\begin{gather*}
\{\sigma(h_1)\}\times h_1^{-1}(W\cap B(h_1)),\\
\{\sigma(h_1)\sigma(h_2)\}\times h_2^{-1}\bigl(h_1^{-1}(W\cap B(h_1))\cap B(h_2)\bigr),
\end{gather*}
等等的并，取遍 $E$ 中元素的所有\Term{finite}{有限}序列 $h_1,\ldots,h_n$；所以它是\Term{open-subset}{开集}。因此\Term{projection}{投影} $\pi:P\to Z$ 是\Term{open-morphism}{开映射}。

还要注意，字 $\sigma(h_1)\cdots\sigma(h_n)$ 是 $G$ 的\Term{reduced-word}{既约字}，除非某个 $h_i$ 是 $T$ 的\Term{identity-element}{单位元}，或两个相邻的 $h_i$ 都属于 $T$，或 $(h_i,h_{i+1})$ 等于 $(h_F,h_F^{-1})$ 或 $(h_F^{-1},h_F)$。因此，若 $x\in X'$，且元素
\[
\beta=\bigl(\sigma(h_1)\cdots\sigma(h_n),h_n^{-1}\cdots h_1^{-1}(x)\bigr)
\]
属于 $\{1\}\times X'$，则可假设每个 $h_i$ 都是\Term{translation}{平移} $t_i$。此时等式 $\sigma(t_1\cdots t_n)=1$ \Term{pushout}{推出} $t_1\cdots t_n=1$，故 $\beta=(1,x)$。\Term{equivalence-relation}{等价关系}与 $G$ 的\Term{action}{作用}相容（$G$ 在 $P=G\times X'$ 的第一因子上作\Term{left-translation}{左平移}），所以对每个 $g\in G$，$\pi$ 在 $\{g\}\times X'$ 上的\Term{restriction}{限制}为\Term{injectivity}{单射}。

任取 $z\in Z$，取其代表元 $(g,x)\in P$。由上，$U=\pi(\{g\}\times X')$ 是 $z$ 的\Term{open-neighborhood}{开邻域}，而由 $\pi$ 诱导的\Term{continuous-map}{连续映射} $\{g\}\times X'\to U$ 为开\Term{bijection}{双射}，故是\Term{homeomorphism}{同胚}。这说明 $Z$ \Term{local-scheme}{局部}与 $X'$（从而也与 $C$）\Term{isomorphism}{同构}，所以仍是\Term{locally-boolean}{局部布尔} $k$-\Term{scheme}{概形}。

最后，$G$ 在 $Z$ 上传递地\Term{action}{作用}。事实上，每个 $z\in Z$ 都与某个形如 $\pi((1,x))$ 的元素 $G$-\Term{conjugacy}{共轭}，故只需证明每个 $(1,x)\in P$ 都\Term{equivalence}{等价}于一个形如 $(\sigma(h),0)$ 的元素。若 $x\in C$，可取 $h$ 为\Term{translation}{平移} $t_x$；若 $x=q_F$，可取 $h=h_F$。因此 $Z$ 是带有\Term{discrete}{离散}群 $G$ 的\Term{transitive-action}{传递作用}的\Term{locally-boolean}{局部布尔} $k$-\Term{scheme}{概形}，但 $Z$ 不\Term{separated}{分离}。
\end{SGAresult}
